\documentclass[../main.tex]{subfiles}
\begin{document}
%==========================================TIÊU ĐỀ TẬP========================================
\begin{table}[h]
  \begin{adjustwidth}{-1cm}{}
    \begin{tabular}{>{\centering\arraybackslash}p{6cm}|>{\raggedright\arraybackslash}p{14cm}}
      \multirow{5}{6cm}
      % Cột bên trái: ÉPISODE và số tập
      {\\[-40pt]
        \flaregothic\fontsize{20pt}{20pt}\selectfont \centering
        \color{eptype}HISTOIRE\color{black}\\
        \gangofthree\fontsize{72pt}{72pt}\selectfont
        \color{epnum}6
        \\[18pt]
        % \flaregothic\fontsize{14pt}{14pt}\selectfont
        % \color{epattr}BIENVENUE, \uline{\color{Banpire} Mel} !
        % \flaregothic\fontsize{18pt}{18pt}\selectfont
        % \color{epattr}ÉPISODE PERDU
        \color{black}
      }
       &
      % Dòng thứ nhất: tên tập tiếng Nhật
      {\fotskip\fontsize{18pt}{18pt}\selectfont \ruby{二}{ふた}\ruby{日}{か}めの\ruby{母}{はは}の\ruby{里}{さと}}
      \\[6pt]
      % Dòng thứ hai: tên tập tiếng Anh
       & {\cafeteria\fontsize{18pt}{18pt}\selectfont The Second Day at Mother's Hometown!}               \\[4pt]
      % Dòng thứ ba: tên tập tiếng Việt
       & {\vnmsans\fontsize{18pt}{18pt}\selectfont Mùng Hai Tết Mẹ!}                                     \\[8pt]
      % Dòng thứ tư: Nhân vật xuất hiện
       & {\brian{\fontsize{14pt}{14pt}\selectfont Track used: Lunar Zoo Year - Demonstration Mini-Game}}
    \end{tabular}
  \end{adjustwidth}
\end{table}
%===========================================NỘI DUNG==========================================
\color{black}

\begin{center}
  \uline{Nhà Hikawa\Bốn giờ sáng.}
\end{center}

*TEN TÈN TEN TÉN TÈN\dots*

Chuông báo thức điện thoại của cậu bắt đầu vang lên inh ỏi lúc tờ mờ sáng, phá tan bầu
không khí yên lặng trong căn phòng.

Kuon tỉnh dậy, đầu còn hơi nặng vì giấc ngủ chập chờn đêm trước, cầm lấy chiếc điện
thoại và nhanh chóng tắt chuông trước khi nó đánh thức cả phòng.

Cậu khẽ ngồi dậy, đưa mắt nhìn quanh phòng. Trên giường, bảy cô gái đang ngủ say sưa.
Miko và Rushia mỗi người chiếm một góc, trong khi Noel nằm nghiêng, tay khẽ vắt qua
chân của Roboco. Pekora thì cuộn tròn như một chú thỏ, còn Okayu và Korone nằm sát
nhau, như thể đang chia sẻ một giấc mơ chung.

Cậu nhẹ nhàng, rón rén bước qua từng đống chăn gối rối bời mà các cô gái đã để lại sau
một đêm dài. Nhìn họ ngủ, cậu không khỏi cảm thấy tình huống này thật lạ lùng, nhưng
cũng đầy ấm áp.

Sau khi vệ sinh cá nhân và mặc quần áo, Kuon xuống tầng ba, nơi mà bố mẹ cậu cũng
vừa mới ngủ dậy trước đó. Kaien đang chuẩn bị bữa sáng đơn giản, chủ yếu là vài món
dư thừa từ bữa tối hôm qua.

``{\color{red} Kuon dậy sớm thế?}'' bà cười nhẹ khi thấy cậu bước vào. Trong cơn ngái ngủ, cậu khẽ
cười: ``{\color{red} Sớm sủa gì đâu mẹ\dots\ Bình thường mọi năm giờ này mình đã phải dậy rồi\dots}''

``{\color{red} Mấy đứa vẫn đang ngủ ngon đấy chứ?}'' Ryujin cất giọng hỏi, rít một điếu thuốc lá đầu
ngày. Cậu gật đầu: ``{\color{red} Dạ, mọi người vẫn đang ngủ cả,}'' đôi mắt cậu vẫn còn phảng phất sự
mệt mỏi, nhưng giọng nói tỏ ra rất rõ ràng và có chút kinh nghiệm.

``{\color{red}Không biết tới khi nào họ dậy nhỉ\dots}'' cô em gái Ryuuhou lên tiếng, lim dim mắt trong khi
vẫn đang tập trung vào bữa sáng của mình. Mẹ câu trêu chọc: ``{\color{red} Nên để các cô bé ngủ
  thêm chút nữa. Dậy sớm với người lạ chỗ cũng không dễ đâu.}''

Tất cả bốn người trong gia đình khẽ bật cười với lời châm chọc vô hại đó.

\ct

Ba mươi phút sau khi đã ăn bữa sáng. Người bố kiểm lại đồ đạc xem họ còn quên món
đồ gì không, hai anh em thì đang tất bật tìm quần áo mặc đủ ấm trong ngăn tủ đã được
lau dọn. Cả nhà đều bận rộn nhưng mọi việc diễn ra nhanh chóng nhờ thói quen chuẩn bị
sẵn từ những năm trước.

Đúng lúc này, cậu Tổng tranh thủ gửi tin nhắn cho A-chan, người vừa chỉnh sửa xong
những đoạn phim quay hôm trước. Tin nhắn chỉ vỏn vẹn một câu: ``Em gửi cho chị tọa
độ này,'' theo sau là một liên kết dẫn đến vị trí nhà ngoại của cậu ở \ruby{Neihei}{Ninh Bình}, cách đó hơn một trăm cây sô.

Phản hồi từ nàng ``Tăng ca'' xuất hiện ngay lập tức: ``Okkkkkkkkkkkkkkkkkkkk''. Kuon
bật cười nhẹ, đoán chắc rằng cô ấy đã mệt lả sau khi thức cả đêm để làm việc.

\ct

Khi đồng hồ điểm năm giờ sáng, cả gia đình Kuon đã sẵn sàng cho chuyến đi. Họ nhẹ
nhàng rời khỏi nhà, tránh làm phiền các cô gái vẫn còn say giấc. Kuon đi sau cùng, tay
cầm túi đồ, miệng khẽ lẩm bẩm điều gì đó không rõ.

Bên ngoài, Kawachi vào sáng sớm mùng Hai Tết yên tĩnh đến lạ thường. Mặt trời chưa
ló dạng, ánh đèn đường vàng nhạt phủ lên những con phố vắng vẻ. Không khí lạnh buổi
sáng khiến những làn hơn thở thoát ra tạo thành những dai sương mỏng, hòa quyện với
không gian tĩnh lặng.

Ngồi trong xe, gia đình bắt đầu trò chuyện về chuyến đi về quê ngoại. Người bố ngồi bên
cạnh cậu ở ghế sau bắt đầu mô tả kế hoạch chuyến đi: ``{\color{red} Lát nữa đón ông bà ngoại xong,
  ta sẽ chạy thẳng đến Neihei. Dự kiến sẽ tới nơi tầm chín, mười giờ. Chúng ta sẽ ăn sáng
  trên đường tới đó.}''

Mẹ cậu, người cầm lái, nói tiếp: ``{\color{red} Ừ. Sáng nay nhà cậu Mori chắc cũng chuẩn bị bữa đón
  khách. Kuon này, lát nữa con nhớ chào hỏi họ hàng đàng hoàng đấy nha.}''

Kuon thở dài, tỏ vẻ hơi bất mãn: ``{\color{red} Mẹ ơi, nhà ngoại đông người như thế, con làm sao nhớ
  nổi ai với ai. Con còn chẳng phân biệt nổi ông Hiki-san với Niki-san nữa kìa\dots\ Thậm chí
  phần đông trong số họ con còn chưa gặp bao giờ nữa\dots}''

Ryuuhou cười khúc khích: ``{\color{red} Anh chỉ nói vậy vì anh chẳng chịu quan tâm thôi. Năm nào
  họ cũng nhắc đến anh đó.}'' Cậu con trai càng tỏ vẻ khó chịu hơn: ``{\color{red} Nhắc thì cứ nhắc thôi.
  Nhưng sao mọi người lại cứ kỳ vọng vào con thế? Con không phải cuốn bách khoa toàn
  thư biết hết ai với ai đâu\dots}''

Ông bố gằn giọng: ``{\color{red} Đừng cằn nhằn nữa, Kuon. Họ hàng bên ngoại dù sao cũng là người nhà. Lễ Tết, cái quan trọng nhất là giữ lễ nghĩa. Con cũng phải làm quen dần với mọi
  người đi.}''

Người mẹ nói một cách nhẹ nhàng hơn: ``{\color{red} Đúng thế. Dù chỉ gặp vài lần, nhưng họ vẫn coi
  con là cháu. Chẳng lẽ con lại không thể bỏ chút thời gian để nhớ mặt họ sao?}''

Kuon tiếp tục ca thán: ``{\color{red} Đấy là một chuyện mẹ ạ. Con còn chưa nói đến chuyện năm nào
  con về cũng bị họ hàng hỏi về việc có người yêu chưa, rồi có việc làm chưa, điểm số thế
  nào\dots\ Con có phải là `hoa hậu thân thiện' quái đâu mà phải làm vừa lòng mọi người như
  thế?}''

``{\color{red} Bởi vì họ thật sự quan tâm tới con, thế thôi,}'' Ryujin đáp gọn, mắt vẫn chăm chăm vào
màn hình điện thoại. Cậu rất muốn tranh luận tiếp với hai người, nhưng đành bó tay chịu
trận. Năm nào cậu cũng đưa ra vấn đề này trước bố mẹ, để rồi đều bị họ gạt ra một bên.

``\textit{Quan tâm cái con khỉ mốc! Họ toàn moi móc xỉa xói mình thì có!}'' cậu nén chặt sự bực
bội trong lòng, trước khi thở một hơi thật sâu để lấy lại bình tĩnh.

``{\color{red} Nhưng mà\dots\ cậu Mori là ai đấy nhở? Hôm qua mẹ bảo mà con quên béng mất rồi,}'' cậu
cố nhớ lại người mà người mẹ đã giới thiệu từ tối hôm qua. Cả bố mẹ cậu đều chỉ có thể
thở dài ngán ngẩm, cô em gái thì cười trừ: ``{\color{red} Đấy, thấy chưa? Đã bảo anh chẳng quan tâm
  rồi mà.}''

``{\color{red} Ha, ha, ha. Kệ anh,}'' cậu đảo mắt, chẳng thèm để tâm tới những tiếng cười của cô em gái.

Mẹ Kuon bật cười, còn bố cậu chỉ lắc đầu. Không khí trong xe trở nên nhẹ nhàng hơn khi
mọi người tiếp tục trao đổi về lịch trình ngày hôm nay.

\ct

Ba mươi phút sau kể từ khi khỏi hành, xe dừng lại trước nhà ông bà ngoại của cậu ở Toei,
nơi mà hai người đã chờ đợi sẵn trước cổng văn hóa, với tay cầm những túi quà nhỏ gói
ghém cẩn thận.
``{\color{red} Chào ông bà!}'' Ryuuhou nhanh nhẹn bước xuống xe, chạy lại đỡ đồ cho ông bà. Kuon
theo sau, mặt hơi ngái ngủ nhưng vẫn cố nở nụ cười.

``{\color{red} Sớm thế này mà bọn mày đã tới rồi, chắc là phải dậy từ lúc gà còn chưa gáy nhỉ?}'' ông
ngoại cậu - Jitori (\ruby{自}{じ}\ruby{鳥}{とり}, \ruby{Gi}{Bạch}-\ruby{tô-ri}{Điểu}) bật cười.

``{\color{red} Dạ, bọn con cũng phải tranh thủ sớm một chút để kịp giờ ăn sáng ạ,}'' Ryujin trả lời. Aiko
- bà ngoại cậu - hỏi Kuon: ``Thế con đã ăn sáng chưa?''

``{\color{red}Cháu ăn một chút rồi ạ,}'' cậu con trai cười nhẹ, đáp.

Sau khi mọi người lên xe đầy đủ, chuyến đi lại tiếp tục. Từ đây, con đường dẫn về Neihei
dần hiện ra, vòng vèo qua những cánh đồng và những ngọn núi thấp thoáng trong sương
sớm. Trong khi ông bà trò chuyện vui vẻ với bố mẹ Kuon, cậu chỉ lặng lẽ nhìn ra ngoài
cửa sổ, thầm nghĩ về những người họ hàng mà mình chẳng hề nhớ mặt.

\ct

Chuyến hành trình của gia đình cậu sau đó đã xảy ra một cách êm xuôi, không có sự kiện gây rối nào diễn ra cả (có lẽ vì với họ, các cô gái nhà Tam Giác Xanh vẫn chưa thức dậy), chỉ có tiếng cười nói vui vẻ của đại gia đình nhà Hikawa, cùng với
đó là thi thoảng những tiếng thở dài ca thán của cậu Tổng mỗi lần họ nhắc đến họ hàng
nhà ngoại của cậu.

Tám giờ sáng hôm đó, cả gia đình đã đi được nửa chặng đường. Bữa sáng ít ỏi mà cả gia
đình đã ăn hồi sáng sớm là không đủ, nên cả nhà đã rẽ vào một vào quán phở bên đường
để lót dạ.

Mọi thứ vẫn không có gì ngạc nhiên xảy ra, cho đến khi\dots

``{\color{red}Cho cháu năm bát phở tái chín ạ\dots}'' Kuon tiến tới trước quầy bán và gọi suất ăn, tỏ vẻ
hơi mệt mỏi.

Lập tức, giọng nói của một cô gái trẻ cất lên làm cậu khẽ giật mình: ``{\color{red} Có ngay!}''

Đó là một cô gái với đôi tai thú trắng nhọn, mái tóc trắng dài được tết lệch, đôi mắt xanh
lam như phản chiếu cả bầu trời, và một chiếc đuôi vẫy nhè nhẹ sau lưng, trên đó đính một
ngôi sao nhỏ lấp lánh. Cô gái mặc bộ kimono trắng, màu sắc dần chuyển sang tông xanh
khi xuống dưới, phối cùng váy hakama cạp cao màu xanh và đôi bốt buộc dây đen\footnote{Về chi tiết bộ trang phục mà Fubuki đang mặc, xem hình (c), trang \pageref{setone}.}.

Cô ấy trông rất quen thuộc. Nếu không có chiếc mặt nạ đen từ Comiket\footnote{Xem {\flaregothic ÉPISODE 13}, quyển số Một.} mùa hè năm
ngoái, Kuon chắc chắn đã nhận ra ngay. Cậu nhíu mày, lập tức gọi lại: ``{\color{red} Ơ này\dots}''

Nhưng cô gái chỉ quay lưng đi, nhanh chóng vào nhà bếp, bỏ mặc Kuon với một loạt câu
hỏi chưa kịp cất thành lời. Cậu ngồi xuống bàn cùng gia đình, miệng lẩm bẩm: ``\textit{Nghi
  lắm đây\dots}''

Cậu bước về bàn ăn của năm người, mắt ngước nhìn lại chiếc sạp mà cậu vừa gọi đồ, cảm
thấy có gì đó không ổn.

``{\color{red} Có gì đó không ổn sao, Kuon?}'' Kaien nhẹ giọng hỏi khi thấy cậu đang ngơ ngác nhìn về
phía sạp. Kuon giật mình, quay về phía mẹ đáp: ``{\color{red} D-Dạ, không có gì ạ.}''

Mẹ cậu thấy vậy cũng chỉ gật đầu, không hỏi thêm cậu điều gì nữa.

\ct

Chưa đầy năm phút sau, một cô bồi bàn khác bước ra, bưng năm bát phở tái chín đặt từng
bát một lên bàn của nhà Kuon. Cô gái đó thốt lên: ``{\color{red}Đồ ăn của mọi người đây!}''

Kuon đưa mắt nhìn cô bồi bàn mới đến. Lại một khuôn mặt rất quen thuộc. Mái tóc nâu
cắt ngắn, điểm thêm một vệt xanh nổi bật trước mái, cùng bộ vu nữ cách điệu mang nét
hiện đại\footnote{Về chi tiết bộ trang phục mà Matsuri đang mặc, xem hình (c), trang \pageref{setthree}.}. Cô gái này đi dép zouri đỏ và đeo một chiếc mặt nạ màu cam, giống với chiếc
mặt nạ mà cô đã đeo khi ``đóng vai'' kẻ phản diện.

Vừa đặt bát cuối cùng xuống, cô ấy khẽ cúi đầu: ``{\color{red}Chúc cả nhà ngon miệng!}'' và rồi hớn
hở bước ra chỗ sạp.
Cậu Tổng chăm chú nhìn bát phở mới bưng của mình, rồi bất giác buột miệng, nở một nụ
cười thắm: ``Cảm ơn em, \uline {Matsuri-chan.}''

Vừa dứt lời, cô gái bồi bàn đứng khựng lại, đôi tay vẫn còn đặt trên bàn. Cả cơ thể cô
như cứng đờ trong vài giây. ``Mình biết mà,'' cậu nghĩ thầm, cười mỉm khi nhìn nàng Lễ
hội kia đứng im bặt lại.

Kuon chưa kịp phản ứng gì thêm thì nghe thấy một tiếng thì thầm khe khẽ từ phía sau, đủ
để cậu nhạy bén nhận ra. Giọng điệu ấy rất giống Haato: ``\hl{Thấy chưa, tôi đã bảo kiểu gì
  Kuon-senpai cũng sẽ phát hiện ra mà.}''

Matsuri bối rối, mặt hơi đỏ lên, lắp bắp: ``S-Sao anh lại\dots?'' Kuon nhếch miệng cười, chỉ
vào chiếc mặt nạ và chậm rãi đáp: ``Biết đúng không? Anh chỉ biết một người có trang
phục kỳ quặc như vậy thôi. Chẳng có ai ngốc đến mức mặc như thế vào mùa đông lạnh
đến 16\degree C cả!''

Matsuri cười khì, cố gắng tỏ ra bình thản: ``L-Lộ rồi à\dots'' Đúng lúc này, Fubuki - người
đã ghi đơn cho gia đình Kuon ban nãy - cũng đi ra, gãi đầu ngượng ngùng: ``Em không
nghĩ là anh có thể đoán được đấy\dots''

Kaien nhìn hai cô nàng với vẻ ngạc nhiên, đồng thời không giấu khỏi sự thích thú: ``Ủa,
Matsuri-chan đấy à? Cháu đến từ khi nào vậy?''

Matsuri quay sang cúi đầu chào mẹ của Kuon, giọng lễ phép hơn hẳn: ``Từ nửa tiếng trước cô chú với ông bà của Kuon ạ\dots''

Ryujin, từ nãy giờ giữ vẻ mặt bình thản, đột ngột nói: ``Chú cũng ngạc nhiên khi Fubuki-chan
cũng có mặt ở đây đấy.''

Kuon quay phắt lại, nheo mắt nhìn về phía đằng sau tấm ván ở gần sạp bếp: ``Con cũng
đang định nói này\dots\ Và cả Haato-chan với Ayame-chan ở phía sau nữa.''

Ngay khi bị chỉ mặt gọi tên, hai người bước ra từ phía bếp và góc khuất sau nhà, lấp ló
như thể đang lén lút làm gì đó. Cả nàng Quỷ sừng và nàng Song tính đều cười, giơ hai
tay lên trời, như thể họ vừa bị bắt tại trận.

``Trời đất, cả gia đình anh đoán chuẩn quá, yo!'' Ayame tạo dáng, nhảy dang tay trước mặt
mọi người, trong khi Haato thì chỉ biết vẫy tay chào.

Cậu con trai khoanh tay lại, nhìn cả bốn người: ``Mấy đứa dậy từ khi nào vậy?'' Nàng
Song tính đáp nhanh: ``Lúc sáu giờ sáng. Bọn em có sang nhà anh gọi dậy nhưng không
thấy ai trong gia đình anh ở nhà cả.''

Cậu nhún vai, cười mỉa: ``Thì, bọn anh đi từ sớm mà,'' rồi bất chợt khựng lại, tỏ vẻ hơi
ngạc nhiên: ``Mà mấy đứa vào nhà bằng cách nào vậy?''

Bố cậu húp nước phở, thản nhiên đáp: ``Bố có đưa cho họ chìa khóa nhà sơ cua đấy.''
Kuon thở dài, tặc lưỡi: ``\textit{Chẳng trách\dots\ Tự dưng lại đưa cho họ chìa khóa nhà làm chi
  không biết\dots}''

``Cụ thể là thế nào vậy, Haato-onee-chan?'' Ryuuhou nghiêng đầu hỏi. Thế là, nàng Trung
học bắt đầu kể lại hồi tưởng của cô cho mọi người.

\begin{center}
  \uline{Hồi tưởng\\Sáu giờ sáng.}
\end{center}

Haato cùng Ayame đến nhà Kuon với ý định gọi mọi người dậy. Nàng Song tinh mở cửa
bằng chìa khóa dự phòng mà Ryujin đưa từ hôm trước. Những tiếng lạch cạch của cửa
cuốn vang lên như phá bĩnh đi khoảng lặng yên ắng buổi sớm mai.

Bước vào nhà, họ chạy lên tầng năm, phòng của hai bố con Kuon để gọi họ dậy trước tiên.
Cảnh tượng khiến cả hai phải dừng chân lại ngay khi bước vào là bầu không khí yên tĩnh
kỳ lạ, chỉ có tiếng thở đều đều từ trong phòng.

Nàng Oni khẽ nghiêng đầu, hỏi nhỏ: ``Ủa, sao không thấy ai ngoài này hết vậy?'' Haato
cẩn thận bước vào, ngó nghiêng xung quanh, cuối cùng phát hiện ra một nhóm người đang
ngủ rải rác trong phòng khách. Họ cũng chính là bảy cô gái đã chiến thắng cuộc thi làm bánh
mochi từ sáng hôm qua.

``Chắc họ không biết nhà Kuon-senpai đi từ sớm rồi,'' Ayame thì thầm. Nhìn thấy cảnh
này, nàng Trung học nhún vai, bước đến gần. Cô quay lại chỗ nàng Quỷ sừng, nháy mắt:
``Ayame-chan, chiêng đâu?''

Ayame cười ranh mãnh, lôi từ ba-lô ra một cái chiêng nhỏ cùng hai dùi gỗ. ``Sẵn sàng
chưa?''

Haato gật đầu, ra hiệu, và ngay lập tức\dots

\textbf{*CHENG! CHENG! CHENG!*}

\dots nàng Oni bắt đầu đánh chiêng ầm ĩ, giống như cách mà cô gọi Kuon dậy vào buổi sáng
ngày mùng Một. Cô nàng hét lớn: ``\textbf{Thanh la kêu tiếng rất vang, cheng cheng cheng,
  cheng, cheng, cheng!\footnote{Lấy từ bài \textit{Cộc Cách Tùng Cheng} của Phan Trần Bảng.} Dậy đi nào mọi người ơi! Yo, yo!}''

Cả bảy cô gái giật mình tỉnh dậy, mặt mũi ngơ ngác, mỗi người một kiểu. Pekora hét lên,
hai tay bịt lấy đôi tai thỏ rung rinh của mình: ``\textbf{Pekooooo\~{}!! Cái gì thế hả!?}''

Okayu dụi mắt, ngáp một cái rõ dài: ``Gì vậy\dots\ Sớm thế này mà ồn thế?'' Korone lẩm
bẩm, đầu của cô vẫn còn nửa tỉnh nửa mê: ``Ayame-chan, làm gì mà gấp thế\dots''

Nàng Song tính chống tay lên hông, cười tươi: ``Dậy đi, mấy người! Nhà Kuon-senpai đi
từ lâu rồi đấy!''

Noel nhìn quanh, không hiểu ý cô nói: ``Đi rồi? Đi đâu cơ?'' Rushia, người cũng vừa tờ
mờ dậy, cũng nghiêng đầu: ``Kuon-senpai\dots\ đi về quê rồi sao?''

Haato lấy điện thoại ra, đọc to tin nhắn mà Kuon gửi cho A-chan từ trước đó: ``Anh ấy có
viết: `Em gửi cho chị tọa độ này,' kèm theo đó là vị trí quê ngoại của ah ở tỉnh Neihei.
Chắc chắn họ đã đi thẳng đến đó rồi.''

Miko nhảy dựng lên, vừa cuống cuồng tìm bộ kimono mà cô để tạm, vừa nói: ``Vậy là họ
bỏ chúng ta lại mà không nói gì ne!? Kuon-senpai xấu tính!''

Roboco mỉm cười, vỗ nhẹ vào vai nàng Vu nữ: ``Chắc họ không muốn làm phiền thôi.
Giờ thì chuẩn bị đi!''

Nàng Pháp sư cũng động viên mọi người: ``Có khi gia đình anh vẫn chưa đi xa đâu!'' Và
roofi, không ai bảo ai, mọi người cuống cuồng thu dọn phòng ốc, người thì chỉnh lại quần
áo, người thì chạy ra ngoài xe để kiểm tra.

Rushia vừa quấn chăn vừa lẩm bẩm: ``Gọi dậy kiểu gì mà muốn đau tim luôn\dots\ Lần sau
đừng gọi như thế nữa được không-nanodesu?''

\ct

Sau khi ổn định, họ chạy xuống dưới nhà, với những thành viên còn lại đã chờ đợi sẵn,
dù mặt ai cũng có vẻ vừa ngái ngủ vừa vội vã.

Ở bên trong xe, Okayu thong dong bước vào trong xe, nhìn đồng hồ và ngáp dài: ``Oa\~{}
*ngáp* Cũng tại mình không nghe báo thức\dots\ Nhưng mà, đi gấp thế này thì ăn sáng kiểu
gì đây?''

Mio, ngồi ở xe bên cạnh, hối hả đưa mọi người lên xe, tỏ vẻ uy nghiêm: ``Không kịp ăn
uống gì đâu! Tất cả nhanh chóng lên xe đi!'' rồi quay sang nàng Đeo kính.

``A-chan, phiền chị chỉ đường cho bọn em!'' Mio nói, hối hả đưa mọi người lên xe, nàng
Đeo kính giơ một ngón tay cái lên, tư thế sẵn sàng chiến đấu, mặc do đôi mắt cô thâm
quầng vì thiếu ngủ.

Miko chạy qua, níu tay A-chan: ``Chúng ta sẽ đuổi kịp chứ, ne!? Đường vắng thế này thì
xe anh ấy chạy nhanh lắm đấy!'' Nàng Sói đen đáp: ``Đuổi được! Ta có hai xe mà, chia
nhau ra vừa đuổi vừa phối hợp!''

Fubuki cầm lái, giơ tay lên, cười một cách tự tin: ``Xe của tôi chạy nhanh lắm, mọi người
cứ tin ở tôi!'' Okayu ở xe bên cạnh, nhún vai cười nhạt: ``Xe của tôi cũng có thể lái được
đến Mach mừ\~{} Fubuki-senpai, nhớ đi đường bìn anh nha\~{}''

``Yên tâm, chị có tính toán cả rồi. Chúng ta sẽ đến đúng lúc, cứ chờ xem,'' nàng Cáo trắng
cười lớn.

\ct

Cả hai xe bắt đầu lăn bánh trên con đường vắng vẻ của Kawachi. Do đường phố lúc sáu
giờ sáng vẫn còn vắng, nên cả hai chiếc xe có thể chạy được với tốc độ tối đa.

A-chan ngồi ghế phụ trên xe Okayu, tay chỉ bản đồ trong điện thoại, mắt lờ đờ nhưng vẫn
cố gắng giữ tỉnh táo: ``Đi đường này trước, rồi rẽ trái ở giao lộ tiếp theo\dots\ À không, tôi
nhầm. Rẽ phải mới đúng.''

Mio nhìn cô nàng đang ngái ngủ qua gương chiếu hậu, nhắc nhở: ``Chị có thể tỉnh táo một
chút được không? Chúng ta không muốn bị lạc đâu.''

Nàng Đeo kính thở dài, chỉnh lại kính: ``Tôi đang cố hết sức rồi đây\dots''

Ở bên xe của Fubuki, cô nàng lái xe với khoảng cách vừa đủ với xe của Mio. Qua bộ đàm nội bộ trên chiếc xe cảnh sát, cô nói với Matsuri: ``Matsuri-chan, tôi có ý này\dots\ Tới lúc
chúng ta tạo bất ngờ rồi đấy.''

Nàng Lễ hội lập tức hiểu ý, gật đầu: ``Ừ! Tới quán ăn nào đó trước, cải trang một chút để
không ai nhận ra. Mọi người sẽ ngạc nhiên cho mà xem!''

Cáo trắng cười nhẹ, tập trung lái xe: ``Vậy thì bám sát vào, đừng để họ nghi ngờ. Tới lúc
ăn sáng, chúng ta sẽ hành động.''

Họ không mất quá nhiều thời gian để bắt kịp được chiếc xe đen của gia đình Kuon. Ngay
khi chiếc xe của cậu đang bắt đầu tìm chỗ đỗ để vào quán phở, cả Fubuki và Matsuri đã
chuẩn bị sẵn kế hoạch của mình.

Họ luồn ra cửa sau, nơi mà ít người biết, rồi lục lọi ra được hai chiếc mặt nạ - một đen,
một cam. Cáo trắng đưa cho đối phương chiếc mặt nạ cam và bảo: ``\hl{Bà ở ngoài phục vụ,
  còn tôi thì vào trong bếp. Chúng ta làm cho khéo vào, đừng để lộ sơ hở.}''

Nàng Lễ hội gật đầu, cười tinh quái: ``Yên tâm, tôi làm bồi bàn chuyên nghiệp lắm!''

Và những gì sau đó đã xảy ra như đã thấy ở trên: Kuon gọi đồ, được Fubuki chốt đơn, rồi
Matsuri là người bưng phở cho cả gia đình cậu.

\begin{center}
  \uline{Kết thúc hồi tưởng.}
\end{center}

Cả bốn cô gái đứng trước cậu Tổn, cười ngượng. Haato chống cằm, tỏ vẻ tự hào: ``Đấy,
bọn em đã dậy sớm thế mà vẫn bị anh bỏ lại. Chắc anh muốn chơi khó bọn em rồi á.''

Ayame cười khì: ``Đừng tưởng anh có thể thoát khỏi được bọn em đâu, yo!'' Kuon chỉ có
thể lắc đầu, thở dài: ``Anh chẳng chơi khó gì cả, mấy đứa lề mề thì có\dots''

Thế rồi, cả bốn cô gái cùng trò chuyện vui vẻ với gia đình nhà Kuon, làm bữa sáng trở
nên náo nhiệt hơn hẳn. Sau nửa tiếng đồng hồ, mọi người ăn xong, thanh toán tiền cho
chủ quán, và chuẩn bị rời đi.

\ct

Khi cả nhà vừa bước ra khỏi quán, họ bất ngờ thấy các thành viên còn lại của \hll\ đang ngồi lấp ló trong hai chiếc xe đỗ gần đó, như thể đang ``mai phục''. Một vài người
còn lấy áo khoác che mặt, nhưng chẳng giấu được sự tò mò lóe lên trong ánh mắt.

Kuon nhìn họ chằm chằm, khuôn mặt không cảm xúc, rồi thở dài: ``Trông mấy đứa giống
mấy tên tội phạm dòm ngó thật.''

Pekora thò đầu ra khỏi cửa sổ, lập tức cất giọng châm chọc: ``Kuon-senpai à, anh bỏ bọn
em lại phía sau mà không một lời nhắn nhủ là thế nào hả, peko?''

Rushia từ ghế lái phụ nhảy xuống, bước nhanh về phía cậu, chống nạnh bất bình: ``Đúng
đấy! Anh đi mà không thèm đợi ai cả. Bộ anh nghĩ mình là nhân vật chính chắc?''

Cậu con trai nhún vai, giọng thản nhiên: ``Anh có nhắn A-chan rồi đấy thôi. Với cả, chẳng
phải mấy đứa cũng bắt kịp được xe bọn anh đấy thôi?''

Nàng Pháp sư khoanh tay, mặt vẫn hậm hực: ``Nhờ tài lái xe của Okayu-senpai và Fubuki-senpai
thôi đấy! Lần sau đừng có mà trốn đi như thế!''

Mio đứng cạnh Mel và Sora, chỉ biết cười trừ trước màn quở trách của hai cô gái. Nàng
Dơi mỉm cười nhẹ, nhỏ giọng với nàng Sói đen: ``Vẫn chưa hết ngái ngủ mà mọi người
đã ồn ào thế này rồi\dots''

Nàng Thần tượng gật đầu, ôn tồn nói: ``Thôi nào, dù gì chúng ta cũng có mặt đầy đủ rồi.
Đừng cãi nhau nữa, tiếp tục hành trình thôi!''

Từ phía ghế sau của một trong nững chiếc xe, Subaru ló đầu ra, cười vang: ``Đúng đó! Cãi
nhau làm gì khi chúng ta còn một chặng đường dài phía trước. Đi thôi, Kuon-senpai!''

Ayame, đang ngồi ngay cạnh Subaru, khẽ thở dài, cười khì: ``Cũng tại Pekora-chan với
Rushia-chan cứ làm quá lên đó, yo.''

Nàng Cáo trắng trèo lên xe, đuôi vẫy, nở một nụ cười tươi: ``Dù sao thì mọi người cũng
bắt kịp được anh ấy rồi. Chúng ta đi tiếp chứ? Đường phía trước chắc sẽ thú vị lắm đấy!''

Nhìn các cô gái cười đùa châm chọc vui vẻ, ông bà ngoại của Kuon khẽ bật cười. Jitori
nhìn nhóm người, nói với bà xã của mình: ``{\color{red} Tuổi trẻ ngày nay đúng là sôi động thật.}''

Aiko mỉm cười: ``{\color{red} Quả đúng là một nhóm người người kỳ lạ. Nhưng mà trông cũng vui
  nhỉ, con?}'' Kaien gật đầu: ``{\color{red} Dạ. Nhưng mà họ nói chuyện rôm rả quá, con cũng hơi bất
  ngờ\dots}''

Bà ngoại hướng về phía cậu con trai bằng ánh mắt dịu dàng, nở một nụ cười nồng thắm:
``{\color{red} Đông vui như vậy thì chẳng khác gì một đại gia đình cả. Kuon, con thật may mắn khi
  có những người bạn như họ đấy.}''

Ryujin đứng bên cạnh, nhìn các cô gái rồi thích thú nói: ``{\color{red} Quả là những người kỳ lạ.
  Nhưng mà trông cũng bám lấy thẳng Kuon ra phết đấy.}''

Kuon chỉ có thể thở dài với những lời ngợi khen ``có cũng như không'' từ họ. Ryuuhou,
người vẫn đang kiểm đếm lại xe trước khi khởi hành, quay lại nói lớn: ``Được rồi, mọi
người! Lên xe đi, còn nhiều nơi phải đến nữa!''

Nàng Vu nữ trên xe vẫy tay, cất giọng hối thúc trong vui vẻ: ``Đi thôi nào, mọi người ơi!
Kuon-san! Cẩn thận không là tắc đường đấy ne\~{}!''

Kuon nhíu mày, khẽ chẹp miệng: ``Biết rồi\dots''

Cả đoàn xe tiếp tục lăn bánh trên đường. Tiếng cười nói, chọc ghẹo vang vọng trong
không gian hẹp của xe, bầu không khí càng trở nên náo nhiệt hơn khi tất cả các thành viên
đã tập hợp đầy đủ.

\ct

\pov{Hikwa Kuon}

Đúng như Miko-chan dự đoán, đường cao tốc về quê vào những ngày này càng lúc càng
đông. Những chiếc xe nối đuôi nhau thành một hàng dài bất tận, và các biển số đến từ
khắp nơi trong nước làm tôi cảm giác như đang xem một triển lãm xe di động vậy. May
mắn là, dù đông đúc, dòng xe vẫn chưa đến mức kẹt cứng, nên từ quán phở đến nhà đầu tiên - nơi mà chúng tôi phải rẽ qua để chúc Tết trước khi tiếp tục hành trình - chỉ mất
khoảng hai mươi phút.

Nhắc đến kế hoạch hôm nay, phải nói rằng nó khá\dots\ phức tạp. Chúng tôi sẽ đi chúc Tết
họ hàng nhà ngoại trước, cụ thể là các gia đình ở sống rải rác trong cùng một tỉnh. Sau đó,
cả đoàn mới tiếp tục lên đường về quê chính của nhà ngoại tôi. Nghe thì có vẻ dễ dàng,
nhưng với số lượng người hiện tại - bao gồm gia đình tôi và toàn bộ dàn \hll\ - thì tôi
gần như có cảm giác mình đang tổ chức một chuyến lưu diễn hơn là một buổi chúc Tết
gia đình.

Và ơn giời, ƠN GIỜI, ít nhất là vì đường đông nên hai chiếc xe của Okayu-chan và
Fubuki-chan - vốn rất có xu hướng bày trò nghịch ngợm - chẳng thể làm được điều gì dại
dột. Điều này khiến tôi nhớ đến hôm qua, khi hai chiếc xe đó đã biến cả thành phố thành
trường đua trong một bộ phim hành động. Kết quả là một số người đi đường chắc hẳn
đã nghĩ rằng họ vừa vô tình lọt vào một chương trình thực tế. Tôi chỉ còn biết cầu mong
ngày hôm nay sẽ yên bình hơn.

Sau khoảng hai mươi phút, xe chúng tôi dừng lại trước một căn nhà nhỏ cổ kính nhưng
khang trang, thuộc về một cụ bà đã ngoài tám mươi - mà thú thật, tôi chẳng thể nhớ nổi
bà ấy tên gì. Điều này không có gì lạ, vì tôi vốn chưa bao giờ thực sự quan tâm đến việc
nhớ tên hay mặt bất cứ ai trong họ hàng nhà ngoại (thậm chí là cả nhà nội). Đối với tôi,
họ chỉ là những cái tên xuất hiện thoáng qua trong những buổi tụ họp mà tôi bị bắt buộc
phải tham gia, để rồi khi đi về là chẳng hề đọng lại một ấn tượng đáng nhớ nào với tôi.

Cụ bà chào đón gia đình tôi với nụ cười hiền hậu và những lời hỏi han đầy tình cảm.
Nhưng khi ánh mắt cụ lướt qua dàn người lố nhố đứng bên ngoài - những cô gái với trang
phục rực rỡ nhưng mang phong cách\dots\ chẳng giống ai - cụ hơi nhíu mày, rồi hỏi mẹ tôi:
``{\color{red} Những người ngoài kia là ai đây?}''

Mẹ tôi mỉm cười, giọng ôn tồn như thường lệ: ``{\color{red} Đồng nghiệp của thằng Cò ấy mà.}'' Cụ
bà nhìn tôi, rồi quay lại nhìn đám người bên ngoài đang vẫy tay chào, vẻ mặt pha lẫn giữa
sự ngạc nhiên và bối rối: ``{\color{red} Sao lại đông thế?}''

Tôi cố nhịn cười, nhưng cuối cùng cũng bật ra một câu đầy hài hước: ``{\color{red} Đó là cả công ty
  cháu ạ.}''

Ngay khi câu trả lời rời khỏi miệng tôi, cụ bà sửng sốt đến mức suýt ngất. Tôi biết, tài ăn
nói của tôi tệ hại thế nào rồi mà.

Ông bà ngoại phải nhanh chóng đỡ cụ vào nhà và trấn an cụ rằng đó chỉ là một cách nói
vui. Nhưng trong lòng, tôi biết rõ rằng cụ vẫn đang cố gắng hiểu tại sao tôi lại dẫn theo
gần hai mươi người, mỗi người một vẻ, vào nhà chúc Tết.


Trong lúc cụ bà được dìu vào, tôi quay sang mẹ, thấp giọng hỏi: ``{\color{red} Ờm\dots\ Con nói vậy có
  hơi quá không ạ?}'' Mẹ tôi lắc đầu, mỉm cười nhẹ: ``{\color{red} Không sao đâu. Chỉ là cụ không quen
  với việc gặp nhiều người lạ thôi.}''

``{\color{red} Cũng chẳng trách cụ ấy được. Không sao đâu con,}'' bố tôi vỗ vào vai tôi, trước khi cùng
họ bước vào của nhà cụ bà kia.

Việc tôi dẫn cả ``công ty'' về quê ăn Tết có lẽ là điều chưa từng xảy ra trong lịch sử dòng
họ này, nên việc cụ ấy sốc tới mức suýt ngất là điều mà tôi có thể đoán trước.

Dàn người bên ngoài vẫn đứng chờ. Peko-chan và Ru-chan bắt đầu than phiền (có lẽ vì
tôi tự thoại lâu quá):

\begin{dialogue}
  \speak{Pekora} Sao lâu thế, Kuon-senpai? Mấy cái lễ nghĩa này mệt chết đi được peko\~{}!
  \speak{Rushia} Đúng đó! Em có cảm giác như anh cố tình kéo dài thời gian để thử lòng
  kiên nhẫn của bọn em vậy!
\end{dialogue}

Tôi chỉ cười nhạt. Những lời than vãn như vậy, tôi đã nghe không biết bao nhiêu lần từ
ngày đầu tiên làm việc với họ.

Ở phía xa, Mio-chan, Mel-chan và Sora-chan vẫn giữ thái độ điềm tĩnh, như mọi khi. Tôi
nhìn thấy Sói đen-chan nói nhỏ với Thần tượng-chan: ``Dù sao thì hôm nay cũng là ngày
vui. Cứ để họ ồn ào một chút cũng không sao.''

Nàng Dơi gật đầu đồng tình: ``Đúng thế. Nhưng chị nghĩ Kuon-senpai chắc cũng mệt lắm
rồi\dots''

Cả đoàn cuối cùng cũng chuẩn bị để vào nhà. Tôi hít một hơi thật sâu, bước đi trước để
dẫn đầu, tự nhủ rằng đây chỉ mới là \uline{nhà đầu tiên }trong danh sách dài dằng dặc của hôm
nay.

\ct

Sau một hồi được ông bà ngoại và bố mẹ tôi trấn an, cụ bà cuối cùng cũng bình tĩnh lại.
Như mọi năm, cụ rút ra những bao lì xì nhỏ đã chuẩn bị sẵn, với vẻ mặt tươi cười và giọng
nói ấm áp: ``{\color{red} Cả nhà mỗi người một bao nhé.}''

Cụ đưa từng bao lì xì cho gia đình tôi, nhưng khi nhìn sang các cô gái idol - mà tôi thầm
gọi trong đầu là những ``diễn viên hài ngoài giờ'' - cụ ngập ngừng. Cũng phải thôi, chẳng
ai chuẩn bị đủ lì xì cho cả một đoàn người đông đúc như vậy cả. Cụ đành phải nhẩm lại
số bao lì xì mà cụ có trong tay, trong khi chúng tôi chỉ có thể cười trừ.

``{\color{red} Bọn cháu không cần lì xì từ cụ đâu ạ!}'' Sora-chan nói với giọng nói nhẹ nhàng. Hầu
như tất cả các cô gái chúng tôi đều cùng đồng loạt gật đầu, thể hiện rằng họ có đủ tiền để
trang trải cuộc sống.

Tuy nhiên, khi cụ vừa kịp thở phào vì nghĩ rằng chỉ cần lì xì cho gia đình tôi là đủ, thì gần
như toàn bộ các thành viên \hll\ đã đồng loạt rút ra bao lì xì của mình và mừng tuổi
lại cụ.

Mio-chan là người tiên phong, bước lên trước với nụ cười dịu dàng: ``{\color{red} Cháu mừng tuổi cụ
  ạ.}'' Cụ bà bối rối, nhưng rồi cũng nhận lấy: ``Ừ, cụ xin\dots''

Ngay sau đó, Mel-chan tiến lên với giọng nhỏ nhẹ nhưng đầy nhiệt tình: ``{\color{red} Cháu mừng
  tuổi ạ!}'' Cụ bà, dù đang cầm một bao lì xì trong tay, vẫn vội gật đầu: ``{\color{red} Ừ, ừ\dots}''

Subaru-chan, Haato-chan, và Noel-chan cũng không chịu thua, đồng loạt giơ tay ra:

\begin{dialogue}
  \speak{Subaru} {\color{red} Cả cháu nữa!}
  \speak{Haato} {\color{red} Cháu cũng vậy ạ!}
  \speak{Noel} {\color{red} Đây là tấm lòng của cháu!}
\end{dialogue}

Cụ bà chỉ biết đáp lại bằng một câu ngắn gọn: ``{\color{red}Rồi, rồi\dots\ cứ bình tĩnh\dots}''

Thậm chí đến một người rất thích tiêu tiền vào trò chơi may rủi như Fubu-chan cũng nhiệt
tình lì xì cho cụ ấy. Đó là điều đáng khen.

Khung cảnh lúc này giống như một phiên chợ Tết hỗn loạn. Các cô gái lần lượt tiến lên,
khiến cụ bà không kịp trở tay. Những bao lì xì cứ thế chồng chất trên tay cụ, đến mức tôi
thậm chí còn lo cụ sẽ phải nhờ ông bà ngoại hoặc mẹ tôi đứng ra nhận hộ.

Riêng chỉ có Shion-chan đứng khoanh tay một góc, nhìn mọi người với ánh mắt tò mò
pha chút châm chọc: ``Trông mọi người có vẻ nhiệt huyết nhỉ.''

Tôi quay sang nhìn cô bé, nhún vai: ``Đương nhiên rồi. Có bố con anh dạy cho mà. Đúng
hơn là Mel-chan, Sora-chan và Mio-chan đã huấn luyện kỹ lưỡng. Nhưng mà Fubu-chan
cũng bắt chước được như họ thì mới là điều anh thấy ngạc nhiên.''

Phù thủy-chan cười khẩy, không mảy may động lòng: ``Nhưng lì xì chẳng phải là mất tiền
sao?'' Tôi thở dài, nheo mắt nhìn em: ``Về lý thuyết là như vậy. Nhưng em hỏi thế có
nghĩa là em không muốn làm à?''

Em ấy nhún vai, giọng thản nhiên: ``Thôi, khỏi cần.'' Nghe câu trả lời theo kiểu ``tôi cóc
quan tâm'', tôi chi có thể ngao ngán: ``Haiz, tùy em thôi. Kiệt xỉn vừa thôi chứ.''

Rồi tôi quay sang những người còn lại, giọng nói lớn hơn để át đi tiếng nháo nhào: ``Mọi
người cứ bình tĩnh! Từng người một thôi, không cụ ấy ngộp thở bây giờ!''

Ông bà ngoại tôi và bố mẹ chỉ đứng nhìn cảnh tượng này, không nhịn được cười gượng.
Mẹ tôi thì thầm với bố: ``Có vẻ như mấy đứa ấy hăng quá nhỉ.''

Tôi thầm nghĩ trong lòng: ``\textit{Quá? Đây là quá lắm luôn ấy chứ!}'' Nhưng nói đi cũng phải
nói lại, sự hào hứng của các cô gái cũng khiến tôi cảm thấy bầu không khí Tết thêm phần
vui vẻ.

Cụ bà sau khi nhận hết bao lì xì cuối cùng từ dàn \hll\ thì thở dài một hơi, nhưng
khuôn mặt vẫn ánh lên niềm vui: ``{\color{red}Các cháu đúng là làm cụ bất ngờ. Nhưng mà\dots\ cảm
  ơn nhé.}''

Mio mỉm cười dịu dàng, cúi đầu cảm ơn cụ: ``{\color{red}Dạ, đó là tấm lòng của bọn cháu thôi ạ. Cụ
  nhận là bọn cháu mừng rồi.}''

Tôi chỉ biết nhìn cảnh này mà cười khổ. Đúng là khi làm việc với nhóm người này, không
ngày nào là thiếu những bất ngờ khó đỡ.

\ct

Sau khi rời khỏi nhà cụ bà, cả đoàn tiếp tục chuyến hành trình với hai nhà nữa. Lần này,
rút kinh nghiệm từ ``sự kiện lì xì hỗn loạn'' ban nãy, tôi quyết định chọn cách tiếp cận
khác.

Theo lời của Shion-chan - người luôn khôn lỏi trong việc tránh những thứ tốn sức - cũng
như theo sự quan sát của bản thân tôi, Sora-chan, Mio-chan, và Mel-chan - những người
có vẻ điềm tĩnh và được các cụ yêu quý nhất - được tôi đại diện làm nhiệm vụ lì xì.

Trước khi đến nơi, tôi cố nài nỉ cô nhóc ``kusogaki'' Shion chịu làm giống như mọi người.
Cô nàng ban đầu nhún vai từ chối, mắt liếc xéo tôi: ``Em thấy chẳng ích gì cả. Mất tiền
thì có.''

Tôi nhún vai, cười nham hiểm: ``Thế thì em thử đi. Biết đâu em lì xì người ta, họ lại lì xì
em nhiều hơn thì sao?''

Shion-chan ngập ngừng, vẻ nghi ngờ hiện rõ trên mặt. Nhưng sau một hồi lưỡng lự, em
cũng quyết định thử. Nào ngờ, đúng như tôi dự đoán, khi Shion -chan đưa lì xì ra, các
bác trong nhà mỉm cười hiền từ, rút lì xì lại cho em. Đôi mắt của em ấy sáng lên khi thấy
số tiền lì xì nhận được nhiều hơn hẳn số mình bỏ ra.

Cô nàng quay sang tôi, tay xòe ra ba phong bao lì xì, giọng điệu pha chút ngạc nhiên lẫn
thán phục: ``Kuon-senpai\dots\ Anh đúng là cáo già thật đấy.''

Tôi cười khẩy, đáp lại: ``Đừng quên ai là senpai của em nhé. Giờ thì tiếp tục công việc đi.''

Tuy nhiên, mọi thứ không phải lúc nào cũng suôn sẻ. Khi vào đến nhà thứ hai, một vài
người họ hàng của tôi - mà thú thực tôi chẳng nhớ tên hay mặt mũi ai, cũng chằng phải là
điều mà tôi bận tâm tới - bắt đầu đặt ra những câu hỏi mà chỉ nghe thôi đã muốn thở dài.

\il

Một người phụ nữ trung niên, có lẽ là một trong các dì (về lý thuyết) của tôi, cười giả lả
hỏi: ``{\color{red} Cháu Kuon, năm nay đã có người yêu chưa? Có bạn gái chưa?}''

Tôi nheo mắt nhìn bà ấy, cố gắng giữ nụ cười xã giao: ``{\color{red} Cháu đang bận học và làm việc,
  cũng chưa nghĩ đến chuyện đó ạ.}'' Dù vậy, bà ấy vẫn không buông tha: ``{\color{red} Thế á? Nói như
  cháu á, chắc chẳng ai thèm đâu đấy hở? Kén chọn quá đó!}''

Đó THỰC SỰ là câu trả lời làm tôi muốn tức cái đầu. Tuy nhiên, ngay khi tôi định đáp
lại thì Mio-chan bước tới, mỉm cười dịu dàng nhưng lời nói thì sắc bén hơn hẳn: ``{\color{red} Dạ,
  chuyện tình cảm là chuyện riêng tư, mong dì đừng làm khó Kuon-senpai ạ.}''

``{\color{red} Với cả Kuon-senpai có bạn gái hay không cũng không phải là nghĩa vụ mà dì phải biết
  đâu ạ,}'' Fubuki-chan tiếp lời.

Người phụ nữ kia thoáng đỏ mặt, nhưng không nói thêm gì.

\ct

Sau đó, một người đàn ông khác - có lẽ là một trong các bác - lại chuyển sang hỏi về
chuyện học tập. Đó không phải là chủ đề mà tôi thấy khó chịu gì, nhưng năm nào bố mẹ
tôi cứ khen tôi học giỏi lắm, nên gần như mọi người toàn đặt nặng kỳ vọng vào tôi cũng
đủ khiến tôi đau đầu.

Bác ấy xởi lởi hỏi: ``{\color{red} Cháu đang học đại học, đúng không? Thế bao giờ ra trường? Có
  định học lên nữa không, hay sẽ ở nhà làm việc phụ bố mẹ?}''

Tôi đáp lại một cách khá nhẹ nhàng: ``{\color{red} Cháu cũng không chắc ạ. Chắc là cháu sẽ học
  khoảng hai, ba năm nữa. Còn cháu cũng không chắc là sẽ ở nhà phụ việc bố mẹ đâu ạ-}''

Tôi đang trả lời thì Pekora-chan cắt ngang lời tôi: ``{\color{red} Vì đến bản thân anh ấy còn chưa lo
  xong mà thi làm sao lo được cho gia đình đâu ạ, peko.}''

Tôi thở dài, chỉ về phía em ấy và hoàn thành câu nói: ``{\color{red} \dots như em ấy nói đấy ạ.}'' Ông bác
ấy nghe câu trả lời của chúng tôi, bắt đầu giở giọng xỉa xói: ``{\color{red} Ui dời, cháu kém thế! Thời
  ngày xưa ấy, chú đã bắt đầu kiếm tiền được rồi!}''

Câu trả lời của bác khiến tôi không thoải mái, nhưng trước khi tôi kịp trả lời, Subaru-chan
- với phong cách thẳng thắn thường thấy - đã chen vào: ``{\color{red} Bác ơi, thời của Bác xưa như
  Diễm rồi. Việc học tập của Kuon-senpai chắc chắn là tốt rồi. Bác không cần lo lắng đâu.
  Với lại, anh ấy còn làm việc ở một công ty lớn mà!}''

Người đàn ông kia hơi khựng lại, không ngờ bị nàng Vịt ``phản pháo'' nhanh đến thế.

\ct

Chưa dừng lại ở đó, một người họ hàng khác cũng hỏi han tình hình cuộc sống của tôi.
Nhưng lần này thậm chí còn lỗ mãng hơn, khi người đó lại bắt đầu so sánh tôi với con
của họ. ``{\color{red} Con tôi năm nay mới tốt nghiệp đại học, mà đã được nhận vào một công ty lớn
  với mức lương cao. Cháu thì sao, hở?}''

``\textit{Việc con cháu nhà chú thì liên quan gì đến tôi cơ chứ\dots}'' đó là điều mà tôi muốn thốt ra,
thậm chí còn muống ``đôi co'' với người họ hàng bất lịch sự đó, nhưng trước khi kịp nói
gì, Ayame-chan - với vẻ ngoài hăng như bò tót cùng tính cách sắc sảo - lên tiếng đầy mỉa
mai: ``{\color{red} Vậy sao ạ? Thế thì chúc mừng anh chị và cháu ạ. Nhưng em nghĩ mỗi người có con
  đường riêng. Anh chị đi đường anh chị, còn Kuon-senpai sẽ đi đường của Kuon-senpai,
  đúng không ạ?}''

``{\color{red} Với cả, chú có thấy bất lịch sự không khi nói phả thẳng mặt thế ạ?}'' Shion-chan lên tiếng
bênh vực tôi, dẫu rằng bản thân em ấy cũng là một kẻ thô lỗ, làm cho tôi cảm thấy hơi sai
sai.

Dù vậy, lời nói của hai em ấy như một cú đấm thẳng mặt, khiến người họ hàng kia chỉ
biết cười gượng và ngậm miệng.

\il

Bầu không khí có phần căng thẳng dần trở lại bình thường khi chúng tôi tiếp tục chuyến
đi. Tôi thầm cảm ơn các cô gái vì đã đứng ra đỡ lời cho tôi. Dù đôi lúc họ có hơi phiền
phức, nhưng những khoảnh khắc thế này khiến tôi nhận ra rằng, họ thực sự là một phần
quan trọng của gia đình \hll\ mà tôi gắn bó.

\dots Tôi thề, tôi sẽ KHÔNG BAO GIỜ về quê ngoại thêm một lần nào nữa.

\ct

Chuyến ghé thăm nhà thứ ba của buổi sáng hôm đó hóa ra lại kéo dài hơn dự kiến, chủ
yếu vì không khí ở đây thoải mái hơn hẳn. Cả gia đình chủ nhà lẫn các cô gái đều có
vẻ hòa hợp ngay từ đầu. Những câu chuyện rôm rả, tiếng cười đùa vang vọng khắp nơi, khiến tôi cảm thấy nhẹ nhõm hơn sau những câu hỏi khó chịu ở nhà trước. Mọi người ở
đây cũng không hỏi những câu hỏi khó chịu hơn, mà thay vào đó, họ thật sự quan tâm tới
tôi.

Ít ra ở đây, tôi có thể chỉ mặt được một vài người, nhưng tôi không thể nhớ rõ tên. Một
lần nữa, họ hàng không phải nằm trong vùng ưu tiên của tôi mà. Nhưng dù gì, ở đây, cũng
có những câu chuyện đáng mến và dễ thương hơn. Có thể kể như là\dots

\il

Trong lúc mọi người đang tụ tập ở phòng khách, một đứa bé tầm vài tháng tuổi được bế ra
để mọi người ngắm nghía. Đứa nhỏ bụ bẫm, má phúng phính đáng yêu khiến Sora-chan
và mẹ tôi ngay lập tức ``kích hoạt chế độ bà mẹ trẻ.''

Mẹ tôi, người vốn rất thích trẻ con, đưa tay ra đón bé từ tay bà chủ nhà. Nhưng Sora-chan,
với tốc độ thần tốc và nụ cười tươi rói, cũng nhanh chóng giành lấy.

\begin{dialogue}
  \speak{Kaien} Sora-chan, để bác bế nào!
  \speak{Sora} Cháu xin lỗi bác, nhưng cô để cháu muốn bế bé một chút thôi! Bé đáng yêu quá ạ!
\end{dialogue}

Mẹ tôi không chịu thua, cố gắng thương lượng với em ấy: ``Thế thì để cô bế trước, rồi lát
nữa cháu bế sau. Cô cũng yêu trẻ con lắm!''

Nàng Thần tượng mỉm cười, vẫn ôm chặt lấy đứa bé: ``Không được đâu, cháu đã bế trước
rồi mà!''

Cả hai nhìn nhau, tạo nên một không khí vừa căng thẳng vừa buồn cười, trong khi đứa bé
hoàn toàn không để ý gì, chỉ ngọ nguậy trong vòng tay Sora-chan.

Tôi nhìn cảnh tượng ba người chơi đùa với nhau như thế, lắc đầu thở dài: ``Họ đúng là
hết thuốc chữa mà. Có mỗi một đứa bé thôi mà sao tranh giành nhau ghê thế\dots''

Sở dĩ tôi nói vậy, là bởi tôi cực kỳ ghét em bé. Tưởng tượng đang yên đang lành, bỗng
dưng một đứa trẻ khóc toáng lên chỉ vì\dots\ gì? Đòi ăn, đòi thay tã, rồi còn đòi bế\dots\ Tôi
biết, tôi đã từng là một em bé, nên tôi hiểu được vì sao chúng lại phản ứng như vậy, nhưng
đó không phải là lý do chính tôi ghét chúng!

Tôi ghét chúng vì cứ mỗi lần chúng nó quậy, thậm chí còn phá hỏng cả những thứ mà
chúng ta cho là rất quý giá, thì cha mẹ của chúng lại chỉ buông một câu đầy hờ hững:
``Trẻ con có biết gì đâu!'' Đó là câu nói đã làm tôi cực, CỰC ức chế khi mà một cháu nhỏ
(mà tôi gọi là thằng) đã phá hỏng cái máy tính của tôi trong lúc làm việc. Trẻ con đúng
là không biết gì thật, nhưng nếu các vị không dạy được con thì để chúng tôi dạy hộ thay
các vị cho!

\dots Được rồi, tôi để cảm xúc của tôi lấn át quá rồi, ờm\dots\ Đến đâu rồi nhỉ?

Bố tôi đứng bên cạnh tôi, nhìn cảnh mẹ tôi và Sora-chan thay phiên nhau chơi với đứa trẻ
thì chỉ bật cười: ``Đây gọi là bản năng của người làm mẹ đấy!''

Tôi chỉ có thể thở dài trong ngao ngán, liếc nhìn Miko-chan đang đếm phong bao lì xì.
Tôi buột miệng: ``Ít nhất thì Miko-chan còn ngoan hơn bọn nhãi con kia nhiều\dots''

Dường như em ấy đã nghe tôi nói, nên nàng Vu nữ ngơ ngác nhìn tôi, thốt lên: ``Em thì
liên quan gì đến chuyện này chứ, ne!? Anh nói thế ám chỉ rằng Miko là em bé đúng không
ne!?''

\dots Thôi kệ.

\ct

Trong khi mẹ tôi và Sora-chan đang mải ``tranh giành'' quyền bế đứa bé, tôi cùng bố,
Choco-sensei và nàng Dị giới Aki ở trong bếp, chuẩn bị cho bữa trưa.

Nhà này có rất nhiều khoai tây, chất đầy trong góc bếp. Bố tôi quyết định làm món khoai
tây xào – vừa để thay đổi khẩu vị, vừa vì nguyên liệu sẵn có. Ông nhanh chóng giao nhiệm
vụ cho từng người: ``Kuon, con lấy chai mật mía trong tủ lạnh giúp bố. Choco-san, cháu
làm thịt xiên nướng nhé. Còn Aki-chan, cháu chuẩn bị nước sốt rưới cho món khoai tây.''

Choco-sensei nhận nhiệm vụ, nụ cười tươi tắn: ``Yên tâm đi ạ! Xiên thịt thì cháu làm
quen rồi!'' Em ấy lấy thịt lợn, ướp nhanh với gia vị rồi tỉ mỉ xiên từng miếng thịt vào que.
Tôi liếc nhìn, nhận ra cách em rắc hạt tiêu khá khéo, khiến tôi không nhịn được gật gù:
``Choco-san làm việc nhanh thật đấy.''

Trong khi đó, Aki-chan loay hoay với chỗ hành tỏi và nước tương để pha nước sốt. Tôi
hỏi em ấy: ``Có cần anh giúp một tay không?''
Nàng Dị giới cười nhẹ, mặc dù trông hơi lúng túng: ``Em nghĩ chắc là được. Nhưng mà\dots\ Anh có chắc là mình không cho thêm gì đặc biệt không?''

Tôi lắc đầu: ``Nhà người ta chứ có phải là nhà mình đâu. Làm nhạt nhạt chút, xong rồi
anh sẽ tự điều chỉnh sau.''

Bố tôi thì vừa đảo khoai trên chảo, vừa nói vọng lại: ``Kuon, lấy bát đũa ra bàn giúp bố.
À, chuẩn bị thêm ít bánh chưng nữa nhé. Lát bố sẽ làm món đặc biệt.''

Tôi nghe vậy, liền hiểu ngay ý mà ông ấy đang nhắc tới. Một món ăn đặc sản của tỉnh
Neihei - hoặc chính xác hơn - gia đình nhà ngoại chúng tôi, cũng là thứ ấn tượng với tôi:
``Món bánh chưng ăn với mật mía phải không? Con mong chờ mãi món đó đấy!''

Nàng Y tá nghe chúng tôi nói vậy, tỏ vẻ tò mò: ``Bánh chưng ăn với mật mía? Là sao thế,
Kuon-sama?''

Tôi cười, đáp: ``Là cách ăn bánh chưng hơi khác. Nó cũng dễ hiểu mà: bánh chưng, ăn
với mật làm từ mía. Cứ đợi mà xem, chắc chắn mọi người sẽ thích.''

Nàng Dị giới nghe vậy cũng háo hức: ``Nghe thú vị quá. Bánh chưng đã là món gì đó mới
mẻ rồi, món này có khi sẽ còn lạ lẫm hơn nữa!''

Bầu không khí trong bếp vừa nhộn nhịp vừa thoải mái. Tôi cảm thấy may mắn vì ít nhất
ở đây, không có ai hỏi những câu khó chịu hay so sánh tôi với con cái nhà khác nữa. Đôi lúc, một bữa ăn đơn giản cùng những người bạn như thế này cũng đủ để xoa dịu mọi áp
lực.

\ct

Sau chưa đến một giờ bốn người chúng tôi hì hục trong bếp, bữa trưa nhanh chóng được
dọn ra sau khi các món đã hoàn thiện. Tất cả mọi người tề tựu lại ở hai gian phòng liền
nhau. Như mọi năm, nhà này thường có khoảng hai mươi người cùng ngồi ăn, nhưng năm
nay với sự xuất hiện của cả đoàn idol, chúng tôi phải chia thành hai mâm lớn.

Và đoán xem tôi ngồi ở đâu? Tất nhiên là mâm có các cô gái thần tượng - theo một cách
bất đắc dĩ. Mâm này còn có mẹ tôi, em gái tôi, và một vài người thân khác. Đúng kiểu
``xung quanh tôi toàn là nữ, ey\footnote{Chế từ bài \textit{Hai triệu năm} của Đen Vâu (hát cùng với Biên.)}.''

Mẹ tôi cười trêu: ``Xung quanh toàn người đẹp thế này là sướng quá còn gì.'' Cô em gái
tôi thậm chí còn bông đùa: ``Chẳng trách mà các bác cứ toàn nói trêu rằng anh có bạn
gái\~{}''

Tôi nhăn mặt, tỏ vẻ chống chế: ``Con lại đang thấy áp lực ngang đây mẹ, chứ có sung
sướng gì đâu.'' Thành thật mà nói, việc ăn cơm với các cô gái thần tượng là một niềm
vinh dự mà ai cũng muốn có, nhưng đấy là ông thanh niên ất ơ nào đó ngoài kia, chứ
không phải tôi.

\dots Mà ngẫm lại thì, hình như tôi đã ăn với mười chín thần tượng \hll\ này từ bữa cơm
mùng Một hôm qua rồi nhỉ?

Ngay khi tôi buông ra một câu nói có phần hơi thô, Rushia-chan lập tức nhập vai, cúi gằm
mặt, giọng thút thít như sắp khóc: ``Anh không muốn ngồi với chúng em sao\dots?''

Tôi biết rõ em ấy đang giả vờ tỏ vẻ thương hại để làm tôi cảm thấy tồi tệ, nhưng ngay khi tôi định trả lời, Mel-chan quay sang, nhìn
tôi với ánh mắt như đang xuyên thấu tâm trí tôi: ``Ý anh không phải thế\dots\ đúng chứ?''

Tôi vội vã giơ tay xua xua, cười khổ: ``Không, không phải! Ý anh là\dots\ Được ngồi với
mấy đứa tận hai lần trong ngày là\dots\ hơi ngoài sức tưởng tượng thôi.'' Ngay khi tôi nói
vậy, cả năm cô gái đều giãn mặt, vui vẻ hơn hẳn. ``Đó, anh nói như vậy ngay từ đầu có
phải hơn không!'' Ru-chan nói, đấm mạnh vào lưng tôi, tỏ vẻ khoái chí.

Chậc\dots\ Tôi đúng là kẻ ngốc mà.

Sora-chan mỉm cười hiền hòa, vỗ nhẹ vào lưng tôi động viên: ``Anh không phải ngại.
Chúng em không ăn thịt anh đâu mà.'' Miko-chan ngồi cạnh cũng hào hứng phụ họa:
``Thôi nào, chuẩn bị ăn thôi ne\~{}!''

Chúng tôi cùng nhau mời cơm và bắt đầu bữa cơm trưa. Khi bữa ăn bắt đầu, mọi người
đều tò mò với bát nhỏ đặt trước mặt mỗi người. Trong đó chứa chất lỏng đen sánh, trông
có vẻ không giống bất kỳ món ăn truyền thống nào.

Noel-chan nhíu mày nhìn bát mật, quay sang hỏi tôi: ``Đây là cái gì vậy?'' Tôi đáp gọn:
``Mật mía đấy.''

Roboco-chan nghe vậy, ngơ ngác: ``Mật mã? Nhưng nó có ăn được đâu?'' Tôi bật cười: ``Anh nói `mật mía', không phải `mật mã'. Có cần anh khám lại thính giác cho em không đây, Ngố ơi?''

Nghe xong, nàng Người máy phồng má lên, đỏ mặt: ``Đừng có mà trêu vậy trước mặt mọi
người chứ!''

Haato-chan xen vào, nghiêng đầu nhìn chất lỏng sánh trong bát của mọi người: ``Nhìn nó
giống mật ong nhỉ?'' Mẹ tôi ngồi đối diện chúng tôi, giải thích: ``Gần giống như vậy, nhưng mà mật mía sánh
hơn.''

Subaru-chan gật gù: ``Chắc nó ngọt hơn mật ong nữa, đúng không?'' Tôi đưa mắt ra hiệu:
``Cứ thử đi, mỗi người một bát nhỏ. Anh dám chắc đây sẽ là một trải nghiệm khó quên
đấy.''

Ban đầu, tất cả mọi người nhìn chiếc bát nhỏ ấy với đầy rẫy sự nghi ngờ. Nhưng rồi,
Miko-chi can đảm (?), tò mò chấm đầu đũa vào bát mật và nếm thử. Vừa cho vào miệng,
mặt em ấy khẽ nhăn lại: ``Ngọt quá ne! Ngọt hơn cả kem bánh taiyaki nữa!''

Nhìn vẻ mặt không thoải mái của em ấy, tôi phì cười: ``Thế nên nó mới ăn kèm với bánh
chưng. Anh biết mấy đứa cũng bắt đầu ngán bánh chưng ăn mặn rồi, đúng không?''

Mel-chan và A-san thì bát mật, rồi nàng Dơi lên tiếng tỏ vẻ ngần ngại: ``Nhưng mà liệu
đồ mặn với đồ ngọt trộn với nhau có ổn không?''

Tôi mỉm cười, đáp: ``Nhiều thứ tưởng không liên quan với nhau, nhưng khi kết hợp lại
thành siêu phẩm đấy.'' Thậm chí, A-san cũng phải gật gù đồng tình: ``Ẩm thực nơi này có
nhiều điều bất ngờ lắm đấy. Nó cũng đáng thử.''

Nghe vậy, Noel-chan xung phong thử trước. Tôi cẩn thận nhắc: ``Nói thế chứ chấm ít
thôi, nó ngọt lắm đấy.''

Noel nhẹ nhàng chấm miếng bánh chưng vào bát mật rồi cắn thử. Khuôn mặt lập tức giãn
ra, đôi mắt sáng rỡ, em ấy ``ưm\~{}'' một tiếng dài: ``Ngon quá! Còn dễ ăn hơn cả bánh
chưng thường nữa!''

Thấy phản ứng tích cực, các cô gái khác lần lượt thử món ăn độc đáo này. Một số như
Subaru-chan, Sora-chan, và Miko-chi khen ngợi không ngớt, thậm chí ăn hết bát mật mía.
Ngược lại, vài người như Mel-chan, Ru-chan, và A-san vẫn trung thành với kiểu ăn bánh
chưng truyền thống.

Cũng phải thôi, vì như lời mẹ tôi nói: ``Mỗi người một khẩu vị, miễn là thấy ngon miệng
là được.'' Tôi không bắt ép họ đâu.

Bữa cơm trưa diễn ra trong không khí thân mật và vui vẻ. Tiếng nói cười xen lẫn những
lời khen dành cho các món ăn. Tôi, dù bị bao vây bởi ``hội chị em,'' cũng thấy nhẹ nhõm
khi không ai đề cập đến những câu hỏi khó chịu từ sáng nay nữa.

Vậy đấy, đôi khi những thứ tưởng chừng đơn giản, như bát mật mía với miếng bánh chưng,
lại đủ để gắn kết mọi người lại với nhau trong ngày Tết.

\ct

Sau bữa ăn trưa, chúng tôi lại tiếp tục hành trình chúc Tết. Lịch trình buổi chiều cũng
không khác gì buổi sáng: thêm hai, ba nhà nữa, chủ yếu là họ hàng xa mà tôi còn chẳng
nhớ nổi tên.

Mỗi lần đến nhà mới, mẹ tôi và bố tôi đều dẫn đầu đoàn, tươi cười chào hỏi mọi người,
giới thiệu từng cô gái trong nhóm idol đi cùng. Ông bà ngoại cũng nhiệt tình trò chuyện,
hỏi thăm tình hình năm qua của từng người. Còn em gái tôi thì đặc biệt nổi bật khi được
các bác và các chị em khen là ``xinh xắn và ngoan ngoãn.''

Nhìn những gương mặt xa lạ này, tôi tự hỏi làm sao mà mình có thể nhớ hết được chừng
ấy người. Dù đã cố gắng tỏ ra lịch sự, tôi vẫn không tránh khỏi cảm giác như mình chỉ
đang hoàn thành một nhiệm vụ bắt buộc: chào hỏi, nói vài câu xã giao, nhận lì xì, và\dots\ chấm hết.

``\textit{Chắc mình phải để ý nhiều hơn nữa,''} tôi nghĩ thầm, nhưng rồi lại tự bật cười chế giễu:
``\textit{Kể cả có để ý thì có nhớ nổi đâu. Cả năm có mỗi một dịp về quê gặp họ, còn thì giờ nào
  mà thân thiết chứ?}''

Sau mỗi màn giới thiệu và trò chuyện ban đầu, tôi lại tìm cách lẻn ra khỏi phòng khách.
Điện thoại trên tay, tôi giả vờ bảo mọi người rằng ``có việc gấp từ công ty'' để tránh ánh
mắt dò xét của họ hàng. Tâm trí tôi lúc này không buồn ặt vào câu chuyện rôm rả ở phòng
khách, mà ở những dòng tin nhắn trên màn hình hoặc lướt qua vài tin tức linh tinh.

Tuy nhiên, thành thật mà nói, trong số những người mà chúng tôi ghé thăm cũng có vài
nhân vật khá nổi bật. Chẳng hạn như\dots

\il

Một bác gái tầm ngoài năm mươi, luôn hỏi han bố mẹ tôi về công việc kinh doanh. Đôi
lúc, bác lại quay sang tôi, nói đùa: ``Kuon giỏi thế này chắc sắp thành ông chủ lớn rồi
nhỉ?'' Tôi chỉ cười trừ, không biết đáp với câu nói ấy như thế nào.

Một ông bác khác, chắc tầm sáu mươi, bảy mươi tuổi gì đó, rất thích kể chuyện về thời
xưa, làm cả đám trẻ nhỏ muốn ngồi nghe. Tôi thoáng nghe thấy Subaru-chan, Mio-chan
và Miko-chi cũng bị kéo vào ``cuộc trò chuyện lịch sử'' này, nhưng có vẻ họ khá thích thú.

Hay một cậu em họ, chừng nhỏ tuổi hơn tôi bốn đến năm năm, tỏ ra ngưỡng mộ khi thấy
tôi có đoàn thần tượng hộ tống. Em ấy cứ liên tục hỏi han họ các câu hỏi và được các cô
gái vui vẻ trả lời (dám chắc là họ đã quen với điều đó). Thậm chí còn xin chữ ký nữa.
Aki-chan và Choco-san vui vẻ ký tặng cậu bé, nhưng tôi chỉ đứng ngoài nhìn, không xen
vào.

\il

\ct

Cứ mỗi lần chúng tôi sang một nhà như vậy, các cô gái và cả gia đình (trừ tôi) đều tay bắt
mặt mừng, nói chuyện và chúc Tết cho nhau với vẻ đầy ấm cúng, như thể họ là một đại
gia đình. Riêng tôi, mỗi lần lẻn ra ngoài, tôi lại chọn một góc yên tĩnh để tránh phải tiếp
tục các cuộc trò chuyện không mấy thú vị. Có lúc, tôi thậm chí còn đi ra hẳn sân trước,
đứng ngắm cây cối hay nghịch chiếc điện thoại.

Tất nhiên bố mẹ tôi cũng không đồng lòng gì với phản ứng thờ ơ của tôi với họ. Mẹ tôi
thậm chí còn nhắc nhở: ``{\color{red} Con làm thế là không được đâu. Người ta sẽ nghĩ con không tôn
  trọng họ đấy.}''

Đáp lại, tôi chi cười xòa và nói: ``{\color{red} Thì ít nhất con cũng đã chúc Tết họ và nói chuyện xã
  giao rồi mà mẹ.}''

Mẹ tôi thở dài, nhưng cũng không nói thêm gì. Có lẽ bà hiểu rằng tôi thật sự không cảm
thấy gần gũi với những người này, dù họ về danh nghĩa là họ hàng.

Hành trình chúc Tết buổi chiều kết thúc khi trời đã xế bóng. Chúng tôi chuẩn bị đến ngôi
nhà chính ở quê để cùng gia đình lớn nấu nướng bữa tối - một bữa ăn đặc biệt đậm chất
quê hương, mà mẹ tôi đã nhắc đến từ sáng.

Dù cả ngày có phần mệt mỏi với những cuộc gặp gỡ không mấy ý nghĩa với tôi, tôi vẫn
cảm thấy nhẹ nhõm vì mọi việc đang diễn ra suôn sẻ. Chỉ cần qua được hết ngày hôm
nay, chắc tôi sẽ thoát khỏi ``lịch trình họ hàng'' này.

Bố tôi thông báo rằng điểm dừng chân cuối cùng sẽ là nhà của một người họ hàng gần,
nơi gia đình chúng tôi thường tổ chức bữa tối sum họp. Nghe đến đây, tôi cũng thấy yên
tâm phần nào, ít ra lần này sẽ không phải đối mặt với những khuôn mặt hoàn toàn xa lạ.
Cũng đội ơn giời, họ cũng là người hiểu rõ tính cách tôi nhất, nên họ cũng sẽ không hỏi
tôi những câu hỏi khó đỡ đến mức ức chế như vừa rồi.

Dẫu vậy, có một điều buồn cười là, dù đã tới nhà này vài lần trước đây, tôi vẫn không nhớ
nổi chủ nhân của nó tên Hiki hay Niki. Hai người này là anh em họ, nhưng không chỉ
giống nhau về diện mạo mà còn cả cách nói chuyện, khiến tôi chẳng thể phân biệt nổi. Đã
thế, họ lại cứ thích mặc áo sơ mi giống nhau, như thể muốn thử thách trí nhớ của tôi vậy.

Ngược lại, Ryuuhou lại chẳng bao giờ nhầm lẫn. ``{\color{red} Ông Hiki-san cao hơn, cằm vuông hơn,
  còn ông Niki-san thì có râu quai nón nhẹ,}'' em ấy giải thích, dù tôi không chắc điều đó có
giúp ích gì không. Nhìn kiểu gì thì cả hai vẫn như hai giọt nước trong mắt tôi.

Những lúc như thế này, tôi luôn tự nhủ: ``\textit{Kệ đi, chỉ cần cười xã giao và gọi họ là `ông' đi
  là sẽ ổn thôi.}''

Khi xe dừng lại trước một cung đường nhỏ quen thuộc, tất cả chúng tôi bước xuống, hít
một hơi thật dài. Không khí lành lạnh của buổi chiều muộn khiến tôi thấy dễ chịu hơn sau
cả ngày lê thê.

Con đường được xây với bức tường gạch dựng hai bên, dẫn tới một căn nhà rộng với hai
gian nhà và một khoản sân lát đá trơn rất rộng. Tiến tới gần hơn, tôi đã nghe thấy tiếng
trò chuyện và tiếng trẻ con cười đùa vọng ra từ trong nhà.

``{\color{red} Đến rồi à? Lâu lắm mới thấy đấy, hai bố con!}'' - một giọng nói vang lên, nhưng tôi
không thể chắc đây là của ông Hiki-san hay là của ông Niki-san.

Tôi đáp lại một cách lễ phép như thường lệ: ``{\color{red} Dạ, lâu lắm rồi ông ạ.}'' Các cô gái đi theo
tôi cũng chào mọi người một cách lễ phép.
Đi vào trong, tôi thấy không gian quen thuộc khiến mình cảm thấy thư thái hơn. Mọi người trong nhà đang tất bật chuẩn bị cho bữa tối, nhưng không khí ở đây thoải mái hơn
những nhà trước. Thay vì phải căng thẳng với những câu hỏi ``hóc búa,'' tôi có thể thư
giãn và ngồi lặng lẽ một góc nếu muốn.

Dù vậy, tôi vẫn để ý rằng bố mẹ và em gái tôi hòa nhập rất nhanh. Mẹ tôi vui vẻ trò chuyện
với các cô trong bếp, còn bố tôi giúp dọn bàn ăn. Ryuuhou, tất nhiên, cũng không phải
ngoại lệ, khi cô nhóc cũng hớn hở chào mọi người và nói chuyện với các bác, ông bà tại
đây.

Còn tôi thì sao? Tất nhiên, nếu tôi lại cắm đầu điện thoại thì sẽ cực kỳ chán, và chẳng
khác gì Tết mấy năm trước cả. Vì thế, tôi cùng các thành viên \hll\ cùng nhau chơi
những trò chơi dân gian mà tôi biết, hoặc đã lâu rồi không có dịp để chơi. Điều này không
chỉ giúp mọi người thư giãn mà còn khiến không khí trở nên sôi động hơn.

\ct

``Được rồi, chúng ta sẽ bắt đầu với trò nhảy lò cò đi,'' tôi đề xuất, khiến mọi người tò mò.

``Nhảy lò cò? Là trò gì vậy, Kuon-senpai?'' Sora-chan nghiêng đầu hỏi, vẻ bối rối hiện rõ
trên khuôn mặt.

Tôi giải thích: ``À, là trò mà tụi trẻ con hay chơi. Vẽ mấy ô vuông trên đất, rồi phải nhảy
qua từng ô bằng một chân.''

Tôi đưa hai cục phấn cho Subaru-chan và Noel-chan, rồi cùng nhau vẽ một dãy ô vuông
dài trên sân. Các ô vuông được đánh số từ 1 đến 10, nhưng thay vì sắp xếp thành một
đường thẳng như bình thường, chúng tôi xếp thành một con đường ngoằn ngoèo, giống
như một bàn cờ tỷ phú thu nhỏ. Mỗi vài ô còn được nối với nhau thành hình zig-zag
(dích-dắc) để tăng độ khó.

``Trông thế này sao? Nhưng trông nó có khác gì cờ tỷ phú đâu,'' Haato-chan nhận xét,
khiến một số người cũng phải gật gù đồng tình.

``Có thể nói như vậy. Anh gọi nói là\dots\ nhảy lò cò phiên bản \textit{hologra},'' tôi nói, vứt cục
phấn ra ngoài, cười hóm hỉnh. Một số người là ``chuyên gia phá hoại'' như Ayame-chan
và Ru-chan thì tỏ vẻ hào hứng với trò chơi quen thuộc (với tôi) nhưng lạ lẫm này.

\begin{dialogue}
  \speak{Ayame} Chơi liền! Chơi liền, yo!
  \speak{Rushia} Trông nó có vẻ hay đấy.
  \speak{Matsuri} Kuon-senpai hôm nay là bày trò có vẻ hay ho lắm ha?
\end{dialogue}

Nhưng một số người thì cũng tỏ ra ngờ vực.

\begin{dialogue}
  \speak{Mio} Trông nó có vẻ\dots\ hơi nguy hiểm.
  \speak{Choco} Cô chưa từng chơi qua trò này bao giờ.
  \speak{Aqua} Huống chi là phiên bản \textit{hologra} này\dots
\end{dialogue}

Tất cả mọi người đều đang bàn tán xôn xao. Aki-chan giơ tay, hỏi: ``Vậy, luật chơi cụ thể
là gì vậy, Kuon-senpai?''

``Luật đơn giản thôi,'' tôi giải thích. ``Mỗi người sẽ lần lượt nhảy bằng một chân qua từng
ô, không được đụng vạch kẻ. Nếu đụng vạch, phải quay lại điểm xuất phát. Nhưng ở các
ô nối với nhau, người chơi phải nhảy liền hai chân cùng lúc. Ai đi hết con đường trước sẽ
thắng.''

``Tưởng dễ mà khó phết đấy,'' Mio-chan gật gù. Korone-chan giơ hai tay lên nắm lại, thốt
lên đầy hứng thú: ``Nhưng mà trông cũng hay lắm á!''

``Đúng vậy,'' tôi nháy mắt. ``Nhưng nếu ai làm rơi đồ vật mà anh ném vào ô số nào đó,
phải chịu một hình phạt đặc biệt.''
``Hình phạt gì vậy?'' Roboco-PON hỏi, tỏ vẻ thắc mắc trong sự hiếu kỳ. Tôi mỉm cười
đầy bí ẩn, đáp gọn ghẽ: ``À, lát nữa chơi thì biết.''

Khi trò chơi bắt đầu, đúng như tôi dự đoán, có hàng loạt pha ngớ ngẩn đến mức khiến cả
sân náo nhiệt bởi những tiếng cười không dứt. Ai cũng tranh nhau nhảy, nhưng việc giữ
thăng bằng trên một chân trong khi ô thì nhỏ, đường thì cong queo, chẳng dễ chút nào.
Và tất nhiên, những hình phạt kỳ lạ do tôi nghĩ ra càng làm trò chơi trở nên ``đậm chất
\textit{hologra}'' hơn bao giờ hết.

Tôi có thể kể ra một vài tình huống dở khóc dở cười mà các thành viên \hll\ đã trải
qua. Chẳng hạn\dots

\il

Khi đến lượt của Mio-chan, em ấy trông có vẻ khá tự tin. Từng bước nhảy của em ấy đều
được thực hiện thoăn thoắt và khéo léo, khiến chúng tôi ngạc nhiên.

``Ba, bốn, năm\dots\ Tám, chín, mười\dots\ Ờm, đoạn này nhảy sao đây?'' Sói đen-chan dừng lại
ở ô được đánh dấu sao.

``Đáp bằng hai chân,'' tôi nói. Mio đảo mắt, rồi làm như những gì tôi nói, và tiếp tục lò cò
ở các ô tiếp theo: ``Mười một, mười hai\dots\ ể\dots?''

Mio-chan lại dừng lần nữa. Ru-chan đi theo em ấy ở rìa ngoài, nghiêng đầu hỏi: ``Sao
vậy?''

Sói đen-chan chỉ về hai ô ở phía trước, thảng thốt: ``Ô-Ông thần nào vẽ hai ô cách xa thế
này!?'' Trước mặt của nàng Sói đen là hai ô số tám, nhưng khoảng cách giữa chúng xa
đến mức phải dạng chân hết cỡ mới đáp được. Thế nhưng, điều này lại đặt em ấy vào tình
huống khó xử: làm vậy thì chẳng khác nào ``phơi hàng'' cho mọi người xem!

Dám chắc đây là trò đùa của Miko-chan. Tôi nhìn thấy nàng Vu nữ đang rúc rích cười, rõ
ràng là thủ phạm.

``Cứ nhảy đại thôi, Mio-chan!'' Roboco từ đằng sau, hối thúc em ấy. Đồng thời những
người thích vui đùa như Koro-chan, Fubu-chan và Ayame-chan thì cứ luôn miệng cổ vũ
em ấy. Không còn cách nào khác, Mio-chan tặc lưỡi, thốt lên: ``Mấy người thật là\dots\ Đành
phải làm vậy! \textbf{Hây!}''

Nàng Sói đen cố nhảy, dạng hai chân ra để đáp đúng vào ô số tám. Nhưng tình thế trở nên hài hước khi em ấy đáp xuống với phần hạ bộ gần chạm đất, khiến em ấy la toáng lên:
``\textbf{Rồi như nào nữa!? Ai đó giúp em với!}''

Miko-chi chạy về phía trước mặt, giơ ra chiếc máy ảnh mà em ấy đã chuẩn bị từ trước,
tủm tỉm cười: ``Rồi, tạo dáng hay lắm ne! Cho chị xin vài kiểu!''

Thấy nàng Vu nữ chụp ảnh một cách đầy nhiệt huyết, khi cô thậm chí còn ``lăn xả'', quyết
``hy sinh vì nghệ thuật,'' tất cả chúng tôi không thể nào nhịn cười, còn bản thân Mio-chan
thì đỏ mặt, thốt lên tới mức sắp khóc tới nơi: ``\textbf{Mấy người làm ơn cứu tôi đi!}''

Cả đám cười ngất, còn tôi thì vội vàng kéo em ấy dậy trước khi mọi chuyện vượt quá mức.

\ct

Tới lượt Okayu-chan và Korone-chan - cặp đôi quậy nhất trong cả công ty. Nàng Mèo là
người thực hiện bước nhảy, trong khi bạn của mình đi theo giám sát.

``Mười bốn, mười lăm\dots'' Okayu-chan đếm, theo sau là Koro-chan: ``Mười sáu, mười bảy,
mười tám\dots''

Bất chợt, nàng Khuyển thốt lên: ``Gượm đã!'' đầy đột ngột, khiến bạn thân của em ấy suýt
ngã.

Nàng Mèo hoàn hồn, quay mặt về người bạn của mình hỏi: ``Sao thế, Koro-san?''

Nàng Chó nâu chỉ về dòng chữ trên ô vuông mà bạn của mình đang đứng, trên đó viết
dòng chữ: ``Trở về vạch xuất phát,'' do chính tôi viết.

Okayu-chan ngạc nhiên: ``Ể\~{}? Nhưng mà tôi sắp về đích rồi mà!'' Korone-chan cười
gian: ``Sao bà không đi đường kia? Ogayu chơi lại nha\~{}''

Tôi đứng từ xa, cười khì: ``Mấy đứa làm như thể đang chơi đổ xúc xắc vậy.'' Okayu-chan
nheo mắt nhìn ô số với vẻ bất mãn, nhưng cuối cùng vẫn phải quay trở lại điểm xuất phát,
trong tiếng cười vui vẻ của mọi người.

\ct

Aqua-chan là người tiếp theo, với Fubu-chan là người đếm ô số.

``Hai mươi hai, hai mươi ba, hai mươi tư\dots'' nàng Hầu gái vừa đếm vừa nhẩm số ô được
ghi trên đó.

Bất chợt, Cáo trắng thốt lên chen ngang: ``Rồi! Em có muốn mua nhà không?'' khiến
chúng toi đổ mồ hôi hột.

``Ể? Vậy cho em một căn\dots'' em ấy nói, bắt đầu lục lọi trong túi tìm tiền. Rushia-chan
thốt lên bất bình: ``Chúng ta có chơi cờ tỷ phú đâu, Aqua-senpai!! Đừng có dễ bị dụ khị
chứ!''

Fubu-chan nháy mắt đầy tinh quái, còn Akutan thì bối rối đến nỗi gật đầu đồng ý mua nhà
ngay mà không cần nghĩ. Cả đám phá lên cười, khiến Aqua đỏ bừng mặt, lí nhí: ``E-Em
chỉ đang thử Fubuki thôi mà!''

\il

Các bạn còn nhớ luật chơi mà tôi đưa ra chứ? Hình phạt dành cho người làm rơi đồ vật
mà tôi đưa cho ở ô số nào đó? Chúng ta có ba người bị chịu phạt rồi đây.

\il

Pekora-chan cẩn thận cầm viên sỏi mà tôi đưa, và chỉ định em ấy nhắm vào ô số bốn trong
lúc nhảy lò cò.

``Khó thế thì làm sao em làm được, peko!'' em ấy có thốt lên như vậy. Nhưng mà, đây là
trò chơi mà, đúng chứ? Thắng thua quan trọng gì.

Không còn cách nào khác, em ấy ném đi trong lúc cố gắng giữ thăng bằng. Viên sỏi bay
thẳng, nhưng rồi trượt khỏi mục tiêu, rơi vào ô số ba.

``Ê, sai rồi nha, rơi nhầm ô rồi kìa!'' Shion-chan thốt lên. Noel-chan cười: ``Thế thì phải
chịu hình phạt thôi, đúng không nè?''

``Mấy người chơi gian! Vừa nhảy lò cò vừa ném sỏi, ai mà làm nổi chứ, peko!?'' Thỏ-chan
tiếp tục ca thán.

Tôi gật đầu tỏ vẻ hứng thú, mặc kệ những lời trách móc từ em ấy. Vì tôi là quản trò của
trò chơi này, tôi có quyền quyết định đưa ra hình phạt. ``Pekora-chan, hình phạt của em
là\dots\ hóa thân thành thỏ trong vòng năm phút!''

Peko-chan ngỡ ngàng: ``Ể!? Nhưng em vốn dĩ là thỏ rồi mà, peko!'' Aki-chan bổ sung:
``Không phải chỉ làm thỏ bình thường đâu. Em phải vừa nhảy lò cò vừa kêu `peko peko'
mỗi bước đấy.''

Mọi người cười nghiêng ngả khi nàng Thỏ bắt đầu nhảy, vừa kêu ``peko peko'' với vẻ mặt
méo mó. Có lẽ em ấy đang tự hỏi tại sao lại đồng ý chơi trò này.

\ct

Tới lượt Matsuri-chan thực hiện cú ném, và tôi chỉ định em ấy ném ô số bảy. Cú ném
được thực hiện, nhưng viên sỏi lại rơi ngay vào\dots\ mép ô số sáu. Chưa vào được hẳn ô
thứ sáu nữa.

``Sai rồi, phải là ô số bảy cơ\~{}'' nàng Dị giới thốt lên, ``bu bu'' một tiếng. Matsuri-chan tặc
lưỡi: ``Chết thật đấy\dots\ Vậy tôi sẽ phải thực hiện hình phạt gì đây?''

Lần này tôi để Shion-chan đưa ra hình phạt: ``Matsuri-senpai sẽ phải thực hiện `Thử thách
hay nói thật'!''

Tất cả chúng tôi vỗ tay hưởng ứng, trong khi bản thân nàng Lễ hội ngơ ngác: ``Ể? L-Là
sao?''

``Luật đơn giản lắm,'' tôi tiếp lời giải thích. ``Em sẽ phải nói `thử thách' hay `sự thật', sau
đó bọn anh sẽ cho đề bài. Nhiệm vụ của em là phải thực hiện đúng yêu cầu đó.''

``Nghe có vẻ dễ nhỉ!'' Matsuri-chan mắt sáng lên, tỏ vẻ mừng rỡ. ``Em sẽ chọn\dots\ Sự thật!''

Noel-chan thấy vậy cười gian: ``Chị nói thế rồi là không được chọn lại đâu đấy! Được
rồi, giờ chị hãy kể lại điều xấu hổ nhất mà chị đã làm!''

Vẻ mặt vui tươi của Matsuri-chan đột ngột trở nên méo xệch vì ngỡ ngàng với yêu cầu kỳ
quặc này. ``C-Cái gì chứ?'' em ấy ngập ngừng, ``Chị có làm điều gì đáng xấu hổ đâu!''

``Thế à?'' tôi cười nhoẻn. ``Thế cái hồi mà em loay hoay tìm chỗ để `xả lũ' hồi mấy tháng
trước-\footnote{Xem {\flaregothic ÉPISODE 23}, quyển số Hai.}''

Lập tức, Matsu-chan chạy tới, bịt miệng tôi lại, đồng thời mặt của em ấy cũng đỏ hỏn lên:
``\textbf{Đừng có kể lại khoảnh khắc ấy!!}''

Cả đám cười rộ lên, tò mò với câu chuyện mà tôi định kể. Trong khi đó, Matsuri-chan
đỏ bừng mặt, lúng túng biện minh cho sự kiện ngày hôm đó, đồng thời lườm về phía tôi,
người vẫn đang vừa cười vừa giả vờ huýt sáo làm ngơ.

\ct

Cuối cùng là Choco-sensei. Tôi ra hiệu phải cho em ấy ném viên sỏi vào ô số mười, nhưng
cuối cùng trật lất sang ô số mười một.

Chưa cần để ai lên tiếng gì, nàng Y tá đã thở dài, sẵn sàng nhận hình phạt từ chúng tôi.
Aki-chan đặt ra hình phạt: ``Hình phạt đây: Choco-sensei phải đoán xem một người ở đây
đang nghĩ gì!''

Thấy thế, Sensei mỉm cười một cách đầy tự tin: ``Thật sao? À, cái này thì dễ mà. Cô có
thể hiểu được mấy em đang nghĩ gì ngay lúc này đấy.''

``Vậy cô thử đoán xem Kuon-senpai đang nghĩ gì đi,'' Haato nói. Tôi ngạc nhiên: ``Anh
á? Sao ở đây có bao nhiêu đứa mà lại đi chọn anh?''

``Vì anh là người duy nhất mà bọn em vẫn không thể đọc vị được đó,'' Ayame-chan đáp
lời, tay vẫn đang ghi lại mọi khoảnh khắc. Em ấy nói cũng phải, hầu như trong tất cả
những người mà tôi đã từng gặp, không ai đọc nổi xem tôi đang nghĩ đến cái gì, ngoại trừ
bố mẹ tôi.

Trông Sensei mỉm cười đầy tự tin như vậy, nhưng tôi biết chắc đây sẽ là một câu đố khó
nhằn đây. Choco-san nhìn vào biểu cảm dường như không cảm xúc của tôi, rồi mạnh dạn
đoán: ``Kuon-sama đang nghĩ rằng\dots\ Cô là người đẹp nhất ở đây, đúng không?''

Không một tiếng trả lời. Tất cả mọi người nín thở chờ tín hiệu từ tôi. Tôi nhắm mắt lại,
tỏ vẻ nghiêm trọng hóa vấn đề.

``Có vẻ như cô ấy đoán đúng rồi\dots'' Mio-chan nói, tỏ vẻ căng thẳng. Fubu-chan lại không
nghĩ vậy: ``Tôi không nghĩ thế đâu, Mio. Ở công ty này ngoài A-chan và YAGOO ra còn
ai đoán được anh ấy chứ.''

Tôi khẽ gật đầu: ``Cũng\dots\ gần đúng. Em thay bằng cụm `khiêu gợi' thì sẽ chính xác hơn
đấy.'' Câu trả lời của tôi khiến mọi người nhìn về phía Sensei, khẽ ngượng ngùng trước bộ cánh mà em ấy đang diện\footnote{Xem lại hình g, trang \ref{setthree}}. Mọi người ngẩn ngơ nhìn nhau, rồi khẽ bật cười, còn Choco-sensei thì cười khì, tỏ vẻ
hãnh diện với bộ mà em ấy đang mặc.

\il

Tiếp đến là trò tạt lon, một trò chơi mà tôi nghĩ hầu như ai từng trải qua tuổi thơ ở quê
cũng đều biết. Trò này không cần dụng cụ cầu kỳ, chỉ cần một chiếc lon, vài đôi dép hay
giày, và một khoảng không gian đủ rộng là đủ. Đối với các cô gái, họ đều tròn mắt như
thể tôi vừa giới thiệu một thứ gì đó từ hành tinh khác. Cũng phải thôi, quê tôi mới chơi
trò này, còn họ thì không.

``Tạt lon? Là cái gì thế?'' Aki-chan ngơ ngác hỏi. Noel tiếp lời, tỏ vẻ băn khoăn: ``Có
phải giống như ném đồ vật vào mục tiêu không?''

Shion-chan cũng không giấu nổi sự tò mò: ``Nếu thế thì dễ quá\dots\ nhưng nghe kỳ kỳ thế nào
nhỉ.''

Tôi bật cười trước sự ngơ ngác của mọi người. Cũng đúng thôi, trò này có lẽ chỉ phổ biến
ở quê, còn với các cô gái từ thành phố lớn hay đến từ những thế giới tưởng tượng như
Aki-chan, hẳn đây là điều mới mẻ.

``Ừ, gần đúng đấy, nhưng mà thêm chút hành động chạy trốn nữa. Để anh giải thích luật
nhé.''

Sau khi nghe tôi hướng dẫn qua luật chơi, mọi người bắt đầu tỏ vẻ hào hứng, đặc biệt là
những người có vẻ ưa thích vận động, như Subaru-chan, hay những người thích táy máy
tay chân như Ayame-chan hay cặp đôi Oka-Koro.

Tôi xung phong làm người trông lon đầu tiên, một phần vì muốn các cô gái được thử sức
với vai trò thú vị hơn, phần khác\dots\ là để tôi có thể mặc sức trêu chọc họ.

Nhưng vấn đề lớn nhất lại là đôi dép. Ai cũng nhanh chóng tháo giày hoặc dép của mình
ra để làm ``vũ khí,'' còn tôi thì phải đi tìm đôi nào thích hợp để chạy. Sau cùng, tôi mượn
tạm đôi dép của mấy bác ở nhà, vừa thoải mái vừa bền, và cũng dễ để chạy trên sân gạch
trơn.

``Kuon-senpai nhìn hợp với dép này ghê\~{}'' Shion-chan giở giọng trêu chọc. Pekora-chan
cũng hùa vào: ``Nhìn cứ như ông chú nào đó đấy nhỉ, peko!''

Tôi lườm nhẹ các em, nhưng đành nhún vai. ``Ông chú thì ông chú, nhưng lát nữa bắt
được đứa nào thì đừng khóc nhé.''

Mọi người bật cười, tiếng giày dép vang lên trên nền sân khi họ chuẩn bị sẵn sàng cho
cuộc chơi. Tôi xếp chiếc lon ngay ngắn giữa sân, đôi mắt liếc qua từng người như đang
dò xét đối thủ.

``\textbf{Nào, bắt đầu thôi!}''


Lượt đầu tiên diễn ra với không khí vừa hồi hộp vừa hào hứng. Tất cả mọi người đồng
loạt ném giày hoặc dép của mình về phía chiếc lon. Tiếng giày dép bay vèo vèo xen lẫn
tiếng cười khúc khích, nhưng trong số đó, chỉ có chiếc giày của Mel-chan là chạm được
vào chiếc lon, làm nó lăn lăn về phía trước một đoạn không quá xa.

Nhanh như chớp, tôi lao tới, nhặt lấy chiếc lon trước khi bất kỳ ai kịp phản ứng, rồi đặt
nó trở lại vị trí ban đầu một cách thản nhiên. Nàng Dơi vừa mới chạy lên được vài bước,
liền đứng khựng lại, đôi mắt mở to ngạc nhiên, rõ ràng không ngờ tôi có thể phản ứng
nhanh đến vậy: ``Ể!? Nhanh thế!''

Tôi phì cười, chạy lùi về phía sau: ``Ném mạnh lên nữa, lon đi xa mới tính chứ. Còn bây
giờ thì coi như em mất cái dép rồi.''

Mel-chan nhíu mày, vẻ mặt ngơ ngác. ``Mất? Ý anh là sao?''

``Tức là cái dép đó không được dùng nữa ở lượt này. Em phải chờ đến khi có ai đó ném
trúng lon rồi mới được lấy lại. Giờ thì lùi ra sau đi, nhường chỗ cho người khác.''

Em ấy lùi lại một cách miễn cưỡng, vừa đi vừa lẩm bẩm điều gì đó nghe không rõ, có lẽ
là đang thầm trách mình sao không ném mạnh hơn. Những người khác thì bắt đầu cười
rộ lên, cảm thấy trò chơi ngày càng thú vị hơn.

Tôi đứng thẳng dậy, cầm chiếc lon trong tay một cách tự tin, mắt đảo qua từng đối thủ.
``Tiếp theo, ai muốn thử sức nào? Đừng để anh bắt được ha!''

\ct

Lượt thứ hai bắt đầu với không khí sôi động hơn. Mọi người tiếp tục ném giày dép của
mình vào chiếc lon, ai nấy đều cố gắng hết sức để chiếc lon bay xa hơn. Lần này, chiếc
dép của Miko-chan đã trúng đích một cách hoàn hảo, khiến chiếc lon lăn xa thêm vài mét.

Nhưng trước khi tôi kịp phản ứng, một sự cố bất ngờ xảy ra: chiếc dép của Sora-chan bay
lệch hướng và đập thẳng vào mặt tôi! Tôi loạng choạng, mất thăng bằng và ngã nhào ra
phía sau.

\begin{dialogue}
  \speak{Kuon} Ái da!
  \speak{Sora} Ấy chết! Xin lỗi anh, Kuon-senpai!
\end{dialogue}

Tiếng cười rộ lên khắp nơi trước tình huống hy hữu này. Nhưng Miko-chi lại không hề
chú ý, mà thay vào đó, em ấy hét lên đầy hào hứng: ``Đổ rồi, ne!''

Nhưng rồi, em ấy đứng đó, ngơ ngác nhìn mọi người: ``\dots Giờ thì phải làm gì nữa nhỉ,
ne?'' Shion-chan hối thúc: ``Chạy lên lấy dép đi, đồ ngốc! Anh ấy bắt chị đi bây giờ!''

``Nhưng mà dép chị ở trên kia, ne!'' nàng Vu nữ chỉ về chiếc dép của mình đang ở ngay
cạnh chiếc lon. Aqua-chan thúc giục: ``Nhảy lò cò như ban nãy!''

Miko-chi nhảy lò cò đầy vội vã để nhặt chiếc dép tông của mình, vừa nhảy vừa nhìn ngó
xung quanh như sợ bị bắt (và sợ tất mình bị bẩn nữa). Tôi nhanh chóng bật dậy, nhặt chiếc
lon và đặt nó trở lại vị trí cũ. Thấy tôi đã kịp phản ứng, Miko-chan định quay người chạy
ngược lại. Nhưng quá muộn, tôi đã vươn tay chạm vào vai em ấy.

``Bắt được rồi! Miko-chan trông lon đi!'' tôi thốt lên ngay khi tóm được vai của em ấy.
Nàng Vu nữ ngạc nhiên, không hiểu em ấy đã bị bắt bằng cách nào: ``Ể? なんで、にぇ？
(Tại sao chứ, ne?) Bằng cách nào-''

Ru-chan vỗ vào trán, bất lực thở dài: ``Chị chậm chân thì chị dính chưởng thôi-nanodesu.''

Miko-chan nhăn nhó, khuôn mặt lộ rõ vẻ ấm ức. Tôi cười nhạt, chọc thêm một câu: ``Anh
nghĩ em nên hy sinh cái tất của mình đi. Còn không thì cứ cởi ra cho đỡ bẩn.''

Nghe vậy, em ấy triệu ra chiếc gậy và liên tiếp đập vào người tôi, mắng: ``Tại sao anh
không nói sớm chứ ne!?''

Tiếng cười giòn tan lại vang lên, và lần này, mọi người đều đồng tình rằng Miko-chi trông
lon sẽ là một cảnh tượng thú vị nhất trong ngày!

Lượt thứ ba bắt đầu với không khí háo hức và một chút căng thẳng, đặc biệt là từ Mikochan,
người vừa bị bắt làm ``người trông lon'' đầy oan ức. Em ấy vẫn còn nhăn nhó, mặt
đỏ bừng vì ngượng, nhưng vẫn cố tỏ ra nghiêm nghị để thực hiện vai trò mới của mình.

``\textbf{Ném, ne!}''

Theo lệnh của Miko-chi, tất cả mọi người đồng loạt ném giày dép của mình về phía chiếc
lon. Tiếng rầm rầm vang lên khi các vật dụng bay tới tấp. Nhưng bất ngờ, giữa cơn mưa
giày dép, chiếc chùy của Noel-chan lại lao vút lên trời, vượt qua chiếc lon, và đáp xuống
rất gần chỗ của đối phương, khiến em ấy giật nảy mình.

``\textbf{Em định giết chị sao, Noel!?}'' Miko-chan thảng thốt hét lên, hú vía sau cú ném trên.
Nàng Hiệp sĩ cười trừ: ``Xin lỗi\~{}!''

Miko-chan vội vàng lùi lại, nét mặt vừa bực bội vừa sợ hãi. Tôi đứng từ xa mà không nhịn
được cười. Quay sang Noel, tôi thắc mắc: ``Sao không dùng giày dép của em để ném?''

``Em không cởi ra được\dots\ Dép geta này chặt quá\dots'' nàng Hiệp sĩ giơ chân lên, cố kéo
chiếc dép của mình ra nhưng không thành. Cả nhóm phá lên cười khi nhìn thấy cảnh
tượng đó.

``Thế thì lần sau đừng mang vũ khí hạng nặng tham gia trò chơi dân gian nữa! Em có thể
lấy mấy cái dép của các bác được mà!'' tôi thốt lên, khẽ bật cười với tai nạn khó đỡ này.

Trong khi đó, nàng Vu nữ phồng má, giậm chân tỏ vẻ dỗi hờn: ``Thật là, mọi người không
ai nghiêm túc cả! Lon thì ở đây, không phải trên trời đâu, đồ ngốc\~{}!''

Cả nhóm lại một lần nữa tràn ngập tiếng cười, còn Miko-chan thì thở dài. Dù trách móc
mọi người, em ấy vẫn có chút nhẹ nhõm vì chưa ai nhặt lon đủ nhanh để em phải chạy
đuổi theo. Nhưng trong lòng, có lẽ em ấy đang âm thầm hy vọng mình sẽ không phải
trông lon quá lâu nữa.

\ct

Trò ``tạt lon phiên bản \textit{hologra}'' nhanh chóng trở thành một màn tấu hài đúng nghĩa khi mọi người
bắt đầu nghĩ ra những chiêu trò độc lạ để trêu chọc và phá lối chơi của nhau.

Ví dụ như\dots

\il

\dots  lần mà Robo-PON làm người trông lon, tình huống lại xoay chuyển một cách không ai
ngờ đến.
Khi đó, nàng Người máy đứng thẳng, đôi mắt máy lóe lên tia tập trung cao độ. Khi Melchan
vừa cúi xuống nhặt chiếc dép của mình, ngay khi em ấy quay người đi, Robo-PON
bất ngờ rút một cánh tay ra khỏi khớp nối (rõ ràng đây là tính năng của robot), sau đó
dùng chính cánh tay ấy quăng thẳng về phía nàng Dơi.
``Bắt được rồi nha, Mel-chan!'' em ấy thốt lên đầy phấn khởi, trong khi em ấy dường như
vẫn đang đứng ở nguyên tại chỗ.
Mel-chan giật mình la lên, tay nắm chặt chiếc dép nhưng chân đã lùi lại vài bước: ``Đừng
dùng cánh tay của chị ném vào em chứ! Sợ muốn ngất xỉu luôn đây này!''
Rushia-chan từ xa hét lớn: ``Như thế là chơi gian đấy! Không ai bảo được phép tháo cánh
tay ra cả!''
Nàng Người máy nhín vai, nhặt lại cánh tay của mình và gắn lại vào vị trí cũ, tỏ vẻ hơi
bất mãn: ``\textit{hologra} mà! Mọi thứ đều có thể xảy ra cả!''
Tất cả chúng tôi đáp lại: ``\textit{Chơi như bình thường cho bọn em nhờ!}'' Roboco chỉ có thể
lầm lũi quay về trông lon và tiếp tục với công việc của mình. ``Mọi người đúng là xấu
tính\dots''

\ct

Khi Fubuki-chan trông lon, em ấy đeo sẵn một chiếc mặt nạ cáo trông đầy tinh nghịch.
Chỉ cần ai đó tiến lại gần chiếc lon, nàng Cáo trắng lập tức nhảy ra từ chỗ nấp, miệng kêu
``Kon-kon!'' khiến nạn nhân sợ xanh mặt mà lùi lại.

Tôi thở dài: ``Fubu-chan, đừng dọa người chơi khác nữa! Đây không phải là lễ hội Halloween
đâu!''

Fubu-chan bỏ mặt nạ ra, lè lưỡi châm chọc: ``Làm gì có dọa! Em đang chỉ tận dụng lợi
thế sẵn có thôi mà\~{} Kon-kon!''

Tình huống còn trở nên hài hước hơn khi Fubuki mải tập trung đi dọa người khác mà quên
mất chiếc lon ngay trước mặt đã bị Aqua-chan hất bay từ lúc nào.

Nàng Hầu gái nhặt chiếc dép của mình, nhanh chóng lùi về phía sau, cười lớn: ``Hehe!
Lon đổ rồi! Fubuki-senpai thất bại rồi nha!''

Đáp lại những tràng cười của chúng tôi, nàng Cáo trắng bực tức đặt lại chiếc lon về vị trí
cũ, đồng thời phồng má lên đầy phụng phịu.

\ct

Khi Ayame-chan trông lon, thậm chí mọi chuyện còn hỗn loạn hơn nữa. Ngay sau câu
hiệu lệnh của nàng Quỷ sừng, những tiếng vèo vèo của giày dép được ném ra vang lên khắp nơi.

Chiếc giày của Shion-chan tông trúng chiếc lon, nhưng nó lại bay thẳng lên trên, lọt sang
khu vườn của nhà bên kia:``Ấy chêt, nhỡ tung cái lon cao quá rồi.''

Sora-chan nhìn chiếc lon đang bay, thốt lên: ``Nó rơi về phía bên kia bức tường rồi kìa!''

Tôi thở dài: ``Không sao, để anh kiếm cái lon khác-'' nhưng chưa kịp dứt lời, Ayame-chan
đã hăng hái: ``Để em!''

Em ấy trèo lên bức tường gạch, nhảy xuống vườn của nhà bên kia. Tất cả chúng tôi dõi
theo em ấy, chờ đợi một điều gì đó.

``Liệu Ayame-sama sẽ ổn chứ?'' Choco-sensei hỏi, khuôn mặt tỏ vẻ lo lắng. Ru-chan vỗ
vào vai Sensei: ``Em nghĩ chị ấy sẽ ổn thôi.''

Cả đám chưa kịp định thần thì nàng Quỷ sừng đã vội vàng chạy tới, một tay cầm chiếc
lon, một tay\dots\ cầm nguyên một nắm đất ẩm ướt đầy giun đất còn đang ngoe nguẩy.

Choco-sensei, vốn dĩ là một người rất sợ sâu bọ, hét toáng lên, vội nấp đằng sau lưng tôi:
``\textbf{E-Eeek! Sâu bọ kìa!! Kinh quá!!}''

``Giun đất thôi mà em,'' tôi ngoái đầu lại giúp em ấy bình tĩnh lại, để rồi trố mắt ngạc nhiên
trước thứ mà Ayame-chan đang cầm trên tay: ``Mà khoan, Ayame-chan, em hốt kiểu gì
vậy!?''

``Em xới cả đất bên đó để lấy cái lon này mà, yo,'' em ấy đáp, tay vẫn đang cầm nắm đấm
đó, không hề tỏ vẻ sợ sệt gì.

Tôi vỗ trán, ngán ngẩm nói: ``Thế này thì phá vườn của người ta rồi còn gì\dots''

Sora-chan cố gắng giữ bình tĩnh nhưng không giấu nổi nụ cười méo xệch khi nhìn nàng
Quỷ sừng, người giờ đã lấm lem đất cát từ đầu đến chân, như thể vừa hoàn thành một
buổi đào khoai chuyên nghiệp.

Em ấy cười trừ với bộ dạng nhếch nhác của em ấy, hỏi: ``Chị biết là em làm thế để lấy
chiếc lon rồi, nhưng mà đất của nhà người ta mà em xới cả lên thế thì phiền lắm đấy!''

``Nhưng cái lon nó nằm sâu dưới đám cỏ kia mà\dots\ Em làm thế là để mọi người có cái chơi
thôi, yo!'' nàng Oni đáp lại, cố gắng đưa chiếc lon cho tôi.

Shion-chan cười khúc khích, đứng từ xa hét vọng lại: ``Tôi tưởng bà là Quỷ sừng chứ có
phải là Quỷ làm vườn đâu, Ayame!''

Ngay lúc ấy, chủ nhà bên kia bất ngờ ló đầu ra từ cửa sau, vẻ mặt ngạc nhiên khi thấy
một nàng Oni nhỏ nhắn đang đứng trên bờ tường, tay đầy đất. Tôi vội vàng chạy tới giải
thích, cúi đầu xin lỗi lia lịa: ``{\color{red} Chúng cháu xin lỗi, chỉ là lỡ làm rơi cái lon qua đây thôi ạ!
  Chúng cháu sẽ dọn sạch chỗ này ngay lập tức!}''

May mắn thay cho chúng tôi, chủ nhân của khu vườn ấy là một người hiền lành. Bác ấy
nhìn chúng tôi đang chơi đùa rất vui vẻ, chỉ cười rồi phẩy tay: ``{\color{red} Không sao, mấy đứa cứ
  chơi đi. Nhưng lần sau nhớ đừng leo tường nữa, nguy hiểm lắm!}''

Tôi thở phào nhẹ nhõm vì được bác ấy nhắc nhở một cách nhẹ nhàng. Tuy vậy, Ayame-chan
vẫn không khẽ lí nhí: ``Nhưng mà\dots\ em đã lấy được cái lon rồi mà\dots''

Chúng tôi bật cười, quyết định tiếp tục chơi trò chơi này.

\ct

Trò chơi tạt lon được tiếp tục với nhiều lượt chơi nữa. Đến lượt Haato-chan trông lon, em
ấy tỏ ra vô cùng quyết tâm, hai tay chống nạnh, mắt liếc trái liếc phải liên tục như một
``thủ môn'' lão luyện, sẵn sàng bắt bất kỳ ai dám tiến lại gần chiếc lon. ``Đừng hòng qua
mắt được Haachama-sama! Cái lon này là của tôi!''

Mọi người bắt đầu ném giày dép. Bất ngờ, một chiếc giày của Shion-chan trượt qua lon
và bay thẳng vào\dots\ đầu Haato-chan, khiến em ấy loạng choạng suýt ngã, lấy bức tường
gạch cạnh đó để bám víu.

``Ái da, đau quá! Cái quái gì vậy!? Ai ném vào tôi vậy!?'' Haachama-chan hét lên, tay
xoa trán, nơi mà chiếc giày của Phù thủy đáp trúng. Chủ nhân của cú ném ấy cười gian:
``Hì hì, xin lỗi, tay em trơn!''

Đang còn đau điếng, Haato lại nghe thấy tiếng cười khúc khích sau lưng. Quay lại, cô
thấy Ayame-chan và Ru-chan đang lén lút chạy tới lấy giày. Em ấy định vồ lấy họ, nhưng
nàng Oni bất cẩn vấp phải lon, khiến nó bay lên và tông vào\dots\ đầu Haato-chan thêm lần
nữa.

``Đừng có nhắm vào đầu tôi nữa chứ!!'' Haato-chan ôm lấy đầu, tỏ vẻ đau đớn (mặc dù
cái lon đó rất nhẹ). Ayame-chan cười khì: ``Xin lỗi nha, Haachama-sama! Nhưng mà cái
lon này dai như đỉa vậy, cứ nhắm vào chị mãi!''

Lời trêu chọc của em ấy khiến Haato-chan nổi cơn thịnh nộ. Em ấy cầm luôn cái lon, rượt
đuổi theo cô nàng vừa châm chọc đó: ``A\textbf{yame, đứng lại đó!! Mày dám trêu chụy, chụy
  đây sẽ biết tay mày!!!}''

Chúng tôi bật cười với cảnh nhàng Song tính đang đuổi theo nàng Quỷ sừng mưu mô đó.

\ct

Đến lượt Mio-chan trông lon. Em ấy đứng nghiêm túc, trông hệt như một chiến binh bảo
vệ ``lãnh địa.'' Sói đen-chan giơ nanh vuốt ra, trông có vẻ thực sự nghiêm túc với trò chơi
dân gian này: ``Được rồi, các em cứ thử xem! Chị đây không dễ bị qua mặt đâu!''

``Chỉ là trò chơi thôi mà, sao phải làm căng thế\dots'' tôi ngạc nhiên. Tất cả chúng tôi chờ
đợi hiệu lệnh từ em, nhưng mất ngờ thay\dots

*CỐP!*

Chiếc bốt của Fubu-chan được ném ra rất mạnh trúng vào chiếc lon làm phá bĩnh đi bầu
không khí đi. Điều tệ hơn, chiếc lon lăn xa hơn mà mọi người tưởng. Mio-chan tỏ vẻ bất
bình: ``Fubuki! Tôi chưa ra hiệu lệnh mà!'' rồi hí hoáy lao theo nhặt.

Cáo trắng-chan lè lưỡi: ``Xin lỗi, xin lỗi nha Mio\~{}'', dường như em ấy không tỏ vẻ hối lỗi
gì. Sói đen-chan chạy theo chiếc lon đó, nhưng vừa mới cúi xuống thì Akutan ở phía xa ném chiếc dép của mình lần nữa, bay thằng vào\dots\ lưng của Sói đen, đẩy em ngã úp mặt
xuống đất.

``Ái-cha! Aqua, chị vừa làm gì vậy!?'' Mio-chan bò dậy, quở trách người vừa vô tình ném
chiếc dép đi. Nàng Hầu gái ngạc nhiên, giơ hai tay ra phía trước: ``Ể!? C-Chị không cố ý
đâu! Thật mà!''

Trong lúc em ấy còn vừa càm ràm vừa tức tối phủi bụi trên quần áo, bỗng từ đâu một chú
mèo hoang tò mò chạy đến, nhặt luôn chiếc lon trong miệng và chạy mất hút.

``Á! Cái lon của chúng ta!'' Khuyển-chan thảng thốt. Lập tức, Mio-chan, Okayu-chan,
Koro-chan và Fubuki-chan đuổi theo chú mèo đó.

\begin{dialogue}
  \speak{Mio} Đứng lại! Trả cái lon đây!
  \speak{Fubuki} Con mèo này đúng là vụng trộm ghê ha\~{}
  \speak{Okayu} Meo, meo, meo\~{} Mau trả lại cho bọn chị cái lon đi nào\~{}
  \speak{Korone} Gâu, gâu!
\end{dialogue}

Chúng tôi nhìn bốn cô nàng \textit{GAMERS} rượt theo con mèo đó, không ai dám cười với tình
cảnh tréo ngoe này. Ayame-chan thậm chí còn thốt lên: ``Quả thật đây là\dots\ tạt lon phiên
bản truy đuổi!!''

\ct

Sau vài phút, cả bốn người quay lại với chiếc lon méo mó trong tay. Mio-chan tỏ vẻ mặt
thê thảm: ``Trò này đúng là không hợp với chị chút nào! Chị sẽ không chơi nó nữa đâu!!''
Shion-chan châm chọc: ``Sao vừa nãy chị tưởng em hăng máu lắm mà?'' Sói đen-chan
đỏ mặt tía tai, nhưng cũng không buồn cãi thêm câu nào nữa. Chúng tôi lại phá lên cười,
quyết định dừng trò chơi dân gian này lại.

\il

Chúng tôi dự kiến vẫn còn khoảng bốn mươi phút nữa trước khi đến giờ cơm tối, vừa đủ
để có thể chơi một trò chơi nữa.

Trò chơi dân gian cuối cùng của ngày hôm nay không phải đến từ Etsunan chúng tôi, mà
là một trò chơi truyền thống của chính Nhật Bản, nơi mà các thành viên \hll\ sẽ quen
thuộc hơn.

Khi mọi người vẫn đang bàn tán rôm rả về những trò đã chơi từ lúc chúng tôi đến nhà của
ông Niki-san (tôi được Ryuuhou nói vậy) đến giờ, tôi đứng lên, vỗ tay một cái để thu hút
sự chú ý: ``Mọi người, còn một trò cuối cùng nữa mà anh muốn đề xuất. Trò này không
phải của Etsunan, mà là của Nhật Bản các em!''

Nghe đến đây, cả nhóm đều hướng ánh mắt tò mò về phía tôi. Một số người đã đoán ra,
đặc biệt là các thành viên từ các Thế hệ đầu tiên, nhưng vẫn để tôi tiếp tục giải thích.

``Ồ, một trò từ đất nước chúng ta à? Nó là gì vậy?'' Sora-chan hào hứng hỏi tôi. Fubuki-chan mạnh
dạn đoán: ``Có phải là kendama không!?''

Tôi khẽ lắc đầu: ``Không, trò này thú vị hơn nhiều. Cần có hai người chơi-'' Tôi chưa kịp
dứt câu thì nàng Cáo trắng lại thốt lên: ``Hanetsuki (羽根突きhoặc 羽子突き)!''

Tôi khẽ giật mình, gật gù: ``\dots Em đoán nhanh thật đấy.'' Nghe vậy, các thành viên ưa thể
lực - một lần nữa - lại tỏ vẻ hứng thú với trò chơi quen thuộc này.

Riêng chỉ có Aki-chan là vẫn nghiêng đầu khó hiểu. ``Hanetsuki? Là trò gì thế?''

Tôi mỉm cười, rút từ đằng sau ra hai chiếc vợt gỗ sơn màu truyền thống, gọi là hagoita,
cùng một quả cầu lông nhỏ được trang trí bằng lông vũ sặc sỡ. ``Em biết hai cái này chứ?''

Ngay lập tức, Dị giới-chan nhận ra ngay: ``A, em biết hai cái này!'' Tất cả mọi người cũng
dần nhận ra những công cụ quen thuộc cho trò chơi đơn giản này.

Về cơ bản, Hanetsuki là một trò chơi đánh cầu truyền thống của Nhật. Nói nôm na thì
giống như cầu lông hiện đại, nhưng không có lưới. Mọi người sẽ đứng đối diện nhau,
dùng những chiếc vợt này để đánh qua đánh lại quả cầu lông. Ai không đỡ được sẽ bị
phạt.

``Vậy chúng ta sẽ phạt như thế nào đây, Kuon-senpai?'' Mio-chan hỏi. Tôi nhìn ngó quanh
sân vườn, tìm xem có gì để ăn phạt không. ``Thông thường, người thua sẽ bị vẽ một dấu
đen lên mặt bằng mực. Nhưng vì ở đây không có mực, cho nên là\dots''

``Nên sao?'' Subaru-chan hỏi. Tôi tiếp tục ngó quanh sân, mắt tia tới một chiếc nồi lớn ở
trong bếp với phần dưới đã bị áng một màu đen kịt của than: ``\dots ta sẽ ăn phạt bằng nhọ
nồi đằng kia.''

Cả nhóm cười ồ lên khi nghe đến phần phạt, và ánh mắt của các thành viên bắt đầu sáng
lên vì hào hứng. ``Nghe thú vị đó! Em muốn chơi ngay!'' Fubu-chan là người nhiệt tình
nhất trong nhóm.

Ru-chan lên tiếng thắc mắc: ``Nhưng Kuon-senpai này, anh có định đặt ra quy tắc nào đặc
biệt không?''

Tôi gật đầu, cẩn thận giải thích thêm: ``Về cơ bản là cũng giống như bình thường thôi.
Hai người đứng cách nhau khoảng ba mét, không được để quả cầu rơi xuống đất. Ai đánh
trượt, để cầu rơi hoặc đánh ra ngoài, thì người đó thua và sẽ bị bôi nhọ nồi này. Ngoài ra,
đánh cầu quá mạnh hoặc không thể đỡ được thì cũng bị xử thua.''

Tất cả các cô gái đều gật gù, hiểu rõ về sự thay đổi đôi chút trong trò chơi này. Tôi tiếp
tục: ``Nhưng, đây là Hanetsuki phiên bản \textit{hologra}, cho nên luật chơi có thể thay đổi một
chút.''

Tất cả mọi người đều giương mắt nhìn tôi, tỏ vẻ tò mò với luật chơi mới. ``Mọi người
có thể tự do dùng bất cứ kỹ năng gì mà mình có để làm khó đối phương, nhưng miễn là
không làm quá khó.'' Nghe đến đây, cả mười chín người họ càng tỏ vẻ háo hức với trò
chơi này.

``Vậy là trò này cần một chút kỹ năng đấy chứ, không phải cứ đánh bừa là được! Em nói vậy có đúng không, Kuon-sama?'' Choco-sensei
nói, đập tay vào lòng bàn tay kia, tỏ vẻ hiểu rõ cách chơi.

``Chính xác. Trò này không chỉ vui, mà còn mang ý nghĩa may mắn nữa. Theo truyền
thống, nó giúp xua đuổi những điều không tốt lành khi bước sang năm mới. Mấy đứa
chắc cũng hiểu rõ ý nghĩa của nó rồi chứ?''

Tất cả mọi người gật đầu lia lịa. ``Vậy chúng ta bắt đầu thôi!'' Ayame-chan nổi máu lên,
giơ nắm đấm lên cao. Tất cả các cô gái cũng đều hưởng ứng.

Tôi mỉm cười khi thấy mọi người đều hăng hái như vậy. Với sự đồng lòng của cả nhóm,
chúng tôi sẵn sàng bước vào trò chơi cuối cùng, Hanetsuki, với không khí náo nhiệt của
ngày đầu năm mới, ngày mùng Hai Tết.

\ct

Khác với trò chơi Hanetsuki thông thường, lần này chúng tôi sẽ chơi theo cặp, tức là sẽ
cần bốn cái vợt. Tôi đã tính trước được tình huống đó, nên đã chuẩn bị đủ tất cả công cụ
để chơi. Còn về nhọ nồi, tôi xin được các bác ở trong bếp cái nồi to nhất, cũng chính là
cái nồi mà tôi để ý tới ban nãy.

Vì ở đây có mười chín người (tôi và A-san không chơi), nên chúng tôi sẽ chia thành bốn
cặp đấu: Fubuki-Mio đấu với Okayu-Korone, Sora-Mel đấu với Haato-Matsu-Choco, Aqua-Shion
đấu với Noel-Roboco, và Pekora-Rushia-Aki đấu với Ayame-Miko-Subaru.

\il

Ở lượt trận đầu tiên, tôi gọi nó là ``nội chiến GAMERS,'' khi Cáo trắng và Sói đen sẽ phải
đối đầu với Chó và Mèo.

Tôi đưa quả cầu cho Sói đen, rồi đứng giữa hai đội làm quản trò. Cả bốn người nhìn về
phía tôi, chờ đợi tôi ra lệnh.

``\textbf{Bắt đầu!}''

Trận đấu bắt đầu, mọi người đều tập trung theo dõi từng cử động
của cả bốn người trên sân. Mio-chan đưa cho đồng đội quả cầu và phát đầu tiên, đường
cầu bay nhẹ nhưng nhanh, hướng thẳng đến Okayu-chan.

``Cầu dễ quá, thế này thì không đủ làm khó em đâu, Fubuki-senpai\~{}'' Miêu-chan thốt lên,
đỡ cầu một cách điêu luyện bằng một cú đánh chéo góc, đưa cầu bay về phía Sói đen.

``Đỡ chéo à? Chị không ngán đâu, Okayu!''

Mio-chan tung vợt đỡ mạnh, đẩy quả cầu bay thẳng về đối phương - Koro-chan. Khuyểnchan
cười lớn, nhảy lên đón cầu bằng cú đánh xoáy: ``Xem đây! Cú xoáy siêu cấp của
Korone!''

Quả cầu xoay tít, bay nhanh về phía Cáo trắng-chan, khiến em ấy phải nhảy lên để đỡ.
Fubu-chan cười khẩy: ``Xoáy mạnh đấy, nhưng vẫn chưa đủ trình đâu!''

Em ấy xoay người, đập mạnh quả cầu trở lại sân đối phương. Okayu-chan dễ dàng đỡ
được, nhưng bất ngờ đồng đội của em ấy lại có vẻ đoán sai hướng rơi và va vào em ấy.

``Koro-san? Cầu đang rơi về chỗ tôi mà?'' Okayu-chan ngạc nhiên khi nhìn thấy bạn của
mình bị ngã. Korone-chan đứng dậy đầy lúng túng, cười trừ. Cặp đôi nhìn quả cầu rơi
bẹp xuống đất, trong khi cặp đôi Cáo-Sói đập tay khi mang về điểm đầu tiên.

Đồng nghĩa với việc, cặp đôi Chó-Mèo sẽ có cho mình những nét vẽ đầu tiên từ nhọ nồi.
Tôi và A-san phết vài đường trên đáy nồi, rồi bôi lên mặt của cặp đôi nghịch ngợm này.

\begin{dialogue}
  \speak{Kuon} Một đường lên trán của Korone-chan\dots
  \speak{Korone} Ái cha\~{}
  \speak{A-chan} Và một đường lên má của cô, Nekomata-san.
  \speak{Okayu} Tiếc ghê\~{}
\end{dialogue}

Sau khi đã bôi nhọ nồi lên mặt của hai em ấy xong, tôi thông báo số điểm: ``Tỷ số bây
giờ là 1-0 nghiêng về FubuMio!'' Tôi thậm chí còn nhắc nhở thêm: ``Đội nào được năm điểm
trước thì đội đó thắng!''

Nghe vậy, cặp đôi Chó-Mèo liền nhanh chóng quay trở lại trò chơi, với việc Okayu-chan
nhặt cầu lên, chuẩn bị phát lại. Korone-chan hít sâu, tập trung hơn cho lượt tiếp theo.

``Ogayu, lần này chúng ta sẽ lật ngược thế cờ!'' Korone-chan nói với đồng đội của mình.
Okayu-chan khẽ gật đầu, phát quả cầu lên về phía đối phương.

Quả cầu bay vút về phía Mio, nhưng lần này đường bay khá cao. ``Chị đỡ được, chị đỡ
được- Ể!?'' Sói đen thậm chí có phần hơi ngỡ ngàng khi quả cầu phát lên lại quá tầm với
với em ấy.

Mio-chan nhảy lên đỡ, nhưng rồi bất ngờ mất thăng bằng và ngã xuống sân: ``Ái da!
Okayu-chan phát cầu cao thế\dots\ Mà nè, rơi cầu kìa, Fubuki!''

Cáo trắng-chan nhanh chóng lao tới, đỡ cầu ngay trước khi nó chạm đất, cứu nguy cho
đội của mình. ``Mio! Đừng lo, có tôi đây rồi!''

Khuyển-chan không hề bỏ lỡ cơ hội, lao vào với cú đập mạnh, nhưng vì quá nhiệt tình,
em ấy đã đập trượt và ngã bổ nhào xuống đất. Cả hai người của đội đối phương, cùng với
đồng đội của em ấy đều đổ mồ hôi hột với sự hậu đậu của em ấy.

``Ui da! Sao tôi lại đỡ trượt thế này!?'' em ấy lồm cồm bò dậy, nhìn quả cầu chạm xuống
đất với ánh mắt thất vọng.

``Vậy là Koro-chan và Okayu-chan tiếp tục bị bôi nhọ nồi nha!'' tôi và A-san đồng thanh,
phết sẵn hai ngón tay vào chiếc nhọ nồi nữa, và lần này chúng tôi bôi lên những vị trí khác
trên khuôn mặt của bọn họ.

\begin{dialogue}
  \speak{A-chan} Inugami-san bị bôi vào hai bên má này.
  \speak{Korone} *cay cú* Hức\dots
  \speak{Kuon} Okayu-chan dính hai đường ở mắt nha!
  \speak{Okayu} Chán òm à\~{}
\end{dialogue}

Trận đấu tiếp tục với những pha phối hợp lộn xộn nhưng không kém phần hài hước của
cả hai đội, làm khán giả cười nghiêng ngả.

Sau hàng loạt tình huống hài hước, cả hai đội bắt đầu điều chỉnh chiến thuật. Những tiếng
hò reo từ khán giả càng làm tăng thêm nhiệt huyết cho các tuyển thủ. Okayu-chan phát
cầu lần nữa, lần này đánh nhẹ nhưng tính toán đường bay rất hiểm, khiến quả cầu bay
ngang.

``Úi trời! Cầu sát lưới thế này!'' Sói đen-chan than phiền, khi thậm chí em ấy lao tới,
nhưng không thể ngăn được quả cầu chạm xuống đất. Đội OkaKoro có được điểm đầu
tiên.

``Như vậy là đội này có một điểm rồi! GAMERS hàng thật luôn á!'' nàng Khuyển thốt
lên, giơ chiếc vợt gỗ lên cao.

``Hai đứa mới có điểm đầu tiên thôi đấy\dots\ Đội bạn có hai rồi,'' tôi đứng ngoài cuộc, khẽ
nhắc nhở họ. Ngay sau đó, tôi bật cười: ``Nhưng không có nghĩa là cặp FubuMio thoát
khỏi hình phạt đâu nhỉ!''

Tôi và A-san đưa chiếc nồi tới họ, trát nhọ lên khuôn mặt hai người họ, như cách mà
chúng tôi làm với đội kia.

\begin{dialogue}
  \speak{Mio} *phụng phịu* Tại Okayu đập khó quá chứ bộ\dots
  \speak{Fubuki} Bọn này sẽ không chịu thua đâu!
\end{dialogue}

Lần này, Fubuki-chan phát cầu, để cho nàng Khuyển ở đội kia đỡ. Nhưng lần này, thay
vì đánh mạnh như cách mà em ấy định làm, Korone-chan tung một cú lốp cầu cao hòng
đưa cầu ra khỏi tầm với của cả hai đối thủ.

``Haha! Đừng bao giờ coi thường sức mạnh của đội OkaKoro này!'' em thốt lên, tung quả
cầu lên cao chót vót.

Ở phía đối phương, dường như Fubuki-chan đã chuẩn bị bài sẵn, nên trông có vẻ không
bị bất ngờ. Em ấy bật cao, xoay người trên không, đánh một cú đập cầu cực mạnh xuống
sân của Okayu-Korone khiến họ không trở tay kịp, mang về thêm một điểm nữa cho đội.

Nàng Miêu nhìn đòn đánh của tiền bối mình, tỏ vẻ khâm phục: ``Fubuki-senpai đánh đòn
này khá gắt đấy\~{}'' Cáo trắng-chan tỏ vẻ tự cao: ``Mấy đứa biết chị là trưởng nhóm của
GAMERS mà, đúng chứ! Không dễ gì đùa với chị đâu!''

Tỷ số giờ là 3-1. Trận đấu tiếp tục với những pha xử lý kỹ thuật đỉnh cao giữa hai đội.
Quả không hổ danh là các thành viên trong nhóm \textit{hololive GAMERS} có khác.

Mio-chan lần này giữ vị trí phòng thủ, tỏ vẻ tự tin hơn hẳn. Em ấy đã có thể tận dụng
chiều cao để chắn các cú đập cầu hiểm hóc từ đội đối phương, nhất là từ một Okayu mưu
mẹo.

``Chị không phải chỉ là sói đâu, còn là bức tường bất bại nữa!'' em hét lên, cố gắng chống
trả các đòn từ họ. Okayu-chan không hề chịu lép vế, sử dụng chiến thuật đánh vào góc
sân trống, buộc cho bên tấn công là Fubu-chan phải di chuyển liên tục: ``Đỡ thử cái này
xem nào, Fubuki-senpai!''

Cả hai đội liên tục so kè lẫn nhau. Đã có những lần mà Mio-chan và Fubuki-chan dồn dập lên tấn công, nhưng rồi Okayu-chan và Korone-chan cũng đáp lại bằng những đòn
khó đánh không kém.

Không bao lâu sau, tỷ số của trận đấu này là 4-4. Trận đấu đang bước vào giai đoạn căng
thẳng nhất. Chỉ còn lại một lượt phát cầu nữa.

Mỗi bên đều đã toan tính cho lượt quyết định này. Ai nấy đều tỏ ra căng thẳng trong một
trận đấu đánh cầu này.

Bất giác, Korone-chan nảy ra một ý tưởng mới, khiến đồng đội bối rối: ``Ogayu! Lùi lại,
một mình tôi sẽ đấu với hai người họ!''

Đây có lẽ là quyết định khá bất ngờ và có phần mạo hiểm từ nàng Khuyển đang hăng máu
này. Miêu-chan nhìn em ấy, tỏ vẻ nghi hoặc: ``Bà chắc chứ, Koro-san?''

Không cần đợi em ấy có câu trả lời, Korone-chan cầm chắc vợt bằng cả hai tay, tạo thủ thế như đang chơi một trò chơi hành động. Ngay khi Mio-chan phát cầu, Koro-chan đã
không do dự mà phát lại quả cầu đó.

``Một mình chống cả thế giới à? Ghê đấy!'' Fubuki-chan nói, tỏ vẻ hăng máu hơn hẳn.
``Để xem em có thể chống trả được với các đòn siêu tốc độ này không!''

Ngay lập tức, cặp đôi Fubuki-Mio liên tục phát nhanh cầu về nàng Chó, nhưng bất ngờ
thay, đối phương có thể tung được mà không gặp chút khó khăn gì.

``Hay dữ!?'' cả ba người họ trố mắt ngạc nhiên. Đáp lại, nàng Chó nâu nói: ``Có giỏi thì
phát nữa đi!''

Cứ như vậy, cặp đôi Cáo-Sói liên tiếp phát nhanh cầu về nàng Chó, và Koro-chan cũng
tung liên hoàn các cú đánh cầu, đồng thời chạy khắp sân, khiến cho đối phương phải vất
vả chống đỡ.

\begin{dialogue}
  \speak{Fubuki} Bọn này sẽ không chịu thua đâu!
  \speak{Korone} Gâu, gâu! Đúng rồi! Đây là đòn liên hoàn của em đây!
  \speak{Mio} *phụng phịu* Tại Okayu đập khó quá chứ bộ\dots
\end{dialogue}

Mặc dù thế, biến cố của trận đấu đã xảy ra: trong một lúc nhiệt tình, Korone-chan đã
không còn giữ được độ chính xác, dẫn đến một cú đánh sai địa chỉ bay thẳng về phía khán
giả, nơi mà các cô gái khác đang chứng kiến trận đấu của họ.

Cả Okayu-chan và Korone-chan nhìn quả cầu bay đi với những cặp mắt ngỡ ngàng. Đúng
lúc đó, đội đối phương mừng như bắt được vàng.
``Cảm ơn vì món quà, Korone!'' Mio-chan mừng rỡ, hừng hực khí thế. ``Fubuki!''

``Có liền!'' Cáo trắng nhanh chóng tận dụng cơ hội, tung cú đập cực mạnh xuống sân đối
thủ, ghi điểm cuối cùng cho đội.

Tôi thổi còi (mà tôi mượn từ Subaru-chan), ra hiệu kết thúc trận đấu. A-san tuyên bố: ``Tỷ
số chung cuộc 5-4, nghiêng về đội Shirakami-san và Ookami-san!''

Đội Cáo-Sói mừng rỡ, mặc dù mặt họ kha khá những vết nhọ nồi trên khuôn mặt. Tôi tiến
tới hai người, đưa cho họ chiếc nồi rất to mà tôi và A-san đã dùng tới.

``Tại sao anh đưa cho bọn em thứ này?'' Mio-chan hỏi. Tôi cười, đáp: ``Phần thưởng của
hai đứa đấy.''

Fubu-chan nghiêng đầu: ``Là sao?'' A-san tiếp lời tôi: ``À, đội thắng sẽ có quyền bôi nhọ
nồi bao nhiêu tùy thích lên mặt của đội thua cuộc.''

Vừa dứt lời, cả hai người họ liền nhận chiếc nồi, đồng thời đồng thanh: ``Cảm ơn hai anh
chị nhiều!'' và lao tới cặp đôi còn lại, những người đang vừa ngồi bệt xuống sàn vừa thở
hổn hển.

Khi họ ngước mặt lên, họ thấy Fubuki-chan và Mio-chan đã leo lẻo những ngón tay được
trát sẵn nhọ nồi, nở những nụ cười gian ác nhưng hiền hậu.

``Chuẩn bị chịu phạt nào hai đứa!'' Mio-chan thốt lên, theo sau là Fubu-chan: ``Dám thách
đấu với tụi này\dots\ quả đúng là sai lầm!''

Và sau đó\dots\ chắc các bạn cũng biết rồi đấy. Cặp đôi thắng cuộc đang bôi lên mặt cặp đôi
còn lại với vẻ mặt thích thú, trong khi những người đang bị bôi nhọ nồi thì cười lăn lộn.

\begin{dialogue}
  \speak{Fubuki} Chết nè, chết nè!
  \speak{Korone} *cười* Ahahaha! Buồn quá, Fubuki-senpai!
  \speak{Mio} Okayu-chan nhận lấy này!
  \speak{Okayu} *cười* Ôi trời ơi, hai người ơi\dots\ Mặt mèo của tôi trở nên dị hợm rồi!
\end{dialogue}

Nhìn bốn người họ vờn nhau vui vẻ, chúng tôi cùng với các khán giả đều vui lây, vì cuối
cùng trận đấu căng thẳng đó lại kết thúc bằng niềm vui.

Cuối cùng, bốn người họ bắt tay nhau đầy hữu hảo, khẳng định rằng nhóm của họ vẫn sẽ
trường tồn.

\begin{dialogue}
  \speak{Mio} Hai đứa cũng chơi khá tốt đấy.
  \speak{Fubuki} Công nhận là hai đứa cũng ghê gớm lắm ha\~{}
  \speak{Okayu} Hai người đúng là những bậc game thủ chuyên nghiệp thật á\~{}
  \speak{Korone} Bọn em vẫn sẽ không bỏ cuộc đâu! Lần sau bọn em sẽ chiến thắng hai
  người!
\end{dialogue}

Như vậy, cặp đấu đầu tiên đã kết thúc. Bốn người họ sau đó ngồi nán lại để chuyển sang
các cặp đấu tiếp theo.

\ct

Trận tiếp theo, là đội Sora-Mel đấu với bộ ba Haato-Matsuri-Choco. Tôi biết, hai đấu ba,
nhưng ít ra còn tốt hơn việc sẽ phải để một người cho ra rìa. Mặc dù bề ngoài, đây là một
trận đấu không cân sức, nhưng thực ra trận đấu không hề dễ dàng cho bất kỳ bên nào, bởi
mỗi người đều sở hữu những chiêu thức ``đặc trưng'' có thể gây khó dễ cho đối phương.

Ngay từ đầu, Haato-chan tung cầu về phía hai cô nàng được đánh giá là hiền lành nhất
với một nụ cười kỳ quái, khiến chúng tôi cũng cảm thấy lạnh gáy. ``Haachama-mode, bắt
đầu đây\~{}''

Quả cầu bay xoáy và đổi hướng không theo bất kỳ quy luật nào mà chúng tôi thường thấy:
lúc thì bay thẳng, lúc thì lại chúi xuống, thậm chí nó còn\dots\ tạo các đường zigazag khiến
cho việc đoán định đường bay như một cực hình.

``S-Sora-senpai! Cẩn thận!'' đồng đội của em ấy thảng thốt kêu lên, báo động cho em ấy.
Nhưng rồi, Sora-chan mim cười nhẹ nhàng, ánh mắt sáng lên một cách đáng sợ.

Không lẽ nào\dots\ Sora-chan đang xuất ra một chiêu thức đáng sợ nào đó!?

``Để chị xử lý.''

Cô nhảy lên và tung một cú đập cầu mạnh đến mức quả cầu phát ra âm thanh xé gió, bay
thẳng về phía Haato-chan. Cả ba người họ không kịp phản ứng với tốc độ kinh hoàng của
quả cầu kia.

\begin{dialogue}
  \speak{Matsuri} S-Sora-senpai!?
  \speak{Haato} Quả cầu nhanh thế-
  \speak{Choco} Cẩn thận kìa, Haato-sama!
\end{dialogue}

Quả cầu đập trúng vợt của Haachama-chan, bật ngược ra ngoài sân, làm em ấy ngã nhào
xuống đất, thốt lên ``Hự!'' một tiếng. Đội ba người họ đã mất điểm đầu tiên.

Sau cú đánh nhanh như chớp đó, cả ba người họ - tất nhiên - đều bị chúng tôi quệt nhọ
vào mặt. Đội nhanh chóng lấy lại tinh thần, với Choco-sensei là người phát cầu.

``Đừng lo, có cô ở đây rồi.''

Em ấy giơ cầu lên trên cao. Nhưng thay vì dùng vọt gỗ như lẽ thường, em ấy lại\dots\ ép quả
cầu vào ngực, rồi hất nó lên trời trước sự ngỡ ngàng của đội bạn và các khán giả cũng như
trọng tài như chúng tôi.

``Ể!? Dùng thế cũng được à!?'' Mel-chan thảng thốt, nhìn quả cầu đang bay về phía hai
người họ.

``Như thế chẳng phải là phạm luật sao?'' Thần tượng-chan tỏ vẻ lo ngại, trong khi đang ở
tư thế hứng lấy quả cầu kia.

Tôi ra dấu cho trận đấu tiếp tục, bằng cách lắc đầu, hàm ý rằng cú tâng cầu đó là hợp lệ.
Choco-sensei thấy vậy, mỉm cười đầy giann xảo: ``Miễn là đáp được sân của hai chị thì
đâu có phạm luật, nhỉ?''

Cả hai người của đội đối phương cố gắng đỡ lấy quả cầu đang rơi đó, nhưng dường như
vẫn còn lơ đễnh sau màn phát cầu vừa rồi, nên cả hai người họ không may va vào nhau
và ngã ra đất, nhìn quả cầu đáp xuống với vẻ bất lực.

Tôi thổi còi, ra hiệu cho Sora-chan và Mel-chan phát cầu ở lượt tiếp theo. Tỷ số bây giờ là 1-1, và tất nhiên đội đỡ hụt lại bị trát nhọ lên mặt.

Lượt tiếp theo, Mel-chan là người tung quả cầu lên. Ngay lập tức, Haato-chan đáp trả
bằng một cú đánh mạnh, đi kèm với tiếng hét ``Haachama-chama!'' thương hiệu của em
ấy.

Quả cầu bay lên cao một cách kỳ lạ, khiến cả Sora-chan và Mel-chan rơi vào thế bí, khi
quả cầu dường như sẽ rơi vào mép của sân. Đúng lúc này, Dơi-chan đã chơi một chiêu
mà khiến chúng tôi trố mắt ra: treo người ngược lên một cành cây gần đó!

``Bắt được rồi! Treo người lên là lợi thế của tôi mà!''

Mel-chan đánh đánh cầu trả về phía đội bạn, với một lực mạnh tương đương. Dù với
lợi thế có sân lớn hơn và có tận ba người, đội Haato-Choco-Matsuri cũng đã phải luống
cuống chạy khắp sân để đỡ được quả cầu đó.

Cả hai bên tiếp tục đánh qua đánh lại như vậy, cho đến khi một cơn gió thổi qua. Ngay
lúc này, Haato-chan reo lên: ``A! Mel-chan không mặc quần lót kìa!'' Câu nói bất ngờ
của Song tính-chan khiến cho cả hai cô gái của đội kia giật mình, với nàng Ma cà rồng
vội che lại phần hạ bộ bằng bộ kimono đang chổng ngược của mình. Mel-chan đỏ mặt,
ngượng ngùng nói: ``Haato-chan!! Đừng có làm mất mặt mình chứ!'' trong khi đồng đội
của em ấy cũng che mặt lại vì xấu hổ.

Đến cả các cô gái ngồi làm khán giả và cả tôi đều ngượng chín mặt. A-san và tôi vội
ngoảnh ra chỗ khác, quyết không liếc vào ``chỗ đó''. Tôi chỉ có thể nghe thấy tiếng *bụp*
của quả cầu được đáp xuống ở sân của Sora-chan và Mel-chan.

\begin{dialogue}
  \speak{Mel} Đ-Đừng có nhìn mình chứ\dots
  \speak{Kuon} Ờm\dots\ Tỷ số bây giờ là 2-1, nghiêng về nhóm ba người.
  \speak{A-chan} Yozora-san, cô có thể xuống được không\dots?
\end{dialogue}

Nàng Thần tượng và nàng Ma cà rồng chỉ có thể thất thểu chịu phạt bằng việc để tôi và
A-san vẽ nhọ nồi đen vào mặt. ``L-Làm ơn đừng có chơi bài chổng ngược thế khi đang
mặc váy chứ\dots'' tôi nói trong sự thở dài.

``E-Em biết mà!!'' Mel-chan nói, vẫn với khuôn mặt đỏ hỏn.

Trận đấu tiếp tục, với Sora-chan là người phát cầu. Nhưng rồi, ngay khi Matsuri-chan
định đỡ quả cầu, dường như em ấy đã chuẩn bị sẵn một bài nào khác.

Em ấy rút từ đâu đó một mảnh pháo giấy khá nhẹ. ``Chơi hội chợ thì phải có lửa chứ!''

Em phóng pháo giấy ra, ngay sau khi đỡ quả cầu được phát. Quả pháo giấy phát nổ, những
quả tua rua bùng ra, làm rối mắt cả hai người họ. Tưởng rằng hai người họ sẽ không thể
đỡ nổi quả cầu trong mớ giấy đó\dots

``Matsuri-chan vẫn còn non lắm.''

\dots mắt của em ấy lóe lên, phát hiện ra quả cầu trong mớ đó. Em ấy sử dụng một chiêu
đặc biệt. Cú đánh của cô nhẹ nhàng nhưng khiến quả cầu quay vòng liên tục trên không
trung, khó mà phán đoán được quỹ đạo.

Quả cầu vút đi khỏi mớ giấy sắc mày, bay quay vòng trên sân đối phương. Haato-chan
trố mắt ngạc nhiên: ``C-Cầu trông như bị ma ám vậy!''

``Mau đỡ nó đi!'' Matsuri thảng thốt, theo sau là Choco-sensei đang cố đoán vị trí quả cầu
sẽ rơi xuống.

Cả ba người đã cố gắng đỡ nó, nhưng đều đánh trượt, làm quả cầu rơi xuống ngay ở chân
họ. Đội Sora-Mel đã gỡ hòa được 2-2.

``Cái này gọi là nghiệp quật đấy!'' tôi ở bên ngoài bình luận, với ba cô nàng đang thẫn thờ,
chuẩn bị có thêm những vết nhọ nồi nữa lên mặt. ``Chơi tốt lắm, Sora-chan, Mel-chan!''

*TUÝT!*

Ở lượt tiếp theo, Haato bật ``chế độ Haachama'' lần nữa, nhưng lần này cô còn\dots\ cầm hai
vợt trên tay.

``\textbf{Các người không thể hạ gục được Haachama này đâu!}'' em ấy tuyên bố đầy mạnh
mẽ, phát cầu cho đội đối phương. Haato-chan đánh cầu với tốc độ không tưởng, nhưng
lại vô tình khiến quả cầu văng trúng một trong hai người đồng đội, khiến đội mình không
thể phối hợp.

``Haato-chan! Bà làm cái gì vậy!?'' Matsuri-chan xoa vào lưng mình, nơi mà quả cầu vừa
mới đập vào, còn nàng Y tá cũng ngỡ ngàng không kém: ``Em định tự làm khó mình đấy
sao? Em có phải là Korone-sama đâu!''

Như vậy, tỷ số là 3-2, nghiêng về đội của hai cô nàng. Tới lúc này, Lễ hội-chan quyết
định chơi lớn.

Em nhảy lên không trung, xoay người như một nghệ sĩ xiếc, tung cú đánh trả cầu mạnh
mẽ về phía Mel-chan, khiến cho đối phương bị bất ngờ.

``Lễ hội phải rực rỡ chứ nhỉ\~{}!''

Đòn đánh này đã giúp đội của em ấy đưa trận đấu về thế cân bằng một lần nữa, khi Melchan
không thể nào đỡ được quả cầu đó.

Trận đấu tiếp tục với sự căng thẳng tột độ, cả hai bên đều tung hết những quân bài tẩy của
mình để đánh bại đối phương, biến trận Hanetsuki này thành một trận đấu nhiều chiêu trò nhất.

Ở lượt quyết định, Haato-chan tung một cú lốp cầu cao nhằm gây áp lực cho cả hai cô
nàng, nhưng Sora-chan đã nhảy lên với một cú đập cầu mạnh uy lực, ghi điểm cuối cùng.

``Và đây là\dots\ đòn kết thúc!'' *BỐP!*

Tiếng còi từ tôi vang lên, ra hiệu kết thúc trận đấu. Tỷ số chung cuộc cho trận đấu này là
5-3, nghiêng về đội Sora-Mel.

Đội Haato-Matsuri-Choco mặc dù đã thua cuộc, nhưng họ trông không có vẻ cay cú gì,
ngược lại đang nở những nụ cười mãn nguyện.

\begin{dialogue}
  \speak{Choco} Mệt thế này thì chắc cô giảm được vài ký đấy\dots
  \speak{Haato} Hai người chơi tốt lắm! Nhưng lần sau, tôi sẽ phục thù, nhớ đó!
  \speak{Matsuri} Vui thật đó! Mà\dots\ tôi ngửi thấy mùi thịt nướng kìa!
\end{dialogue}

Tất cả chúng tôi bật cười với lời cảm thán của em ấy. Mà nhắc đến thịt nướng, không biết
bữa tối hôm nay có gì nhỉ\dots?

Dù sao thì, tôi quyết định để Sora-chan và Mel-chan bổi nhọ nồi thỏa thích lên mặt của ba
người. Nhìn cả năm người họ đùa giỡn với nhau, thậm chí Matsuri-chan và Haato-chan
còn bôi lại lên mặt hai người họ, trông không giống như hai đội đã có một trận đấu kịch
tích và mưu mô tới mức nào.

Ta sẽ chuyển đến cặp thi đấu thứ ba!

\ct

Mọi người tập trung cổ vũ nhiệt tình khi hai đội bước vào sân. Đó là trận đáu giữa Aqua-
Robooc và Shion-Noel. Tôi gọi nó là ``cặp bạn thân đối đầu nhau'' và ``hai người đại ngốc
nhất đấu nhau.''

Một bên là Aqua-chan và Robo-PON, đại diện cho sự nhanh nhẹn (dù không biết có đúng
không) và công nghệ hiện đại, bên còn lại là Shion-chan và Noel-PON, kết hợp giữa phép
thuật và sức mạnh cơ bắp.

Hầu gái-chan bắt đầu khởi động cơ tay, tỏ vẻ đầy quyết tâm với đối thủ của em ấy: ``Hầu
gái đỉnh nhất sẽ không thua đâu! Sẵn sàng chưa, Shion?''

Shion bật cười đầy tự tin, tay xoay xoay cái vợt gỗ của mình: ``Bà nghĩ phép thuật của tôi
là thứ mà mấy người có thể đỡ được sao? Tôi sẽ cho hai người xem!''

Roboco-chan mỉm cười nhẹ nhàng, đôi mắt sáng lên những màu ánh sáng rực rỡ: ``Chế
độ tăng cường sức mạnh đã kích hoạt. Thử xem ai nhanh hơn nào\~{}''

Noel-chan siết chặt cây vợt trong tay, thể hiện rõ cơ bắp của mình: ``Đừng lo, em sẽ bảo
vệ sân nhà. Shion-senpai, chị cứ tập trung tấn công đi!''

Tôi thổi còi, ra hiệu trận đấu bắt đầu, và ngay từ lượt đầu tiên, Shion-chan đã tung một
đòn phép thuật xoáy khiến quả cầu biến thành một quả bóng lửa nhỏ.

``Ể!? Shion-chan, bà làm cái trò mèo gì vậy!? Nóng quá!!'' nàng Hầu gái thốt lên, cố
tránh quả cầu lửa ra xa. Đúng lúc này, đồng đội của em ấy nhanh chóng di chuyển, dùng
cánh tay cơ khí đập mạnh quả cầu qua lưới, làm nó trở lại trạng thái bình thường.

``Chế độ làm nguội kích hoạt. Xin lỗi nhé, Shion-chan!''

Chưa kịp để nàng Phù thủy ngạc nhiên, Noel-chan đã thốt lên: ``Để em đỡ! Ya!'' Hiệp sĩ
kịp thời đỡ lấy bằng cú đánh mạnh như búa bổ, đưa quả cầu về phía đội đối phương.

Cú đánh quá mạnh khiến cho quả cầu bay lên rất cao, làm Hầu gái-chan hớt hải tìm chỗ
đón quả cầu rớt xuống. ``C-Cái gì thế này!?''

Em ấy cố gắng nhảy lên, nhưng lại lỡ tay đánh trượt, làm quả cầu rơi xuống sân. Căp đôi
Shion-Noel có được điểm đầu tiên.

``Noel-chan thế mà cũng cừ đấy!'' Shion-chan nói, giơ ngón tay cái lên. Nàng Hiệp sĩ
cười gượng: ``Chị cứ quá khen\dots''

Ở bên kia, Aqua-chan phồng má lên, trong khi Roboco-chan thì cười khì: ``Được rồi, lượt
tiếp theo chị sẽ đỡ được!''

Lượt tiếp theo, cô nàng Hầu gái thể hiện sức mạnh của một ``hầu gái game thủ'' tối thượng.
Không hề nản chí, em tung cú phát cầu với một tốc độ nhanh đến mức tạo thành âm thanh
*VÚT!* trong không khí.

``Hầu gái thì phải phục vụ nhanh cho các chủ nhân, \textbf{đúng chứ!?}'' em ấy thốt lên như vậy,
ngay khi sử dụng vợt của mình đập quả cầu đó.

Noel-chan may mắn thay đã đỡ được quả cầu, nhưng nó lại bật quá cao. ``Chết rồi!''
Roboco-chan tận dụng thời cơ này để thực hiện cú đập mạnh, khiến cho đối phương
không kịp trở tay. ``Nhận lấy!''

Cú đập mạnh đó cũng giúp đội gỡ hòa được 1-1. Bây giờ, cả hai đội đã có những vệt nhọ
nồi trên mặt ở ngay hai lượt chơi đầu tiên.

Lượt đánh tiếp theo, Phù thủy-chan tiếp tục sử dụng phép thuật của mình làm lợi thế. Lần
này, em ấy biến quả cầu thành một quả cầu nước.

``He he. Tưởng gì chứ, quả này bọn chị đỡ được!'' nàng Người máy cười một cách đầy tự
tin khi quả cầu đã ngay sát họ.

Tuy nhiên, điều mà họ không hề ngờ tới là quả cầu nước đó đã xối nước xuống ngay dưới
chân họ, khiến cả hai người trượt ngã vì sân trơn khi cố đỡ nó.

Roboco thốt lên, tỏ vẻ bất bình với đội bạn: ``Sân trơn quá! Shion-chan, thật là không
công bằng\~{}!''

Đáp lại, Shion-chan chỉ nở một nụ cười đầy ẩn ý: ``Ai bảo chị không chuẩn bị chứ? Ma
thuật luôn chiến thắng! Có đúng là bọn em chơi bẩn không, Kuon-senpai?''

Em ấy cùng đồng đội của mình quay mặt về phía tôi, vẫn nở một nụ cười gian xảo. Tôi
lắc đầu xua tay, ra hiệu rằng đây không phải lỗi, đồng thời thổi còi ra hiệu cho đội của
Akutan và Robo-PON phát cầu.

``Thấy chứ? Đến cả trọng tài cũng nói đây không phải là chơi bẩn mà.''

Aqua-chan và Roboco-chan nhìn tôi bằng những ánh mắt thất vọng, tỏ vẻ như muốn phản
đối quyết định của tôi. Nhưng tôi lắc đầu và ra hiệu họ phát cầu lên.

Không còn cách nào khác, hai cô nàng phải làm như vậy trong sự ấm ức. Tỷ số bây giờ
là 2-1 dành cho đội Phép thuật và Cơ bắp.

Lần tung cầu tiếp theo, Roboco-chan tiếp tục sử dụng cơ thể rô-bốt của mình để gây khó
cho đối phương.

``Hãy xem công nghệ của chị có đọ được phép thuật của em không!''

Robo-PON kích hoạt chế độ nhảy cao, tung cú đập khiến quả cầu bay với quỹ đạo cong
queo và khó đoán. Cả hai người của đội đối phương ngớ người nhìn quả cầu đó, không
biết phải đỡ kiểu gì.

\begin{dialogue}
  \speak{Noel} Ể\dots\ Ể?
  \speak{Shion} Này! Roboco, chị đập cầu kiểu gì vậy!?
\end{dialogue}

Đáp lại những con mắt ngơ ngác nhìn quả cầu bay qua cả hai mà không thể đỡ được,
Akutan chỉ cười mỉm: ``Miễn là thắng được mà, đúng chứ?''

Tỷ số bây giờ là 2-2. Trận đấu bây giờ đã bắt đầu trở nên gay cấn hơn, khi thành viên cả
hai đội bắt đầu tung ra những đòn phép và đòn cơ để triệt hạ nhau.

Lần này, Noel-chan đập quả cầu với sự nhẹ nhàng, và tất nhiên Akutan đỡ được mà không
có vấn đề gì.

Nhưng rồi, dường như Hiệp sĩ-chan bắt đầu đẩy nhanh tốc độ, khiến cho nàng Hầu gái
bắt đầu phải chạy vào quanh sân, trong khi Roboco-chan chỉ đứng đó mà không biết phải
xử trí như thế nào.

``Cứ từ từ mà chạy theo quả cầu, Aqua-senpai!'' Hiệp sĩ-chan vừa nói, vừa liên tiếp đánh
cầu trả về cho nàng Hầu gái đang thở hồng hộc. Sức bền của em ấy quả thật là không thể
so bì được với cô nàng nội hướng này.

``Cố lên nào, Aqua-chan! Chị tin tưởng vào em!'' đồng đội của em ấy vẫn luôn miệng cổ
vũ, trong khi bản thân nàng Hầu gái vẫn đang mệt bở hơi tai, gắt lên: ``Đừng có đứng đó
nữa! Mau giúp em đỡ hộ đi!''

Nhân lúc em ấy đang loạng choạng chạy đi chạy lại, bạn của em ấy đã tung một cú đánh
xoáy khó chịu, khiến cho cả hai người không thể nào đỡ nổi, với Hầu gái-chan ngã chỏng
quèo xuống đất thở dốc. Aqua-Roboco bị rơi vào thế bị dẫn trước 3-2.

Tuy nhiên, bàn thua đó làm cho Akutan không hề nản chí. Em ấy từ từ đứng dậy, lau mồ
hôi toát ra trên rán, khởi động ``chế độ stressing tối đa.''

``Lần này chị sẽ không bị em dễ dụ khị như vậy đâu, Noel! Lên đi!'' em ấy thốt lên, đồng
thời không quên quay về đồng đội: ``Roboco-senpai, chị cũng mau phụ em đi, đừng có
đứng ngây ra thế!''

Cũng giống như lượt trước, Hiệp sĩ-chan tiếp tục dùng sức bền của mình để ép Akutan.
Nhưng với chế độ này, nàng Hầu gái có thể chạy với tốc độ chóng mặt, có thể bắt kịp
được với sức bền của đối phương, đồng thời phối hợp ăn ý với đồng đội (vốn cũng có sức
bền tốt) để ép ngược lại.

\begin{dialogue}
  \speak{Aqua} Sẵn sàng chưa, senpai?
  \speak{Roboco} Được rồi\dots\ Đỡ lấy! Ya!
\end{dialogue}

Roboco-chan đập quả cầu rất mạnh với tốc độ cực nhanh, tới mức mắt thường không thể
nhìn được.

\begin{dialogue}
  \speak{Shion} Đ-Đòn này là!
  \speak{Noel}Cái quái gì-!?
\end{dialogue}

\textbf{*RẦM!*}

Quả cầu đáp xuống ngay đầu của nàng Phù thủy, khiến cô nàng xanh mặt vì sợ. Đồng đội
của cô cũng không khá khẩm hơn: Noel-chan thì run rẩy. Tỷ số bây giờ là 3-3, và cả hai
đội đã chi chít vết nhọ nồi trên mặt, từ quanh mắt cho đến kín trán.

Lượt đấu tiếp theo, Roboco-chan tiếp tục thử nghiệm một trong những tính năng mà em
ấy đang có: tự động nhắm mục tiêu. Lần này, không ai nói một câu gì, chỉ tập trung vào
đánh cầu.

Robo-PON nhắm vào một trong những điểm ở dưới mặt sân, căn chỉnh đầy kỹ lưỡng,
đo cả tốc độ gió. Tuy vậy, công việc đo đạc của cô khiến cho không chỉ Shion-chan và
Noel-chan, mà còn cả khán giả bắt đầu sốt ruột. Tôi và A-san vẫn đang theo sát cô nàng
Rô-bốt hậu đấu ấy.

``Này, chừng nào chị mới chịu đánh thế, Roboco?'' Phù thủy-chan bắt đầu mất kiên nhẫn.
Hiệp sĩ-chan tiếp lời: ``Chỉ là trò chơi dân gian thôi mà, đừng có mà câu giờ như thế chứ,
senpai!''

Nhưng ngay khi hai người của đối phương than phiền, Robo-PON bất ngờ tung cú đánh
mạnh vào góc sân của Shion-Noel trước sự ngỡ ngàng của cả hai tuyển thủ đội bạn và cả
khán giả.

Tôi thổi còi, ra hiệu cho đội của hai cô gái đang đứng trời trồng kia phát cầu, đồng nghĩa
với việc đội Aqua-Roboco được điểm.

``N-Này! Đánh bất ngờ thế ai mà đỡ được chứ!'' Noel-chan thốt lên, tỏ vẻ không phục.

Akutan cười lớn: ``Ai bảo hai người không tập trung vào cơ! Đội chị được điểm rồi nha\~{}''

A-san lên tiếng: ``Về lý thuyết thì họ nói đúng đó, Shirogane-san. Phải luôn sẵn sàng cho
mọi tình huống chứ.''

Cả Phù thủy-chan và Hiệp sĩ-chan đều chỉ có thể rên rỉ trong sự bực dọc với kiểu chơi có
phần không công bằng này.

Nhưng biết sao được, Hanetsuki phiên bản \textit{hologra} mà, việc gì cũng có thể xảy ra mà, nhỉ?

Lượt tiếp theo, Shion-chan quyết định sử dụng phép thuật của mình để làm quả cầu biến
mất, làm cho hai người họ lúng túng, với nàng Hầu gái thốt lên: ``Ê! Shion-chan, bà chơi
gì kỳ vậy!''

Phù thủy-chan nhoẻn miệng cười: ``Thì vừa nãy đội bà cũng chơi bài bẩn đấy. Tôi chơi
lại đúng bài hai người thôi.''

Nhưng có vẻ như, phép thuật có mạnh đến mấy cũng không thể chiến thắng được công
nghệ. Quả cầu thực sự biến mất, nhưng bóng của nó vẫn còn. Roboco-chan đã dự đoán chính xác điểm rơi, đập cầu trả lại đối phương. ``\textbf{Đỡ lấy!}''

Quả nhiên, vì quả cầu không thể nhìn thấy trước mắt, nên Noel-chan không thể nào đoán
được nó ở đâu. ``Ơ-Ơ, này!? Shion-senpai, chị định làm khó em sao!?''

Khi nó xuất hiện trở lại thì đã đáp trên sân của họ từ khi nào, với điểm số cuối cùng dành
cho đội Người máy-Hầu gái. Hồi còi kết thúc trận đấu từ tôi vang lên, với chiến thắng
thuộc về Akutan và Robo-PON.

Ngay sau hồi còi mãn cuộc, Aqua-chan thở hổn hển, ngồi xụp xuống sân, nhưng vẫn nở
một nụ cười nhẹ: ``Ta thắng rồi\dots\ Tôi thắng bà rồi đấy, Shion-chan!''

Trong khi đó, Roboco-chan đã bôi sẵn tro nồi từ hai bàn tay của em ấy từ khi nào. Em
mỉm cười, ngoáy ngoáy hai tay: ``Chuẩn bị sẵn sàng ăn phạt nha\~{}''

Cả Shion-chan và Noel-chan đều đổ mồ hôi hột với cô nàng đang vui vẻ kia. Mọi chuyện
sau đó, chắc cũng không quá khó đoán.

Ngay lúc này, Akutan và Roboco-chan đang vừa bôi trát nhọ nồi lên mặt của hai người,
vừa cù léc họ, khiến cho đội thua vừa tức vừa cười.

``R-Rồi mọi người sẽ thấy! Hahaha! Lần sau tôi sẽ\dots\ haha, luyện thêm phép mạnh hơn!
Hahaha!'' Shion-chan vừa thề thốt, vừa lăn ra cười trước những bàn tay đang chuyển động
khắp người em ấy.

Noel-chan cũng không khá hơn, nhưng dù gì em ấy trông cũng vui hơn: ``Hahaha! Trận
vừa rồi chúng ta chơi vui thật đấy!''

Tôi nhìn căn bếp, thấy rằng mọi người đã ở công đoạn cuối chuẩn bị bữa tối. Vừa đủ để
tổ chức thêm một trận nữa.

Tôi để mặc bốn người họ chọc lét và trát nhọ lẫn nhau, gọi sáu người cuối cùng lên.

\ct

Trận cuối cùng là trận ba đấu ba, PekoRushiAki đấu với Ayame-Miko-Subaru. Cũng vì
đây là trận cuối cùng, tôi quyết định thay đổi luật chơi một chút: thay vì đạt năm điểm
trước là thắng, số điểm đã được nâng lên là bảy.

Không khí trở nên sôi động khi sáu người đứng đối diện nhau trên sân. Đây là trận đấu
cuối cùng, và dĩ nhiên, cả hai đội đều có sự kết hợp của những cá tính mạnh mẽ nhất.

Và đây có lẽ là trận Hanetsuki\dots\ kỳ quặc nhất mà tôi từng được thấy.

Cũng giống như ba trận trước, cả sáu cô gái đều ra oai và dọa dẫm bằng lời nói lẫn nhau.
Peko-chan thì ưỡn ngực ra, tỏ vẻ ngạo mạn: ``Đội này có Usada Pekora đây, chiến thắng
là chắc chắn rồi, peko!''

Rushia-chan trông khá căng thẳng nhưng cũng không kém cạnh gì, khi ánh mắt em ấy
lại rực lên sự quyết tâm: ``Hummu. Lần này Fandead của tôi sẽ giúp đội ta giành chiến
thắng! Lên nào, thuộc hạ của ta!''

Aki-chan thì chỉ cười nhẹ, một tay cầm chiếc vợt gỗ, tay kia vuốt tóc bay phấp phới trong
gió đông: ``Hãy chuẩn bị cho điều bất ngờ nhé\~{}''

Đội kia cũng đang hăng máu không kém. Ayame-chan nhếch môi cười đầy tự tin, tay cầm
chiếc vợt, nhưng đằng sau lưng vẫn đang giắt sẵn hai thanh katana: ``Ta đây sẽ chặn mọi
đòn đánh, yo!''

Miko-chan đứng cạnh họ, nhắm mắt tỏ vẻ tự tin, toát ra sự nghiêm nghị thường thấy:
``Cầu nguyện chiến thắng cho đội ta nào, ne!''

Subaru-chan với phong thái huấn luyện viên thể thao, đấm tay vào lòng bàn tay, ánh mắt
kiên định: ``Đừng lơ là! Lên chiến thuật và tập trung thi đấu, mọi người!''

Tôi thổi còi khai cuộc, với Pekora-chan thực hiện cú phát cầu cho đối thủ. Ngay từ lượt
đầu tiên, em ấy đã tung một quả cầu được\dots\ bọc bởi những củ cà rốt: ``Peko! Vũ khí bí
mật của tôi đây, peko!''

Quả cầu dần rơi xuống. Miko-chi cố gắng đỡ nó, nhưng rồi\dots

``À-rế? Sao nó hhẹ hều vậy, ne?'' nàng Vu nữ tỏ vẻ ngạc nhiên khi quả cầu không thể bay
cao lên được. Miko-chan cố gắng đuổi theo quả cầu đó, nhưng rồi cứ mỗi lần em ấy cố
tâng lên cao, quả cầu chỉ dập dìu trên không rồi cắm đầu xuống. ``C-Cái gì vậy ne!?''

Thế rồi, việc gì đến cũng phải đến, Miko-chan vấp ngã khi cố gắng liên tục tâng quả cầu
lên. Tỷ số bây giờ là 1-0.

``Pekora-chan chơi bài hay đấy!'' Ayame-chan gật gù nói, nở một nụ cười gian, trong khi
đó Miko-chan thì bày tỏ bất bình: ``Hay ho gì chứ, ne!''

Lượt phát cầu tiếp theo, Subaru-chan là người phát cầu cho đội bạn. Quả cầu được tung
lên rất cao, vượt quá tầm với của cả Pekora-chan, Rushia-chan và Aki-chan.

``Subaru-chan chơi hăng quá! Cao thế này thì làm sao bọn chị đỡ được chứ!'' Aki-chan
thảng thốt, cố gắng chạy theo quả cầu.

``Để em lo vụ này! Các Fandead của ta, lên nào!'' Ru-chan triệu hồi một bộ xương từ dưới
mặt đất lên, tay của nó cầm chiếc vợt gỗ. Bộ xương đó đỡ quả cầu, phát mạnh nó ngược
lại về phía sân đối phương.

Đến lượt Vịt-chan chạy ra xa đỡ lấy quả cầu. ``Rushia-chan tưởng thế mà gớm phết đấy!''
Em ấy đánh trả lại về sân đối phương.

``Lên nào, thuộc hạ của ta!''

Một bộ xương khác được triệu hồi lên, đập trả về phía sân của Ayame-Miko-Subaru. Và
tất nhiên, nàng Vịt đỡ được quả cầu đó.

Hai bên cứ trao qua trao lại quả cầu, và mỗi lần như vậy, nàng Pháp sư lại triệu hồi
thêm một bộ xương nữa, trong khi nàng Tomboy cũng không kém cạnh gì mà chạy qua
chạy lại như một vận động viên điền kinh.

Cho đến khi Subaru-chan đập quả cầu ấy rất mạnh lần nữa, lúc này sân của đối phương
đã chi chít hàng đống bộ xương. ``Nhận lấy này, shuba!''

Quả cầu được phát lên lần nữa, nhưng với phía sân đầy những bộ xương, tất cả bọn chúng
đều đuổi theo quả cầu đó. Vì số lượng bộ xương mà Ru-chan triệu ra là quá lớn, tất cả
chúng đều chen lấn xô đẩy nhau, cố gắng tìm cách đập trúng quả cầu đó.

``N-Này! Fandead, các ngươi đang làm gì vậy!? Đừng có đổ xô tất cả lên thế chứ!'' chủ
nhân của những bộ xương thảng thốt kêu lên, nhưng dường như chẳng ai trong bọn chúng
chịu nghe.

Hệ quả là, tất cả chúng đều xô lẫn nhau, ngã nhào ra đất, và không một ai trong só những
bộ xương đó đập trúng quả cầu cả, biếu không một điểm cho đội bạn.

Ru-chan thở dài, búng tay một cái, toàn bộ đám thuộc hạ của em ấy cũng lủi về xuống đất
trong sự hân hoan của đội đối phương.

``Đúng là lắm thầy nhiều ma mà, ne\~{}'' Miko-chan giở giọng châm chọc. Rushia-chan thốt
lên: ``T-Thì sao chứ!?Đừng có cười nhạo như thế!''

Tỷ số được cân bằng về 1-1. Tôi thổi còi, ra hiệu cho đội của Pekora-Rushia-Aki phát
cầu lên.

Lượt tiếp theo, Aki-chan là người đập cầu cho đối phương. Tôi nhìn thấy mắt của nàng
Oni sáng lên, có vẻ như em ấy lại có bài vở gì rồi.

Ayame-chan đạp trả lại quả cầu đầy mạnh mẽ, khiến nó bay xoáy như được trúng một
nhát kiếm: ``Thử đỡ xem sao, yo!''

Quả cầu quay với tốc độ cực nhanh và hiểm đến sắc bén, tới mức không cả Ru-chan và
Pekora-chan không kịp trở tay.

\begin{dialogue}
  \speak{Pekora} Quả cầu bay hiểm hóc quá, peko!
  \speak{Rushia} C-Cái gì thế này-nanodesu!?
\end{dialogue}

Tuy vậy, quả cầu mà nàng Dị giới thổi lại đi trúng vào vị trí của Subaru-chan. Nàng Vịt
không hề do dự gì mà tung cú đập mạnh với tốc độ như vận động viên điền kinh khiến
quả cầu rơi xuống sân đội Peko-chan, mang về một điểm trong đội.

Cứ như thế cả hai đội tiếp tục thi đấu với nhau với một bầu không khí căng thẳng. Người
triệu hồi, người chơi bài lừa, người đáp trả lại, khiến cho sức nóng của trận đấu càng tăng
thêm. Mỗi lần đội Thỏ - Pháp sư - Dị giới đạt được điểm, đội Vu nữ - Quỷ sừng - Vịt lại
có được điểm gỡ, biến trận đấu này thành một cuộc rượt đuổi tỷ số.

Trong một lượt phát cầu, Thỏ-chan tiếp tục dùng chiêu bài cà rốt của mình để tung hỏa
mù cho đối phương. Em ấy tung lên quả cầu bọc bên trong cà rốt, cùng với đó là những
củ cà rốt giả.

``Peko, peko\~{} Cà rốt thật ở đây này!''

Tiếp đó, Aki-chan sử dụng cặp bím tóc biết bay của mình để thổi tung chúng lên cao,
nhằm cho đối thủ rối bời vì không biết quả cầu thật sẽ bay theo hướng nào.

``Nhận lấy!''

May mắn thay, Ayame-chan sử dụng hai thanh katana của mình để chém những củ cà rốt
giả, rồi ngay tắp lự, lấy vợt gỗ đập quả cầu mạnh xuống đất, giành về một điểm trong đội.

Lượt tiếp theo, Rushia-chan tiếp tục triệu hồi Fandead, rồi kết hợp với nàng Dị giới để
thổi những quả cầu và những cục đá để tấn công đối phương.

``N-Này! Mấy người làm cái trò gì vậy ne!?'' Miko-chi vừa tránh những cục đá bay tới,
vừa xác định xem quả cầu thật sẽ rơi ở đâu.

Đúng lúc này, Subaru-chan mang dáng dấp của một huân luyện viên trưởng, liên tục hướng
dẫn hai người đồng đội phối hợp một cách chặt chẽ: ``Miko-chan, dùng vợt gỗ đứng chắn
trước! Ayame, xác định xem quả cầu nằm ở đâu, sau đó sử dụng đòn đánh mạnh vào!''

Ba người bọn họ phối hợp thế nào, đội bạn cũng không chịu thua, khi kết hợp đòn đánh
của mình để vừa tấn công, vừa đánh trả quả cầu xuống sân bên kia.

Tại một lượt phát cầu, Aki-chan đã tung một cú đánh cầu sử dụng một trong hai bím tóc
của mình, làm quả cầu bay vòng vèo, không thể xác định nổi đường đi của quả cầu.

``N-Ne!? Aki-chan chơi ăn gian!'' nàng Vu nữ chạy đi chạy lại, cố đón quả cầu đang rơi
xuống kia. Nhưng cuối cùng, quả cầu đã rơi xuống ngay chân của Miko-chi, giành luôn
một điểm từ họ.

Lượt cuối cùng, cũng là lượt quyết định cả trận đấu này, khi hai đội đã hòa nhau với tỷ số
6-6.

Miko-chi thực hiện cú phát cầu về phía Pekora-chan. Ngay lúc nàng Thỏ đã có thể đập
được trúng quả cầu kia\dots

``Tôi bắt được, tôi bắt được, tôi bắt được, peko!!''

\dots nhưng không ngờ quả cầu lại đột ngột thay đổi hướng đi khi chếch thẳng xuống dưới
như hàm logarit tự nhiên (\(y = \ln(x)\)) khi đi về giá trị không.

``N-Này! Hai người đang lừa tôi hả, peko!?'' Peko-chan tỏ vẻ bất bình. May cho em ấy,
Rushia-chan đã triệu hồi một Fandead của mình lên, đỡ được quả cầu trong gang tấc.

Quả cầu được trả về cho Subaru-chan. Em ấy bâng quơ đánh trúng quả cầu đó, khiến cho
nó đi theo quỹ đạo hypebol (\(y = ax^3\)), để rồi chạm xuống sân của PekoRushiAki, ấn định
chiến thắng sát nút 7-6 dành cho đội Miko-Ayame-Subaru.

Tiếng còi mãn cuộc vang lên từ tôi, và chưa cần A-san đưa chiếc nồi nhọ giờ đã vơi dần,
Ayame-chan đã hí hửng lấy luôn nó và trát lên mặt của cả ba người trong đội thua trận
kia.

Ru-chan rên rỉ trong khi bị Miko-chi trát mặt: ``Ể!? Tại sao chứ? Fandead của tôi không
thể giúp được sao?'' Vu nữ-chan đáp: ``Ai bảo em cứ lạm dụng nó làm chi, ne.''

Trong khi đó, Peko-chan thì bị cả Ayame-chan và Subaru-chan trát nhọ lên khắp mặt,
thậm chí còn ác hơn khi trát lên cả\dots\ đôi tai thỏ của em ấy. ``P-Peko!? Tai của tôi mà!
Đừng có trát lên nó chứ, peko!!''

Aki-chan thì nhẹ nhàng hơn, em ấy cười khì, chấp nhận thất bại với tỷ số sát nút: ``Cơ mà
chị thấy cũng vui mà! Trận đấu giao hữu này cũng đã kết thúc tốt đẹp rồi á\~{}''

Nhưng rồi, biến cố hài hước nữa đã xảy ra: tất cả mười chín người họ đều lao đến và tổ
chức trát nhọ nồi lên lẫn nhau!

\begin{dialogue}
  \speak{Fubuki} Lên nào mọi người ơi!
  \speak{Ayame} Đừng có đùa với quỷ Oni ta đây!
  \speak{Shion} Có còn tôi tạo thêm nhọ nồi không đây?
  \speak{Sora} Cho chị chơi cùng với!
\end{dialogue}

Tất cả bốn trận đấu và hai trò chơi dân gian trước đó đều được tôi và A-san quay lại hết
làm tư liệu làm phim. Nhìn tất cả mọi người đùa vui lẫn nhau sau những trận đấu căng
thẳng, cả hai chúng tôi đều không giấu được những nụ cười.

``Kể ra cậu cũng có nhiều trò đề bày ra vào dịp Tết nhỉ,'' chị đồng nghiệp của tôi nói, thở
dài nhưng lại mỉm cười nhẹ.

``Chắc là vì mọi người ở \hll\ tới chơi nên em mới bày ra mấy trò như vậy đấy, chứ
bình thường là có khi em cũng tự giam mình trong phòng để nghịch điện thoại rồi ạ,'' tôi
mỉm cười đáp lại.

``Đúng thật nhỉ. Lâu rồi họ mới có thời gian thư giãn đến mức thế.''

``Em không nghĩ việc mọi người thi đấu với nhau máu lửa như vậy lại là cách họ giải trí
đâu ạ\dots'' tôi cười trừ.

\il

Tôi cũng muốn bày thêm nhiều trò chơi dân gian nữa, nhưng rồi tiếng bố tôi vang lên từ
trong nhà khiến mọi người phải dừng lại: ``Mấy đứa ơi, vào ăn tối thôi! Mọi người đang
chờ đấy!''

Tất cả chúng tôi đặt vợt xuống, hào hứng cho bữa tối của nhà. Tôi gọi vọng lại: ``Vâng,
đợi bọn con tí!''

Quay lại nhìn mười chín cô gái đang be bét nhọ nồi đen trên những gương mặt vui vẻ đó,
tôi chợt bật cười. ``Trước tiên là mấy đứa đi rửa mặt đã, chứ mấy đứa nhìn như ma ấy!''

Tất cả chúng tôi bật cười với lời bình đó. Aki-chan chỉnh lại tóc tai, mỉm cười dịu dàng:
``Chơi xong đúng là đói thật\dots\ Mong món ăn đặc biệt tối nay quá!''

Mio-chan phủi bụi trên quần áo, khẽ cười mỉm: ``Em nghĩ chắc cũng giống như chúng ta
đã ăn vào mùng Một thôi.''

Noel chêm thêm một câu nói đùa, tay xoa bụng tỏ vẻ đói: ``Kuon-senpai, đừng quên là
em ăn khỏe lắm đấy! Anh đã chuẩn bị tinh thần chưa?''

Tôi cười, đáp lại: ``Miễn là đừng ăn hết phần của mọi người là được.''

Subaru-chan dẫn chúng tôi tới nhà vệ sinh để rửa mặt đầy nhọ nồi, trước khi vào nhà dùng
bữa. Không khí ấm áp từ bếp và tiếng cười nói của đại gia đình nhà ngoại nhanh chóng
lan tỏa, làm dịu đi cả sự mệt mỏi sau những trò chơi náo nhiệt.

\ct

Bữa tối bắt đầu với không khí nhộn nhịp. Không chỉ gia đình tôi và các cô gái, chúng tôi
còn ăn chung với gia đình nhà ông Niki-san, với một vài người mà tôi biết mặt, khác hoàn
toàn với bữa trưa với những người xa lạ kia.

Bánh chưng vẫn nằm ngay ngắn trên bàn, nhưng sự xuất hiện của các món ăn khác như
cơm, giò trứng, canh khoai tây, và thịt bò xào cần tỏi khiến mọi người không khỏi hào
hứng. Tôi nhìn vào bàn ăn, khẽ thở phào: ``May quá, hôm nay không chỉ có bánh chưng.
Chắc chắn sẽ không ai bị ép ăn thứ này cả tuần đâu.''

Mio-chan ngồi đối diện tôi cười khúc khích, vừa gắp miếng thịt, vừa trêu: ``Vậy là anh
không có tình yêu đặc biệt nào với bánh chưng à, Kuon-senpai?''

Tôi nhăn mặt, đáp: ``Không phải thế, chỉ là\dots\ ngày hôm qua, rồi bữa trưa nay đã ăn bánh
chưng rồi. Anh bắt đầu thấy sợ món này rồi á\dots''

Nghe vậy, Fubu-chan tiếp tục giở giọng trêu tôi, nhìn tôi với ánh mắt tinh nghịch: ``Cẩn
thận kẻo bà ngoại anh nghe thấy lại bắt anh ăn thêm đấy\~{}''

Tôi cười vội: ``Em cứ nói quá lên\dots\ Cơ mà bà ngoại anh thỉnh thoảng cũng hay ép anh ăn
thật.''

Bà ngoại tôi nghe thấy thế liền cười nhẹ nhàng nhưng giọng điệu rõ ràng đầy ý nhị: ``Bánh
chưng là tinh hoa của ngày Tết, ai lại nói sợ chứ. Thế mà lúc nhỏ, Kuon lúc nào cũng
giành ăn với Ryuuhou, phải không hả anh Jitori?''

Ông ngoại tôi gật gù đồng tình, xổ thẳng vào mặt tôi một câu đầy hài hước: ``Tao nhớ
chính thằng Cò này còn muốn đòi ăn cả của con Cún, đến mức còn bị bố nó mắng cơ
mà!''

Em gái tôi nghe vậy, cũng nhanh chóng chen vào tiếp sức châm chọc tôi: ``Giờ thì anh sợ
rồi, nên bánh chưng của anh năm nay cứ để em xử lý cho nha\~{}''

Tôi nở một nụ cười khó xử, vì bị cả ông bà ngoại, lẫn em gái, lẫn cả các cô gái \hll\ đều đổ vào châm chọc. Tôi chỉ có thể gượng gạo nói: ``Ừ-Ừ. Mọi người cứ việc. Anh
không ăn bánh chưng nữa đâu.''

Trong khi cả mâm mà tôi ăn đang cười, Pekora-chan đang ngồi ăn ở mâm khác, cũng dỏng
tai nghe những câu nói đùa mà chúng tôi vừa pha ban nãy.

Em ấy cắn miếng giò trứng, vừa xen vào với giọng điệu đầy hứng thú: ``Nhưng bánh
chưng này có cà rốt không, peko? Nếu không có thì\dots\ cũng hơi thất vọng ra đấy, peko!''

Ayame-chan gõ nhẹ vào vai em ấy và bật cười nhắc nhở: ``Peko-chan nói gì thế? Đây là
Tết mà, có phải mùa cầ rốt nhà em đâu!''

Ở một mâm khác, Choco-sensei dường như không để tâm lắm với cuộc hội thoại bánh
chưng này. Em ấy duyên dáng gắp một miếng thịt bò, rồi mỉm cười với mẹ tôi: ``Mà, mấy
món này nhà cô làm ngon thật đấy. Cô chú làm thế nào mà thịt vừa mềm lại thơm thế ạ?''

Mẹ tôi mỉm cười đáp lại: ``Cảm ơn cháu vì lời khen, Choco-chan. Thịt thì phải chọn đúng
loại, khi ướp thì phải thật lâu, rồi khi xào phải thật to lửa. Xong rồi\dots''

Tôi nhìn thấy Aki-chan ngồi bên cạnh tôi, tay chép lia lịa những gì mà mẹ tôi nói. Tôi
biết em ấy với Sensei là những người thích nấu ăn, nên việc học hỏi từ những người có
kinh nghiệm cũng là điều dễ hiểu.

Subaru-chan vẫn giữ tinh thần của một vận động viên chuyên nghiệp, vừa ăn vừa hào
hứng: ``Quả thật những món này là nguồn năng lượng lý tưởng sau khi chơi mấy trò vận
động!''

Mel-chan nhấp bát nước canh khoai tây, thở dài một tiếng tỏ vẻ thỏa mãn: ``Và chúng
cũng lạ miệng nữa!''

Shion-chan chậm rãi thưởng thức món bánh chưng cùng với Ru-chan, rồi quay sang nhận
xét:

\begin{dialogue}
  \speak{Shion} Vị bánh chưng này cũng lạ hơn so với hôm qua.
  \speak{Rushia} Kuon-senpai, bánh chưng này là nhà ngoại anh làm hay là đi mua vậy?
\end{dialogue}

Bố tôi đang ngồi mâm với các bác, các ông thay tôi trả lời: ``Bánh chưng nhà ngoại đây,
tự tay bà Aiko và vợ chú gói đấy. Thế mới giữ được hương vị đặc trưng.''

Matsuri-chan mắt sáng lên: ``Thế ạ? Vậy thì lần tới, chúng cháu sẽ gói thêm nhiều bánh
chưng cho nhà cô chú nữa ạ!''

Sự nhiệt tình của em ấy khiến tất cả chúng tôi bật cười. Cả bữa cơm tiếp tục trong không
khí ấm cúng, với những câu chuyện vui vẻ và tiếng cười không ngớt, như một cái kết trọn
vẹn cho mùng Hai Tết Nguyên Đán.

\ct

Ăn xong, chúng tôi quây quần lại trò chuyện với nhau thêm một chút. Những tiếng cười,
những câu chuyện trong ngày cứ thế nối tiếp không ngừng, làm ấm cả căn phòng giữa
buổi tối mùng Hai ngay tại quê ngoại Neihei nhà tôi.

Chúng tôi ngồi bên ngoài hiên nhà, nhìn ngắm bầu trời sao và hồi tưởng lại những gì
chúng tôi đã làm trong cả ngày hôm nay.

``Hôm nay mấy đứa thấy thế nào? Có vẻ như ai cũng hết sức vui vẻ nhỉ?''

Fubuki-chan lúc nào cũng là người đầu tiên lên tiếng, em cười toe toét: ``Chắc chắn rồi,
đặc biệt là trò Hanetsuki! Nhưng mà này, Mio, tôi không ngờ bà lại chơi nghiêm túc đến
thế đâu!''

Mio-chan gật đầu, cười một cách hiền hòa: ``Thì chúng ta phải có chút cạnh tranh chứ!
Mà Okayu, Korone, hai đứa định bỏ qua vụ thua bọn chị một cách dễ dàng vậy à?''

Korone-chan gãi đầu ngượng ngùng: ``Đúng, đúng! Mio-chan mà chơi nghiêm túc thì
bọn em cũng chịu thôi. Nhưng mà vui lắm, đúng không Ogayu?''

Okayu-chan nhún vai, đáp lại bằng giọng điệu lười biếng thường thấy: ``Cũng không tệ\~{}
Nhưng lần tới tôi muốn ăn gì đó trước khi chơi\dots\ Tại tôi đói quá không tập trung được
á\~{}''

Haato-chan nhảy chồm tới, tỏ vẻ bất bình hơn hẳn: ``Bọn em cũng chơi nghiêm túc đấy,
vậy mà vẫn thua! Choco-sensei, tại sao cô cứ phải dùng ngực để đỡ cầu chứ!?''

Choco-san cười nhẹ nhàng, nhưng khôn giấu được vẻ khiêu gợi: ``Chẳng phải Kuon-sama
nói rằng có thể dùng bất cứ chiêu gì để làm khó đối phương sao? Cô có gì thì dùng nấy
thôi chứ. Haato-sama cũng chẳng phải suýt ngã khi chế độ `Haaachama' của em được
kích hoạt đó sao?''

Nàng Song tính nghe em ấy phản biện lại, cũng chỉ ngập ngừng đáp: ``Q-Quả đúng là
vậy\dots\ Nhưng mà\dots''

Trong khi đó, Matsuri-chan thì lên tiếng khen về chiêu bài ``Người dơi'' của mình: ``Mà,
Mel-chan, bà trèo ngược trên cây trông như con dơi thật đó!''

``Ừm\dots\ Tại tôi thấy cách đó dễ nhìn cầu hơn mà,'' nàng Dơi đáp lại.

``Mỗi tội lại không mặc quần thôi,'' Haachama lên tiếng, làm cho Dơi-chan đỏ mặt: ``Đã
bảo là đừng nhắc tới chuyện đó nữa đi!''

Em ấy vội quay sang Matsuri-chan, người vẫn đang cười một cách khoái chí: ``M-Mà này,
Matsuri-chan, sức cổ động của bà cũng kinh ngạc lắm đấy!''

Ru-chan bất ngờ chen ngang mọi người, phồng má lên, cố tỏ vẻ mặt nghiêm túc đến mức
hơi buồn cười: ``Này, nhưng mấy người không thấy màn triệu hồi Fandead của em là ngầu
nhất à? Sao không ai khen em gì cả!''

Pekora-chan vừa nhai miếng bánh chưng cuối cùng, vừa cười ha hả: ``Ngầu thì đúng là
ngầu thật đấy, nhưng mà cái cách bọn Fandead của bà đuổi cả cầu lẫn tôi, rồi còn húc vào
nhau tranh cầu nữa thì cầng không ổn đâu, peko.''

Shion-chan cũng chêm vào: ``Đúng á. Em nên về nhà luyện thêm cách điều khiển bọn
chúng đi.'' Lập tức, Ru-chan đỏ mặt, cãi lại: ``Phép thuật của chị cũng có giúp đội chị
thắng được đâu!''

``Nhưng ít nhất thi chị vẫn hơn em ở khoản đối đáp đó!'' Phù thủy đáp lại, tỏ vẻ tự tin,
``Em thì làm sao có thể đọ được với chị!''

Tới lúc này, Sora-chan lên tiếng: ``Nhắc mới nhớ, một vài người trong chúng ta đã giúp
Kuon-senpai tránh được những những câu hỏi nhạy cảm nhỉ?''

Ayame-chan tỏ vẻ phấn khích, như thể tự hào với những gì em ấy đã làm: ``Về khoản đó
thì em, Shion và Subaru cùng Mio-chan đã xử gọn vụ này rồi, nhỉ?''

Sói đen-chan nghe vậy, gật gù: ``Nhưng mà phải thừa nhận là có vài người họ hàng của
anh ấy đúng là bất lịch sự thật.''

Subaru-chan tỏ vẻ thông cảm: ``Đúng á! Thậm chí còn có người xóc xỉa anh ấy nữa cơ
mà! Họ đang ghen tuông với anh ấy hay sao ấy!''

Miko-chan nhăn mặt, như thể đang tỏ vẻ khinh miệt với một số người mà em ấy đã gặp:
``Một số người thậm chí còn chúi mũi vào chuyện riêng của chúng mình nữa, ne.''

Roboco-chan thì tiến tới, nhìn tôi với ánh mắt thương cảm: ``Sao anh không thử bật chế
độ tự động trả lời như em được nhỉ, Kuon-senpai?''

Tôi cười xòa: ``Anh có phải là rô-bốt như em đâu. Với cả miệng anh cũng không hề dẻo
đâu.''

Chúng tôi tiếp tục ngồi kể cho nhau những gì mà chúng tôi đã làm trong cả ngày hôm nay.
Sora-chan thì chơi với rất nhiều em bé, đến mức còn bị chúng tôi gọi trêu là Em bé. Tất
nhiên là Sora-chan không thích điều đó.

Noel-chan thì hứng thú với các món ăn truyền thống và phi truyền thống trong hai ngày
Tết vừa rồi, và sẵn sàng ăn chúng cả tháng mà không sợ chán.

Ayame-chan thì cho chúng tôi xem lại những thước phim mà em ấy đã quay riêng để lưu
giữ kỷ niệm. Có những đoạn khi mà chúng tôi xem lại đều bật cười vì thú vị, nhưng có
những đoạn mà khi xem xong cũng đủ chúng tôi ớn người vì những hành động ngớ ngẩn
mà chúng tôi làm.

Chúng tôi cứ tiếp tục như thế, cho đến khi chúng tôi chuẩn bị ra về. Lúc đó đã là chín giờ
tối.

Ông ngoại và bố gọi chúng tôi. Trên tay họ là những món quà mà các bác, ông bà đã tự
tay làm cho chúng tôi.

``Giờ thì đến lúc mang quà quê về rồi. Không thể để khách tay không ra về được, đúng
chứ?'' ông ngoại tôi lên tiếng.
ất công chuẩn bị cho chúng ta rồi. Các cháu cứ
nhận đi nhé, cho mọi người ở đây vui,'' bố tôi tiếp lời.

Tôi gật đầu, rồi cùng mọi người đứng dậy, chuẩn bị đồ đạc ra về, đồng thời bốc dỡ đồ đạc
ra xe.

Quà quê gồm lại có bánh chưng, giò trứng, gạo, và tỷ thứ nữa. Nhìn đống quà chất cao
như núi, ai nấy đều không khỏi trầm trồ.

``Nhiều thế này thì làm sao mang hết về được nhỉ?'' Mio-chan nhìn núi hàng đó, tỏ vẻ băn
khoăn đến độ đổ mồ hôi hột.

May mắn thay, Noel-chan - người luôn tự hào với cơ bắp của mình - đã xắn tay áo, sẵn
sàng phục vụ: ``Không sao! Cứ để đó cho em, mọi người không cần lo!''

Roboco-chan nhấn nút trên cánh tay mình để kích hoạt chế độ tăng cường lực, rồi đi tới đống đồ đó, nhấc lên một núi đồ nữa một cách dễ dàng: ``Chị cũng muốn giúp nữa! Cánh
tay máu này có thể mang nhiều hơn bất cứ ai ở đây mà\~{}''

Nhìn hai người khỏe mạnh nhất bốc vác phần lớn hàng đi, chúng tôi cũng cảm thấy an
tâm phần nào.

Nhưng giờ, chúng tôi chỉ còn lại đúng một núi hàng nữa cần phải bốc vác. Núi hàng này
chỉ cao bằng một phần ba so với núi hàng cao nhất mà Hiệp sĩ-chan đang bê, nhưng cũng
đủ chúng tôi cảm thấy ngại.

Lúc này, mọi ánh mắt nhanh chóng đổ dồn về phía Akutan. Dường như em ấy nhận ra
ngay chúng tôi đang định làm gì, liền giơ tay lắc đầu nguầy nguậy: ``Đ-Đừng mà! Đừng
nhìn em kiểu đó!''

Fubu-chan cười gian xảo, tiến đến đặt tay lên vai Aqua-chan: ``Akutan à\~{}, chị để ý là em
đã ăn nhiều nhất trong bữa tối này đấy, đúng hơm? Phải dùng sức để đền đáp thôi\~{}''

Shion-chan cười lớn, vỗ tay cổ vuc: ``Đúng á! Đến Noel-chan cũng không ăn nhiều đến
thế đâu. Bà múp míp lắm rồi á!''

Hầu gái-chan nghe vậy, đỏ đến chín mặt. Em ấy vùng vẫy hai tay, thốt lên: ``T-Tôi không
có! Tôi không có béo đâu!''

``Vậy thì Akutan sẽ không có vấn đề gì khi xách chúng nhỉ, yo?'' tới lượt Ayame-chan giở
giọng châm chọc kiêm thách thức.

``Tất nhiên là được chứ! Xem tôi đây!'' Akutan tiến tới xách núi đồ còn lại trong sự bực
tức.

Nhưng ngay khi em nâng nó lên, cả thân người của em liền trụng xuống vì sức nặng của
núi hàng đó.

``Cái này\dots\ nặng quá!! Kuon-senpai! Cứu em!''

Không hề bỏ lỡ cơ hội, Peko-chan ở ngay cạnh đó, không những không giúp mà còn nhại
giọng của tiền bối: ``Aqua-senpai à, hầu gái gì mà yếu xìu thế, peko! Nhanh tay lên nào,
peko!''

Đúng lúc đó, một túi bánh chưng bị tuột khỏi tay Aqua-chab, rơi xuống đất với tiếng
*bịch!* thật lớn. Em hoảng hốt cúi xuống nhặt, nhưng tay vẫn ôm đầy túi khác, làm động
tác của em ấy trông không khác gì đang chơi trò xiếc.

\begin{dialogue}
  \speak{Aqua} Thôi chết! Bánh chưng rơi mất rồi- Á! Không, không! Túi trứng!
  \speak{Subaru} Cố lên nào, Aqua-chan! Bà làm được mà!
  \speak{Rushia} Bọn em tin ở chị-nanodesu!
  \speak{Aqua} Đừng có đứng đó cổ vũ nữa! Mau giúp chị đi\~{}!
\end{dialogue}

Thậm chí Mel-chan nhìn cảnh tượng éo le của em ấy cũng không thể nhịn cười: ``Aqua-chan
dễ thương thật đó!'' Nhưng rồi, em ấy nhìn hậu bối đang chật vật nhặt đồ và xách đồ bằng sự thương cảm: ``Nhưng mà\dots\ mọi người ơi, có chắc là giao hết cho em ấy được
không?''

May mắn thay cho Akutan, Noel-chan quay lại với nụ cười rạng rỡ, cầm thêm vài túi lớn
nữa trong núi hàng đó: ``Không sao, Mel-senpai. Để em mang phụ thêm, Aqua-senpai chỉ
cần làm phần nhỏ thôi mà.''

Mặc do số hàng đã vơi đi trên tay của Hầu gái, em ấy vẫn không tài nào xách hết chúng
được. ``Ai đó làm ơn cứu em với\dots''

Hết cách, chúng tôi cùng phụ em ấy xách đồ, rồi xếp chúng lên xe, nhưng không khí vẫn
đầy ắp tiếng cười.

Ngay khi Aqua-chan sau khi xong việc, em ấy ngồi thụp xuống bên cạnh chiếc xe, thở
dài mệt mỏi: ``Đã bảo rồi, em có phải là người chuyên việc này đâu\dots\ Lần sau mọi người
đừng làm khó em nữa, được chứ?''

Tất cả chúng tôi cười với phản ứng dỗi hờn đầy đáng yêu của Akutan. Roboco-chan kiểm
tra nốt những túi hàng mà chúng tôi đã chất lên xe, liền mỉm cười an ủi: ``Làm tốt lắm,
Aqua-chan. Nhưng mà, lần tới thì phải nâng cấp sức mạnh lên nha!''

Nàng Hầu gái chỉ biết cúi đầu bất lực, giữa tiếng cười của mọi người.

\ct

Chúng tôi quay về thành phố trong đêm tối mùng Hai Tết, với những thước phim ghi lại
cái Tết quê ngoại của Hikawa Kuon tôi đây, cùng các thành viên \hll. Một bữa tiệc
hiếm hoi mà không hề có kịch bản rõ ràng, chỉ toàn là những tình huống ``hổ lốn'' nhưng
lại đậm chất Tết Nguyên Đán - một thứ mà tôi chưa hề được trải nghiệm bao giờ trong suốt hơn hai mươi năm nay.

Bắt đầu từ một ngày mệt mỏi khi phải về quê ngoại, cho đến những câu hỏi cực khó nghe
từ những người mà tôi gọi là ``họ hàng'' kia\dots

\dots rồi đên những bữa cơm chan chứa tình thân, đến những tiếng cười rộn ràng của mọi
người\dots

\dots và kết thúc bằng những trò chơi gian đầy sôi động nhưng đầy đối kháng kia\dots
tất cả đều đã được ghi lại và khắc sâu trong trí nhớ tôi như một cuốn phim quý giá.

Đúng thật\dots\ có lẽ tôi đã hết từ để diễn tả chuyến du ngoạn này rồi.

Nhưng nếu có một điều tôi chắc chắn, đó là những khoảnh khắc tưởng như đơn giản này,
những tiếng cười giòn tan và cả những lần ``đấu khẩu'' đầy hài hước, sẽ mãi là một phần
không thể thiếu trong hành trình của tất cả chúng tôi.

Chúng ta vẫn còn ngày mùng Ba phía trước nữa, nhưng tôi tự hỏi rằng\dots\ liệu bao nhiêu
cái Tết sau nữa sẽ có thể tuyệt vời như lần này?

Nhưng trước mắt, dường như gia đình tôi đã có một kế hoạch để tận hưởng mùng Ba ngày mai rồi\dots

\end{document}